
\documentclass[10pt]{article}

\usepackage[a4paper,margin=0.9in,headsep=14pt,footskip=18pt]{geometry}
\usepackage{iftex}
\ifPDFTeX
\usepackage[T1]{fontenc}
\usepackage[utf8]{inputenc}
\usepackage{lmodern}
\else
\usepackage{fontspec}
\setmainfont{Liberation Serif}
\setsansfont{Liberation Sans}
\setmonofont{DejaVu Sans Mono}[Scale=0.92]
\fi

\usepackage{amsmath,amssymb}
\usepackage{graphicx}
\usepackage{booktabs,longtable,array,tabularx,calc}
\usepackage{enumitem}
\usepackage{setspace}
\usepackage{titlesec}
\usepackage{needspace}
\usepackage{xcolor}
\usepackage{hyperref}
\usepackage{fvextra}
\usepackage{fancyhdr}
\usepackage[most]{tcolorbox}
\tcbuselibrary{skins,breakable,listings}
\usepackage{listings}
\usepackage{float}
\usepackage{microtype}
\IfFileExists{xurl.sty}{\usepackage{xurl}}{}
\urlstyle{same}
\DeclareUnicodeCharacter{00D7}{$\times$}
\DeclareUnicodeCharacter{03B1}{$\alpha$}
\DeclareUnicodeCharacter{03B3}{$\gamma$}
\DeclareUnicodeCharacter{03C3}{$\sigma$}
\DeclareUnicodeCharacter{221A}{$\sqrt{}$}

\hypersetup{hidelinks}
\renewcommand{\theHfigure}{\arabic{figure}.\arabic{page}}
\setlength{\parindent}{0pt}
\setlength{\parskip}{0.25em}
\setlength{\tabcolsep}{4.5pt}
\renewcommand{\arraystretch}{1.12}
\setlength{\LTpre}{0.35em}
\setlength{\LTpost}{0.55em}
\setstretch{0.99}
\raggedbottom
\sloppy
\emergencystretch=3em
\hfuzz=100pt
\vfuzz=20pt
\hbadness=10000
\vbadness=10000

\definecolor{SectionBlue}{HTML}{2F4A68}
\definecolor{SectionRule}{HTML}{BCC8D2}
\definecolor{MetaGray}{HTML}{8A8A8A}
\definecolor{BugBlue}{HTML}{184F7A}
\definecolor{BugBlueBorder}{HTML}{6A94B6}
\definecolor{BugOrange}{HTML}{C8832A}
\definecolor{BugOrangeBorder}{HTML}{D8A96C}
\definecolor{BrickRed}{HTML}{C0392B}
\definecolor{OliveGreen}{HTML}{27AE60}
\definecolor{CodeBg}{HTML}{FAFAFA}
\definecolor{CodeFrame}{HTML}{CFCFCF}
\definecolor{DiffAdd}{HTML}{118847}
\definecolor{DiffDel}{HTML}{C53A28}

\pagestyle{fancy}
\fancyhf{}
\fancyhead[L]{\itshape QANet Bug Analysis Report}
\fancyhead[R]{COMP4329 Deep Learning}
\fancyfoot[C]{Page \thepage}
\renewcommand{\headrulewidth}{0.6pt}
\renewcommand{\footrulewidth}{0.6pt}
\setlength{\headheight}{14pt}

\titleformat{\section}[block]
 {\normalfont\bfseries\color{SectionBlue}\fontsize{17}{21}\selectfont}
 {}{0pt}{}
\titleformat{\subsection}[block]
 {\normalfont\bfseries\color{SectionBlue}\fontsize{14}{18}\selectfont}
 {}{0pt}{}

\newcommand{\tightlist}{%
 \setlength{\itemsep}{0.2em}\setlength{\parskip}{0pt}
}
\setlist[itemize]{leftmargin=1.4em,itemsep=0.3em,topsep=0.25em}
\setlist[enumerate]{leftmargin=1.4em,itemsep=0.3em,topsep=0.25em}

\newcommand{\ReportTitle}[3]{%
 \vspace*{0.6em}
 \begin{center}
 {\bfseries\fontsize{23}{28}\selectfont #1\par}
 \vspace{0.22em}
 {\fontsize{13.5}{16.2}\selectfont #2\par}
 \vspace{0.35em}
 {\color{MetaGray}\fontsize{11}{14}\selectfont #3\par}
 \end{center}
 \vspace{0.9em}
}

\newcommand{\ReportSection}[1]{%
 \Needspace{24\baselineskip}%
 \section*{#1}%
 \vspace{-0.35em}%
 {\color{SectionRule}\rule{\linewidth}{0.6pt}\par}%
 \vspace{0.35em}%
}

\newcommand{\ReportSubsection}[1]{%
 \Needspace{24\baselineskip}%
 \subsection*{#1}%
 \vspace{-0.12em}%
}

\newcommand{\FieldLabel}[1]{%
 {\bfseries #1:}\par
}

\newcommand{\SafeIncludeGraphics}[2][]{%
 \IfFileExists{#2}{\includegraphics[#1]{#2}}{%
  \fbox{\parbox[c][0.22\textheight][c]{0.9\linewidth}{\centering\textit{Figure placeholder}\\[0.35em]\small\detokenize{#2}}}
 }%
}

\newtcolorbox{bugbox}[1]{%
 enhanced,
 breakable,
 pad at break*=2mm,
 colback=white,
 colframe=BugBlueBorder,
 boxrule=0.5pt,
 arc=2mm,
 left=3.2mm,
 right=3.2mm,
 top=3.1mm,
 bottom=2.8mm,
 before={\par\needspace{18\baselineskip}\vspace{0.55em}},
 after skip=0.7em,
 attach boxed title to top left={xshift=3.2mm,yshift*=-\tcboxedtitleheight/2},
 boxed title style={
 colback=BugBlue,
 colframe=BugBlue,
 arc=1mm,
 boxrule=0.4pt,
 left=2.0mm,right=2.0mm,top=0.9mm,bottom=0.9mm
 },
 fonttitle=\bfseries\small\color{white},
 title={\parbox[t]{\dimexpr\linewidth-16mm\relax}{#1}}
}

\newtcolorbox{notebugbox}[1]{%
 enhanced,
 breakable,
 pad at break*=2mm,
 colback=white,
 colframe=BugOrangeBorder,
 boxrule=0.5pt,
 arc=2mm,
 left=3.2mm,
 right=3.2mm,
 top=3.1mm,
 bottom=2.8mm,
 before={\par\needspace{24\baselineskip}\vspace{0.55em}},
 after skip=0.7em,
 attach boxed title to top left={xshift=3.2mm,yshift*=-\tcboxedtitleheight/2},
 boxed title style={
 colback=BugOrange,
 colframe=BugOrange,
 arc=1mm,
 boxrule=0.4pt,
 left=2.0mm,right=2.0mm,top=0.9mm,bottom=0.9mm
 },
 fonttitle=\bfseries\small\color{white},
 title={\parbox[t]{\dimexpr\linewidth-16mm\relax}{#1}}
}


\newtcolorbox{rednotebugbox}[1]{%
 enhanced,
 breakable,
 pad at break*=2mm,
 colback=white,
 colframe=BrickRed!45,
 boxrule=0.5pt,
 arc=2mm,
 left=3.2mm,
 right=3.2mm,
 top=3.1mm,
 bottom=2.8mm,
 before={\par\needspace{24\baselineskip}\vspace{0.55em}},
 after skip=0.7em,
 attach boxed title to top left={xshift=3.2mm,yshift*=-\tcboxedtitleheight/2},
 boxed title style={
 colback=BrickRed,
 colframe=BrickRed,
 arc=1mm,
 boxrule=0.4pt,
 left=2.0mm,right=2.0mm,top=0.9mm,bottom=0.9mm
 },
 fonttitle=\bfseries\small\color{white},
 title={\parbox[t]{\dimexpr\linewidth-16mm\relax}{#1}}
}

\lstdefinestyle{diffstyle}{
 basicstyle=\ttfamily\footnotesize,
 columns=fullflexible,
 keepspaces=true,
 showstringspaces=false,
 breaklines=true,
 breakatwhitespace=false,
 tabsize=2,
 xleftmargin=0.4em,
 frame=leftline,
 framerule=1pt,
 rulecolor=\color{CodeFrame},
 framesep=0.6em,
 backgroundcolor=\color{CodeBg},
 aboveskip=0.15em,
 belowskip=0.15em,
 moredelim=[il][\color{DiffDel}]{-},
 moredelim=[il][\color{DiffAdd}]{+},
 morecomment=[l][\color{MetaGray}]{\#}
}

\begin{document}
\ReportTitle{QANet Bug Analysis Report}{Comprehensive debugging analysis of the repaired QANet implementation}{Structured form version with full code diffs}

\ReportSection{1 Introduction and Overview}

This report documents the repair of a QANet implementation and the empirical evidence that the repaired system now behaves like a trainable extractive question-answering model rather than a collection of disconnected code paths. The analysis is organized around two intertwined goals. First, the end-to-end pipeline must become executable and empirically meaningful: input processing, forward propagation, loss computation, optimisation, validation, and checkpoint reload must all operate coherently. Second, the learning mechanisms themselves must match their intended theoretical roles, so that later observations about optimisation, regularisation, convolution, attention, and decoding can be attributed to the model rather than to implementation artifacts.

The later mechanism study is motivated by three hypotheses derived from the repaired architecture. They focus on how information is processed inside the model rather than on superficial hyperparameter comparison: why architecturally identical encoder convolutions (Conv\_0 and Conv\_1) diverge by an order of magnitude in causal importance ($\Delta$F1 ratio $\approx 30\times$) despite sharing the same architecture, whether the directional asymmetry between C2Q and Q2C within CQ Attention persists and intensifies across the architecture change from BiDAF to QANet, and whether the Pointer layer's asymmetric wiring over repeated model-encoder passes reflects genuine functional specialisation or an untested inherited design. The debugging analysis below therefore serves two purposes at once: it explains what had to be repaired, and it establishes why later experiments can be interpreted as evidence about QANet itself.

\ReportSubsection{1.1 Debugging scope}

\begin{itemize}
\item
 \textbf{Structural pipeline errors}: defects that break the execution contract of the training or evaluation pipeline, including tensor-layout mismatches, invalid loss/backpropagation interfaces, incorrect runtime control flow, and checkpoint inconsistencies.
\item
 \textbf{Mechanism-level implementation flaws}: defects where code may run but the mechanism is mathematically wrong, theoretically inconsistent, or architecturally misaligned with the intended QANet design.
\end{itemize}

\ReportSubsection{1.2 Bug Distribution by Category}

\small
\setlength{\LTleft}{0pt}
\setlength{\LTright}{0pt}
\begin{longtable}{
  >{\raggedright\arraybackslash}p{0.10\linewidth}
  >{\raggedright\arraybackslash}p{0.28\linewidth}
  >{\centering\arraybackslash}p{0.06\linewidth}
  >{\raggedright\arraybackslash}p{0.48\linewidth}}
\toprule
\textbf{Section} & \textbf{Category} & \textbf{Count} & \textbf{Bug IDs} \\
\midrule
\endfirsthead
\toprule
\textbf{Section} & \textbf{Category} & \textbf{Count} & \textbf{Bug IDs} \\
\midrule
\endhead
\bottomrule
\endfoot
\textbf{\S 2} & Input Embedding \& Projection & 5 & BUG-021, BUG-012/013, BUG-014/015, BUG-064, BUG-N013 \\
\midrule
\textbf{\S 3} & Convolution Operators          & 3 & BUG-006, BUG-007/008, BUG-009/010/011 \\
              & Encoder Architecture           & 7 & BUG-016, BUG-017/018/019, BUG-N006, BUG-N008, BUG-N009, BUG-N010, BUG-N012 \\
              & Normalization, Activation \& Dropout & 6 & BUG-004, BUG-044, BUG-037, BUG-036, BUG-035, BUG-048 \\
              & Weight Initialization          & 4 & BUG-041, BUG-042, BUG-043, BUG-062 \\
\midrule
\textbf{\S 4} & Attention Mechanisms           & 5 & BUG-005, BUG-051, BUG-038/050, BUG-N003, BUG-049 \\
\midrule
\textbf{\S 5} & \multicolumn{3}{l}{\textit{Structural repairs only; documented inline in Section 5.}} \\
\midrule
\textbf{\S 6} & Output Layer \& Loss Function  & 4 & BUG-003, BUG-020, BUG-N007, BUG-N016 \\
\midrule
\textbf{\S 7} & Evaluation \& Inference        & 5 & BUG-001/002, BUG-N002, BUG-034, BUG-N011, BUG-N018 \\
              & Optimizers                     & 7 & BUG-022, BUG-023, BUG-052, BUG-053, BUG-039/040, BUG-054, BUG-055 \\
              & Learning Rate Schedulers       & 5 & BUG-024, BUG-056, BUG-057, BUG-058, BUG-N001 \\
              & Training Pipeline              & 5 & BUG-025, BUG-026, BUG-059, BUG-N015, BUG-N017 \\
\end{longtable}
\normalsize

\ReportSection{2 Input Embedding: Bug Repairs}

\ReportSubsection{2.1 Input Embedding \& Projection}

{\itshape
Word embeddings, character embeddings, highway transformation, and immediate embedding-side projection.
\par}

\begin{bugbox}{BUG-021 \checkmark{} Context word/char embedding lookup swap}
{\color{MetaGray}\small Stage: 1 \;\textbar\; Category: Runtime Failure / Forward Propagation \;\textbar\; Source: Models/qanet.py L70\par}
\vspace{0.22em}
\FieldLabel{Nature}
Incorrect structural pipeline implementation in input embedding routing.
\par
\vspace{0.35em}
\FieldLabel{Symptom}
Context word IDs are sent through the character-embedding table, while context character IDs are sent through the word-embedding table. This can raise invalid-index errors and, when indices remain numerically admissible, it still corrupts the context representation before encoding.
\par
\vspace{0.35em}
\FieldLabel{Root Cause}
Word IDs must index the word-embedding matrix and character IDs must index the character-embedding matrix. Reversing that routing violates the lookup contract: one failure mode is immediate index invalidity, while the subtler failure mode is a numerically valid but semantically incorrect context representation.
\par
\vspace{0.35em}
\FieldLabel{Corrective Action Taken}
Route Cwid through the word embedding table and Ccid through the character embedding table before entering the shared embedding module.
\par
\vspace{0.35em}
\begin{Verbatim}[commandchars=\\\{\},fontsize=\footnotesize,breaklines=true,breaksymbol={},frame=leftline,framerule=2pt,rulecolor=\color{gray!40},xleftmargin=6pt,formatcom=\singlespacing]
\textcolor{gray}{# Models/qanet.py L63 class QANet(nn.Module):}
qmask = (Qwid == 0)
\textcolor{BrickRed}{- Cw, Cc = self.char_emb(Cwid), self.word_emb(Ccid)}
\textcolor{OliveGreen}{+ Cw, Cc = self.word_emb(Cwid), self.char_emb(Ccid)}
Qw, Qc = self.word_emb(Qwid), self.char_emb(Qcid)
\end{Verbatim}
\vspace{0.22em}
\FieldLabel{Impact on Training}
\textbf{Pre-fix:} This defect fails in two different ways. If word indices exceed the character-vocabulary range, execution stops immediately. If the indices remain numerically admissible, the rest of the network can still run but it is now optimising on an input contract in which word-level and character-level semantics are exchanged, so later attention and pointer updates are internally consistent only with respect to the wrong lexical signal.\par\smallskip
\textbf{Post-fix:} The input boundary again preserves the intended word/character semantics. Subsequent encoder, attention, and output updates therefore act on the correct context signal, so improvements in loss and span metrics can be interpreted as genuine learning rather than compensation for a corrupted embedding route.
\par
\vspace{0.35em}
\end{bugbox}

\begin{bugbox}{BUG-012/013 \checkmark{} Highway transpose fixed}
{\color{MetaGray}\small Stage: 1 \;\textbar\; Category: Runtime Failure / Forward Propagation \;\textbar\; Source: Models/embedding.py L19\par}
\vspace{0.22em}
\FieldLabel{Nature}
Incorrect structural pipeline implementation in highway-network tensor layout.
\par
\vspace{0.35em}
\FieldLabel{Symptom}
The highway sublayer receives the batch axis as if it were the feature axis, which can produce linear-layer dimension failures or, in shape-compatible cases, semantically invalid cross-sample mixing.
\par
\vspace{0.35em}
\FieldLabel{Root Cause}
For tensors shaped {[}B, C, L{]}, the highway MLP should act on the feature axis after transposing to {[}B, L, C{]}. Using \texttt{transpose(0, 2)} moves the batch axis into the feature position, which violates the tensor-layout contract even when the resulting sizes happen to remain numerically admissible.
\par
\vspace{0.35em}
\FieldLabel{Corrective Action Taken}
Transpose dimensions 1 and 2 before the linear layers, then transpose back afterwards.
\par
\vspace{0.35em}
\begin{Verbatim}[commandchars=\\\{\},fontsize=\footnotesize,breaklines=true,breaksymbol={},frame=leftline,framerule=2pt,rulecolor=\color{gray!40},xleftmargin=6pt,formatcom=\singlespacing]
\textcolor{gray}{# Models/embedding.py L17 class Highway(nn.Module):}
def forward(self, x: torch.Tensor) -> torch.Tensor:
\textcolor{gray}{# x: [B, C, L] -> [B, L, C]}
\textcolor{BrickRed}{- x = x.transpose(0, 2)}
\textcolor{OliveGreen}{+ x = x.transpose(1, 2)}
for i in range(self.n):
gate = torch.sigmoid(self.gate[i](x))
\end{Verbatim}
\vspace{0.22em}
\FieldLabel{Impact on Training}
\textbf{Pre-fix:} If the axis mismatch is exposed, the highway layer raises a shape error before any valid update is taken. If it remains hidden by accidental size compatibility, the gating MLP treats the batch axis as part of the feature layout, so sample independence is broken and the embedding transform no longer represents within-example lexical composition. Even with all later modules correct, the model is then built on a contaminated representation basis.\par\smallskip
\textbf{Post-fix:} The highway gates now act on the channel features of each token within each example, which restores sample independence and the intended embedding transformation. Later optimisation and evaluation therefore reflect the behaviour of the model rather than an axis-handling artifact.
\par
\vspace{0.35em}
\end{bugbox}

\begin{bugbox}{BUG-014/015 \checkmark{} Character-tensor permutation before Conv2d fixed}
{\color{MetaGray}\small Stage: 1 \;\textbar\; Category: Runtime Failure / Forward Propagation \;\textbar\; Source: Models/embedding.py L44\par}
\vspace{0.22em}
\FieldLabel{Nature}
Incorrect structural pipeline implementation in character-convolution input layout.
\par
\vspace{0.35em}
\FieldLabel{Symptom}
Embedding.forward fails with a Conv2d channel mismatch because char\_limit appears in the channel position instead of d\_char.
\par
\vspace{0.35em}
\FieldLabel{Root Cause}
PyTorch Conv2d expects input in {[}N, C, H, W{]} order, where the second axis is the channel axis used by the convolution filters. The character tensor originally has shape {[}B, L, char\_len, dchar{]}. Permuting it as (0, 2, 1, 3) incorrectly puts char\_len into the channel slot.
\par
\vspace{0.35em}
\FieldLabel{Corrective Action Taken}
Change the permutation so the tensor becomes {[}B, dchar, L, char\_len{]} before Conv2d.
\par
\vspace{0.35em}
\begin{Verbatim}[commandchars=\\\{\},fontsize=\footnotesize,breaklines=true,breaksymbol={},frame=leftline,framerule=2pt,rulecolor=\color{gray!40},xleftmargin=6pt,formatcom=\singlespacing]
\textcolor{gray}{# Models/embedding.py L37 class Embedding(nn.Module):}
\textcolor{gray}{# ch_emb: [B, L, char_len, d_char]}
\textcolor{gray}{# wd_emb: [B, L, d_word]}
\textcolor{BrickRed}{- ch_emb = ch_emb.permute(0, 2, 1, 3)}
\textcolor{gray}{# [B, d_char, L, char_len]}
\textcolor{OliveGreen}{+ ch_emb = ch_emb.permute(0, 3, 1, 2)}
\textcolor{gray}{# [B, d_char, L, char_len]}
ch_emb = self.drop_char(ch_emb)

ch_emb = self.conv2d(ch_emb)
\end{Verbatim}
\vspace{0.22em}
\FieldLabel{Impact on Training}
\textbf{Pre-fix:} The explicit failure mode is immediate: the convolution expects channels in the second slot and instead receives the character-length axis, so the character path cannot execute. More importantly, this also clarifies the theoretical contract: even if one forced the shapes to agree elsewhere, the filters would still be attached to the wrong semantic axis, so the module would no longer be learning channel-wise character patterns within each token.\par\smallskip
\textbf{Post-fix:} The tensor now matches the \([N,C,H,W]\) contract expected by \texttt{Conv2d}, so the character encoder can learn over the intended feature axis. The downstream pipeline therefore receives a valid character signal rather than a broken or semantically misoriented tensor.
\par
\vspace{0.35em}
\end{bugbox}

\begin{bugbox}{BUG-064 \checkmark{} Pretrained word embeddings frozen}
{\color{MetaGray}\small Stage: 2 \;\textbar\; Category: Deep Learning Mechanism / Regularization Techniques \;\textbar\; Source: Models/qanet.py L44\par}
\vspace{0.22em}
\FieldLabel{Nature}
Inconsistent mechanism-level implementation flaw in pretrained embedding update policy.
\par
\vspace{0.35em}
\FieldLabel{Symptom}
GloVe word embeddings are fine-tuned during training instead of being frozen.
\par
\vspace{0.35em}
\FieldLabel{Root Cause}
The pretrained embedding constructor is called with freeze=False. In the intended QANet setup, pretrained lexical embeddings are held fixed so that the model learns higher-level composition on top of a stable lexical basis rather than drifting the base embedding space.
\par
\vspace{0.35em}
\FieldLabel{Corrective Action Taken}
Set freeze=True for the pretrained word embedding layer.
\par
\vspace{0.35em}
\begin{Verbatim}[commandchars=\\\{\},fontsize=\footnotesize,breaklines=true,breaksymbol={},frame=leftline,framerule=2pt,rulecolor=\color{gray!40},xleftmargin=6pt,formatcom=\singlespacing]
\textcolor{gray}{# Models/qanet.py L41 class QANet(nn.Module):}
self.word_emb = nn.Embedding.from_pretrained(
torch.tensor(word_mat, dtype=torch.float32),
\textcolor{BrickRed}{-}
freeze=False
\textcolor{OliveGreen}{+}
freeze=True
)
\end{Verbatim}
\vspace{0.22em}
\FieldLabel{Impact on Training}
\textbf{Pre-fix:} Pre-trained GloVe representations drift during training, increasing overfitting risk on the relatively small SQuAD dataset and weakening the intended lexical prior. Regularization is applied to the wrong quantity or with the wrong strength, so the model can overfit or suppress the wrong signal. \par\smallskip
\textbf{Post-fix:} Word embeddings remain frozen at their pre-trained values, preserving generalization and aligning the implementation with the QANet paper. 
\par
\vspace{0.35em}
\end{bugbox}

\begin{bugbox}{BUG-N013 \checkmark{} Character embedding cross-word leakage removed}
{\color{MetaGray}\small Stage: 2 \;\textbar\; Category: Deep Learning Mechanism / Embedding Module \;\textbar\; Source: Models/embedding.py L27-50\par}
\vspace{0.22em}
\FieldLabel{Nature}
Inconsistent mechanism-level implementation flaw in character-encoding locality.
\par
\vspace{0.35em}
\FieldLabel{Symptom}
The character embedding path applies a 2D depthwise-separable convolution over both token position and character position. As a result, character features from adjacent words may leak into each other.
\par
\vspace{0.35em}
\FieldLabel{Root Cause}
A per-word character encoder should convolve only across the character axis of one word at a time. A 2D convolution over token position and character position instead mixes lexical content across neighbouring words, which breaks the separation between word-internal composition and later contextual modelling.
\par
\vspace{0.35em}
\FieldLabel{Corrective Action Taken}
Replace the 2D character convolution with a 1D per-word convolution: flatten the batch and token axes together, convolve only over character positions, then max-pool and reshape back. This supersedes BUG-014/015: the earlier permutation repair was a necessary runtime fix for the original 2D path, but the final per-word 1D formulation removes that path entirely.
\par
\vspace{0.35em}
\begin{Verbatim}[commandchars=\\\{\},fontsize=\footnotesize,breaklines=true,breaksymbol={},frame=leftline,framerule=2pt,rulecolor=\color{gray!40},xleftmargin=6pt,formatcom=\singlespacing]
\textcolor{gray}{# Models/embedding.py L30 class Embedding(nn.Module):}
self.drop = Dropout(dropout)
self.drop_char = Dropout(dropout_char)
\textcolor{BrickRed}{- self.conv2d = DepthwiseSeparableConv(d_char, d_char, 5, dim=2, init_name=init_name)}
\textcolor{OliveGreen}{+ # [OLD] 2D DSConv(k=5x5) over [B, d_char, L, char_len] -- cross-word leakage}
\textcolor{OliveGreen}{+ # [FIX] 1D conv(k=5) over char_len only, per-word}
\textcolor{OliveGreen}{+ self.char_conv = DepthwiseSeparableConv(d_char, d_char, 5, dim=1, init_name=init_name)}
\textcolor{gray}{# Models/embedding.py L37 class Embedding(nn.Module):}
\textcolor{BrickRed}{- ch_emb = ch_emb.permute(0, 3, 1, 2)}
\textcolor{OliveGreen}{+ B, L, char_len, d_char = ch_emb.shape}
\textcolor{OliveGreen}{+ ch_emb = ch_emb.view(B * L, char_len, d_char).transpose(1, 2)}
ch_emb = self.drop_char(ch_emb)
\textcolor{BrickRed}{- ch_emb = self.conv2d(ch_emb)}
\textcolor{OliveGreen}{+ ch_emb = self.char_conv(ch_emb)}
ch_emb = self.act(ch_emb)
\textcolor{BrickRed}{- ch_emb, _ = torch.max(ch_emb, dim=3)}
\textcolor{OliveGreen}{+ ch_emb, _ = torch.max(ch_emb, dim=2)}
\textcolor{OliveGreen}{+ ch_emb = ch_emb.view(B, L, d_char).transpose(1, 2)}
\end{Verbatim}
\vspace{0.22em}
\FieldLabel{Impact on Training}
\textbf{Pre-fix:} Even if later contextual layers are correct, the character pathway has already leaked neighbouring-word information into what should be a within-word lexical representation. That gives the model an unintended shortcut: contextual evidence can enter before the designated fusion and encoder stages, which makes later component roles harder to interpret and can distort how lexical evidence is learned.\par\smallskip
\textbf{Post-fix:} Each word is now encoded from its own character sequence only. The character route therefore contributes morphological and subword information without prematurely performing contextual mixing.
\par
\vspace{0.35em}
\end{bugbox}

\ReportSection{3 Embedding Encoder: Bug Repairs}

\ReportSubsection{3.1 Convolution Operators}

{\itshape
Custom Conv1d, Conv2d padding, and depthwise-separable convolution behaviour within the encoder stack.
\par}

\begin{bugbox}{BUG-006 \checkmark{} Conv1d unfold dimension fixed}
{\color{MetaGray}\small Stage: 1 \;\textbar\; Category: Runtime Failure / Forward Propagation \;\textbar\; Source: Models/conv.py L55\par}
\vspace{0.22em}
\FieldLabel{Nature}
Incorrect structural pipeline implementation in Conv1d sliding-window axis selection.
\par
\vspace{0.35em}
\FieldLabel{Symptom}
Conv1d unfolds along the channel dimension rather than the sequence dimension, producing invalid intermediate layouts and causing incorrect convolution behaviour.
\par
\vspace{0.35em}
\FieldLabel{Root Cause}
For a channel-first 1D input {[}B, C, L{]}, convolution slides over the last dimension L, not over the channel dimension C. In PyTorch, unfold(dim,\ldots) extracts moving windows along the specified axis. Using dim=1 therefore extracts windows across channels rather than across positions.
\par
\vspace{0.35em}
\FieldLabel{Corrective Action Taken}
Change the unfold dimension from 1 to 2.
\par
\vspace{0.35em}
\begin{Verbatim}[commandchars=\\\{\},fontsize=\footnotesize,breaklines=true,breaksymbol={},frame=leftline,framerule=2pt,rulecolor=\color{gray!40},xleftmargin=6pt,formatcom=\singlespacing]
\textcolor{gray}{# Models/conv.py L53 class Conv1d(nn.Module):}
L_out = x.size(2) - self.kernel_size + 1
\textcolor{BrickRed}{- x_unf = x.unfold(1, self.kernel_size, 1)}
\textcolor{OliveGreen}{+ x_unf = x.unfold(2, self.kernel_size, 1)}
\end{Verbatim}
\vspace{0.22em}
\FieldLabel{Impact on Training}
\textbf{Pre-fix:} Sliding windows are extracted over the wrong axis, producing invalid tensor layouts for grouped convolution and causing shape/runtime failures. The affected path therefore cannot produce a valid loss, gradient, or validation signal.\par\smallskip
\textbf{Post-fix:} Unfold now operates over sequence length, so convolution windows and grouped reshapes match the intended Conv1d computation. Once the contract is restored, later loss and metric changes reflect model learning rather than pipeline repair.
\par
\vspace{0.35em}
\end{bugbox}

\begin{bugbox}{BUG-007/008 \checkmark{} Conv2d width-padding height mismatch fixed}
{\color{MetaGray}\small Stage: 1 \;\textbar\; Category: Runtime Failure / Forward Propagation \;\textbar\; Source: Models/conv.py L124\par}
\vspace{0.22em}
\FieldLabel{Nature}
Incorrect structural pipeline implementation in Conv2d padding construction.
\par
\vspace{0.35em}
\FieldLabel{Symptom}
2D convolutions with padding fail during width padding because the width-pad tensor height does not match the already height-padded input.
\par
\vspace{0.35em}
\FieldLabel{Root Cause}
After height padding, the tensor height becomes H + 2p. if the width-padding tensor is still allocated with the original H, concatenation fails because the non-concatenated dimensions must match exactly. The implementation used stale height metadata instead of the updated x.size(2).
\par
\vspace{0.35em}
\FieldLabel{Corrective Action Taken}
Allocate the width-padding tensor using the post-height-padding size rather than the original height.
\par
\vspace{0.35em}
\begin{Verbatim}[commandchars=\\\{\},fontsize=\footnotesize,breaklines=true,breaksymbol={},frame=leftline,framerule=2pt,rulecolor=\color{gray!40},xleftmargin=6pt,formatcom=\singlespacing]
\textcolor{gray}{# Models/conv.py L122 class Conv2d(nn.Module):}
pad_h = x.new_zeros(B, C_in, p, W)
x = torch.cat([pad_h, x, pad_h], dim=2)
\textcolor{BrickRed}{- pad_w = x.new_zeros(B, C_in, H, p)}

\textcolor{OliveGreen}{+ pad_w = x.new_zeros(B, C_in, x.size(2), p)}
x = torch.cat([pad_w, x, pad_w], dim=3)
\end{Verbatim}
\vspace{0.22em}
\FieldLabel{Impact on Training}
\textbf{Pre-fix:} Conv2d width-padding concatenation fails because the tensor heights differ, blocking execution of the convolution stack. The affected path therefore cannot produce a valid loss, gradient, or validation signal.\par\smallskip
\textbf{Post-fix:} Width padding now uses the correct post-height-padding shape, so concatenation and Conv2d forward execution proceed normally. Once the contract is restored, later loss and metric changes reflect model learning rather than pipeline repair.
\par
\vspace{0.35em}
\end{bugbox}

\begin{bugbox}{BUG-009/010/011 \checkmark{} Depthwise-separable convolution order restored}
{\color{MetaGray}\small Stage: 1 \;\textbar\; Category: Runtime Failure / Forward Propagation \;\textbar\; Source: Models/conv.py L175\par}
\vspace{0.22em}
\FieldLabel{Nature}
Incorrect structural pipeline implementation in depthwise-separable convolution ordering.
\par
\vspace{0.35em}
\FieldLabel{Symptom}
Layers where in\_ch != out\_ch fail with channel/group mismatches, and even when shapes happen to match the composed operation is still not the intended separable convolution.
\par
\vspace{0.35em}
\FieldLabel{Root Cause}
A depthwise-separable convolution is, by definition, depthwise first and pointwise second. The depthwise stage preserves the channel count and applies grouped spatial filtering independently per channel, while the pointwise stage mixes channels afterward. Reversing the order changes the operator entirely and can invalidate the grouped-convolution assumptions.
\par
\vspace{0.35em}
\FieldLabel{Corrective Action Taken}
Apply the depthwise convolution first and the pointwise convolution second.
\par
\vspace{0.35em}
\begin{Verbatim}[commandchars=\\\{\},fontsize=\footnotesize,breaklines=true,breaksymbol={},frame=leftline,framerule=2pt,rulecolor=\color{gray!40},xleftmargin=6pt,formatcom=\singlespacing]
\textcolor{gray}{# Models/conv.py L173 class DepthwiseSeparableConv(nn.Module):}
def forward(self, x):
\textcolor{BrickRed}{- return self.depthwise_conv(self.pointwise_conv(x))}
\textcolor{OliveGreen}{+ return self.pointwise_conv(self.depthwise_conv(x))}
\end{Verbatim}
\vspace{0.22em}
\FieldLabel{Impact on Training}
\textbf{Pre-fix:} Pointwise-first changes the channel count before grouped depthwise convolution, causing crashes when channel assumptions are violated and producing the wrong operator even when execution survives. If the mismatch becomes explicit, execution stops before a valid loss or metric can be produced. If it remains latent because dimensions stay compatible, downstream attention, optimization, and evaluation are driven by the incorrect semantic signal.\par\smallskip
\textbf{Post-fix:} Depthwise is now applied before pointwise as intended, restoring correct separable-convolution semantics and preserving grouped-convolution channel assumptions. Later layers now receive the intended signal, so downstream loss and span metrics become meaningful evidence of learning.
\par
\vspace{0.35em}
\end{bugbox}

\ReportSubsection{3.2 Encoder Architecture}

{\itshape
Positional encoding, normalization indexing, feed-forward structure, shared encoder design, projection layers, and stochastic depth.
\par}

\begin{bugbox}{BUG-016 \checkmark{} Positional-encoding frequency broadcasting fixed}
{\color{MetaGray}\small Stage: 1 \;\textbar\; Category: Runtime Failure / Forward Propagation \;\textbar\; Source: Models/encoder.py L29-32\par}
\vspace{0.22em}
\FieldLabel{Nature}
Incorrect structural pipeline implementation in positional-encoding broadcast construction.
\par
\vspace{0.35em}
\FieldLabel{Symptom}
EncoderBlock construction usually fails while building positional encodings because the frequency tensor does not broadcast against the {[}C, L{]} position grid.
\par
\vspace{0.35em}
\FieldLabel{Root Cause}
To broadcast over a {[}C, L{]} grid, the frequency vector requires shape {[}C, 1{]}, not {[}1, C{]}. Under PyTorch broadcasting semantics, trailing dimensions align from the right, so unsqueezing on dimension 0 produces the incorrect orientation.
\par
\vspace{0.35em}
\FieldLabel{Corrective Action Taken}
Unsqueeze the frequency tensor on dimension 1 instead of dimension 0.
\par
\vspace{0.35em}
\begin{Verbatim}[commandchars=\\\{\},fontsize=\footnotesize,breaklines=true,breaksymbol={},frame=leftline,framerule=2pt,rulecolor=\color{gray!40},xleftmargin=6pt,formatcom=\singlespacing]
\textcolor{gray}{# Models/encoder.py L30 class PosEncoder(nn.Module):}
\textcolor{BrickRed}{- ).unsqueeze(0)}
\textcolor{OliveGreen}{+ ).unsqueeze(1)}
\end{Verbatim}
\vspace{0.22em}
\FieldLabel{Impact on Training}
\textbf{Pre-fix:} Positional encoding construction fails due to broadcast mismatch, so the model cannot initialize and training or evaluation cannot even begin. The affected path therefore cannot produce a valid loss, gradient, or validation signal.\par\smallskip
\textbf{Post-fix:} The frequency tensor now has the correct {[}C, 1{]} shape, making broadcasting valid and allowing initialization to proceed. Once the contract is restored, later loss and metric changes reflect model learning rather than pipeline repair.
\par
\vspace{0.35em}
\end{bugbox}

\begin{bugbox}{BUG-017/018/019 \checkmark{} Encoder normalization indexing fixed}
{\color{MetaGray}\small Stage: 1 \;\textbar\; Category: Runtime Failure / Forward Propagation \;\textbar\; Source: Models/encoder.py L131\par}
\vspace{0.22em}
\FieldLabel{Nature}
Incorrect structural pipeline implementation in encoder normalization indexing.
\par
\vspace{0.35em}
\FieldLabel{Symptom}
self.norms{[}i+1{]} goes out of bounds on the final convolution iteration, causing an IndexError.
\par
\vspace{0.35em}
\FieldLabel{Root Cause}
Python lists are zero-indexed. If a list has conv\_num elements, the valid indices are 0,... , conv\_num -1. On the last loop iteration, i = conv\_num -1, so i + 1 = conv\_num is out of bounds.
\par
\vspace{0.35em}
\FieldLabel{Corrective Action Taken}
Use self.norms{[}i{]} rather than self.norms{[}i+1{]}.
\par
\vspace{0.35em}
\begin{Verbatim}[commandchars=\\\{\},fontsize=\footnotesize,breaklines=true,breaksymbol={},frame=leftline,framerule=2pt,rulecolor=\color{gray!40},xleftmargin=6pt,formatcom=\singlespacing]
\textcolor{gray}{# Models/encoder.py L119 class EncoderBlock(nn.Module):}
\textcolor{BrickRed}{- out = self.norms[i + 1](out)}
\textcolor{OliveGreen}{+ out = self.norms[i](out)}
\end{Verbatim}
\vspace{0.22em}
\FieldLabel{Impact on Training}
\textbf{Pre-fix:} The final convolution stage indexes one element past the normalization list boundary, so encoder blocks cannot complete forward execution. The affected path therefore cannot produce a valid loss, gradient, or validation signal.\par\smallskip
\textbf{Post-fix:} Each convolution stage now uses its matching normalization module, eliminating the out-of-range access and restoring stable encoder execution. Once the contract is restored, later loss and metric changes reflect model learning rather than pipeline repair.
\par
\vspace{0.35em}
\end{bugbox}

\begin{bugbox}{BUG-N006 \checkmark{} Two-layer FFN restored}
{\color{MetaGray}\small Stage: 2 \;\textbar\; Category: Deep Learning Mechanism / Encoder Architecture \;\textbar\; Source: Models/encoder.py L96-97, L144-148\par}
\vspace{0.22em}
\FieldLabel{Nature}
Inconsistent mechanism-level implementation flaw in feed-forward sublayer design.
\par
\vspace{0.35em}
\FieldLabel{Symptom}
Model encoder blocks have reduced expressive capacity, so the model may show decreasing loss while F1 improves only slowly.
\par
\vspace{0.35em}
\FieldLabel{Root Cause}
An FFN in the Transformer/QANet sense is not just ``an affine map plus an activation''; it is a two-layer hidden MLP of the form W2σ(W1x + b1) + b2. With only one linear layer, the block loses the intended hidden transformation stage and therefore cannot realize the same class of nonlinear mappings.
\par
\vspace{0.35em}
\FieldLabel{Corrective Action Taken}
Replace the single linear layer with a two-layer feed-forward block and place the activation between the two projections.
\par
\vspace{0.35em}
\begin{Verbatim}[commandchars=\\\{\},fontsize=\footnotesize,breaklines=true,breaksymbol={},frame=leftline,framerule=2pt,rulecolor=\color{gray!40},xleftmargin=6pt,formatcom=\singlespacing]
\textcolor{gray}{# Models/encoder.py L97 class EncoderBlock(nn.Module):}
\textcolor{BrickRed}{- self.fc = nn.Linear(d_model, d_model, bias=True)}
\textcolor{OliveGreen}{+ self.fc1 = nn.Linear(d_model, d_model, bias=True)}
\textcolor{OliveGreen}{+ self.fc2 = nn.Linear(d_model, d_model, bias=True)}
\textcolor{gray}{# forward}
\textcolor{BrickRed}{- out = self.fc(out.transpose(1, 2)).transpose(1, 2)}
\textcolor{OliveGreen}{+ out = out.transpose(1, 2)}
\textcolor{OliveGreen}{+ out = self.fc1(out)}
out = self.act(out)
\textcolor{OliveGreen}{+ out = self.fc2(out)}
\textcolor{OliveGreen}{+ out = out.transpose(1, 2)}
\end{Verbatim}
\vspace{0.22em}
\FieldLabel{Impact on Training}
\textbf{Pre-fix:} Even if attention and convolution are implemented correctly, each encoder block still lacks one of its three intended transformations. The block can only apply a single affine map after attention, so channel-wise reparameterisation is much weaker than intended and later behaviour reflects a reduced-capacity architecture rather than the repaired QANet design.\par\smallskip
\textbf{Post-fix:} The encoder block again contains a genuine two-layer position-wise network, so nonlinear channel mixing is restored alongside convolution and attention. This makes later performance and ablation results attributable to the intended encoder family.
\par
\vspace{0.35em}
\end{bugbox}

\begin{bugbox}{BUG-N008 \checkmark{} Shared context/question encoder structure restored}
{\color{MetaGray}\small Stage: 2 \;\textbar\; Category: Deep Learning Mechanism / Encoder Architecture \;\textbar\; Source: Models/Normalizations/normalization.py L8-20 + Models/qanet.py L50-53, 74-75, 79-80\par}
\vspace{0.22em}
\FieldLabel{Nature}
Inconsistent mechanism-level implementation flaw in shared encoder structure and channel-only normalization.
\par
\vspace{0.35em}
\FieldLabel{Symptom}
Context and question branches are encoded by independent projection and embedding-encoder modules, so the two sequences are not constrained to enter attention in a shared feature space.
\par
\vspace{0.35em}
\FieldLabel{Root Cause}
The QANet embedding encoder is shared so that context and question tokens are normalized and transformed by the same operator before attention. Breaking that sharing creates two separate representation geometries, while normalization that is not channel-consistent further weakens the comparability on which attention relies.
\par
\vspace{0.35em}
\FieldLabel{Corrective Action Taken}
Introduce channel-only LayerNorm for channel-first tensors, and merge separate context/question projection and embedding encoder modules into shared components.
\par
\vspace{0.35em}
\begin{Verbatim}[commandchars=\\\{\},fontsize=\footnotesize,breaklines=true,breaksymbol={},frame=leftline,framerule=2pt,rulecolor=\color{gray!40},xleftmargin=6pt,formatcom=\singlespacing]
\textcolor{gray}{# Models/Normalizations/normalization.py L8-20}
\textcolor{OliveGreen}{+ import torch}
 import torch.nn as nn
 from.layernorm import LayerNorm
 from.groupnorm import GroupNorm
\textcolor{OliveGreen}{+}
\textcolor{OliveGreen}{+ class _ChannelFirstLayerNorm(nn.Module):}
\textcolor{OliveGreen}{+  """LayerNorm wrapper for channel-first [B, C, L] tensors."""}
\textcolor{OliveGreen}{+  def __init__(self, d_model: int):}
\textcolor{OliveGreen}{+   super().__init__()}
\textcolor{OliveGreen}{+   self.norm = LayerNorm(d_model)}
\textcolor{OliveGreen}{+}
\textcolor{OliveGreen}{+  def forward(self, x: torch.Tensor) -> torch.Tensor:}
\textcolor{OliveGreen}{+   return self.norm(x.transpose(1, 2)).transpose(1, 2)}
 
\textcolor{BrickRed}{- def get_norm(name: str, d_model: int, length: int, num_groups: int = 8) -> nn.Module:}
\textcolor{OliveGreen}{+ def get_norm(name: str, d_model: int, length: int = 0, num_groups: int = 8) -> nn.Module:}
\textcolor{BrickRed}{-  if name == "layer_norm":}
\textcolor{BrickRed}{-   return LayerNorm([d_model, length])}
\textcolor{OliveGreen}{+  if name == "layer_norm":}
\textcolor{OliveGreen}{+   return _ChannelFirstLayerNorm(d_model)}
  return GroupNorm(num_groups, d_model)

\textcolor{gray}{# Models/qanet.py L50-53, 74-75, 79-80}
\textcolor{BrickRed}{- self.context_conv = DepthwiseSeparableConv(d_word + d_char, d_model, 5, init_name=init_name)}
\textcolor{BrickRed}{- self.question_conv = DepthwiseSeparableConv(d_word + d_char, d_model, 5, init_name=init_name)}
\textcolor{OliveGreen}{+ self.proj_conv = DepthwiseSeparableConv(d_word + d_char, d_model, 5, init_name=init_name)}

\textcolor{BrickRed}{- self.c_emb_enc = EncoderBlock(d_model, num_heads, dropout, conv_num=4, k=7, length=len_c, init_name=init_name, act_name=act_name, norm_name=norm_name, norm_groups=norm_groups)}
\textcolor{BrickRed}{- self.q_emb_enc = EncoderBlock(d_model, num_heads, dropout, conv_num=4, k=7, length=len_q, init_name=init_name, act_name=act_name, norm_name=norm_name, norm_groups=norm_groups)}
\textcolor{OliveGreen}{+ self.emb_enc = EncoderBlock(d_model, num_heads, dropout, conv_num=4, k=7, length=max(len_c, len_q), init_name=init_name, act_name=act_name, norm_name=norm_name, norm_groups=norm_groups)}

\textcolor{BrickRed}{- C = self.context_conv(C)}
\textcolor{BrickRed}{- Q = self.question_conv(Q)}
\textcolor{OliveGreen}{+ C = self.proj_conv(C)}
\textcolor{OliveGreen}{+ Q = self.proj_conv(Q)}

\textcolor{BrickRed}{- Ce = self.c_emb_enc(C, cmask)}
\textcolor{BrickRed}{- Qe = self.q_emb_enc(Q, qmask)}
\textcolor{OliveGreen}{+ Ce = self.emb_enc(C, cmask)}
\textcolor{OliveGreen}{+ Qe = self.emb_enc(Q, qmask)}
\end{Verbatim}
\vspace{0.22em}
\FieldLabel{Impact on Training}
\textbf{Pre-fix:} If the two streams are encoded by different operators, later context--question attention is asked to compare vectors that do not live in the same learned geometry. The model may still converge, but any similarity score is now burdened with compensating for representational mismatch, so the fusion layer mixes ``context vs. question alignment'' with ``branch-specific calibration error.''\par\smallskip
\textbf{Post-fix:} Context and question now pass through the same projection and embedding encoder, with consistent channel-wise normalization. Cross-attention can therefore compare like with like, which restores a coherent basis for later fusion and span prediction.
\par
\vspace{0.35em}
\end{bugbox}

\begin{bugbox}{BUG-N009 \checkmark{} Embedding projection changed to standard Conv1d(k=1)}
{\color{MetaGray}\small Stage: 2 \;\textbar\; Category: Deep Learning Mechanism / Convolutional Modules \;\textbar\; Source: Models/qanet.py L50\par}
\vspace{0.22em}
\FieldLabel{Nature}
Inconsistent mechanism-level implementation flaw in embedding-side projection operator design.
\par
\vspace{0.35em}
\FieldLabel{Symptom}
The embedding-side projection is implemented as a depthwise-separable convolution, so this stage performs local spatial mixing rather than a pure per-position channel projection.
\par
\vspace{0.35em}
\FieldLabel{Root Cause}
This stage only needs to map $(d_{word}+d_{char})$ channels into $d_{model}$ at each token position. A \texttt{Conv1d} with kernel size 1 performs exactly that cross-channel mixing, whereas a wider separable convolution introduces cross-position receptive-field expansion that belongs to later encoder layers, not to the embedding projection.
\par
\vspace{0.35em}
\FieldLabel{Corrective Action Taken}
Replace the depthwise-separable projection with a standard pointwise Conv1d of kernel size 1.
\par
\vspace{0.35em}
\begin{Verbatim}[commandchars=\\\{\},fontsize=\footnotesize,breaklines=true,breaksymbol={},frame=leftline,framerule=2pt,rulecolor=\color{gray!40},xleftmargin=6pt,formatcom=\singlespacing]
\textcolor{gray}{# Models/qanet.py L50}
\textcolor{BrickRed}{- self.proj_conv = DepthwiseSeparableConv(d_word + d_char, d_model, 5, init_name=init_name)}
\textcolor{OliveGreen}{+ self.proj_conv = Conv1d(d_word + d_char, d_model, kernel_size=1, bias=False)}
\textcolor{OliveGreen}{+ initializations[init_name](self.proj_conv.weight)}
\end{Verbatim}
\vspace{0.22em}
\FieldLabel{Impact on Training}
\textbf{Pre-fix:} Even when all later layers are correct, this projection stage already mixes neighbouring tokens. That destroys a clean separation between \emph{channel projection} and \emph{contextual composition}: later convolutions may appear useful or redundant partly because the projection has already performed some of their job.\par\smallskip
\textbf{Post-fix:} The projection is again token-wise and channel-wise only. Local contextual mixing is deferred to the encoder blocks, so later training behaviour and architectural analysis can be attributed to the intended stage ordering.
\par
\vspace{0.35em}
\end{bugbox}

\begin{bugbox}{BUG-N010 \checkmark{} Model-encoder input projection changed to standard Conv1d(k=1)}
{\color{MetaGray}\small Stage: 2 \;\textbar\; Category: Deep Learning Mechanism / Convolutional Modules \;\textbar\; Source: Models/qanet.py L58\par}
\vspace{0.22em}
\FieldLabel{Nature}
Inconsistent mechanism-level implementation flaw in model-encoder input projection design.
\par
\vspace{0.35em}
\FieldLabel{Symptom}
The model-encoder input projection from 4 × d\_model to d\_model is implemented as a depthwise-separable convolution, so the resizer performs local spatial mixing rather than pure per-position channel fusion.
\par
\vspace{0.35em}
\FieldLabel{Root Cause}
The concatenated attention tensor {[}$c, a, c \odot a, c \odot b${]} only needs a per-position channel reduction before entering the model encoder. A pointwise \texttt{Conv1d} performs that channel fusion; a wider separable convolution additionally mixes neighbouring positions, which changes the role of the layer from channel projection to unintended local feature extraction.
\par
\vspace{0.35em}
\FieldLabel{Corrective Action Taken}
Replace the depthwise-separable resizer with a pointwise Conv1d of kernel size 1.
\par
\vspace{0.35em}
\begin{Verbatim}[commandchars=\\\{\},fontsize=\footnotesize,breaklines=true,breaksymbol={},frame=leftline,framerule=2pt,rulecolor=\color{gray!40},xleftmargin=6pt,formatcom=\singlespacing]
\textcolor{gray}{# Models/qanet.py L58}
\textcolor{BrickRed}{- self.cq_resizer = DepthwiseSeparableConv(d_model * 4, d_model, 5, init_name=init_name)}
\textcolor{OliveGreen}{+ self.cq_resizer = Conv1d(d_model * 4, d_model, kernel_size=1, bias=False)}
\textcolor{OliveGreen}{+ initializations[init_name](self.cq_resizer.weight)}
\end{Verbatim}
\vspace{0.22em}
\FieldLabel{Impact on Training}
\textbf{Pre-fix:} The fusion output is supposed to compress \([c, a, c\odot a, c\odot b]\) at each position before the model encoder starts its own contextual reasoning. If this resizer already performs local convolution, then neighbouring positions are mixed before the repeated encoder passes, so later claims about what the model encoder is learning are confounded by work that has already happened one stage too early.\par\smallskip
\textbf{Post-fix:} The resizer now performs only per-position channel fusion. Spatial interaction is left to the subsequent model-encoder blocks, which restores a clean causal boundary between fusion and contextual refinement.
\par
\vspace{0.35em}
\end{bugbox}

\begin{bugbox}{BUG-N012 \checkmark{} Stochastic depth implemented}
{\color{MetaGray}\small Stage: 2 \;\textbar\; Category: Deep Learning Mechanism / Regularization Techniques \;\textbar\; Source: Models/encoder.py EncoderBlock \& Models/qanet.py shared encoder/model-encoder call sites\par}
\vspace{0.22em}
\FieldLabel{Nature}
Inconsistent mechanism-level implementation flaw in stochastic-depth regularization.
\par
\vspace{0.35em}
\FieldLabel{Symptom}
Residual sublayers are regularized with ordinary element-wise dropout rather than layer-wise stochastic depth, so the training path does not follow the paper-specified survival schedule.
\par
\vspace{0.35em}
\FieldLabel{Root Cause}
Element-wise dropout suppresses individual activations within a layer, whereas stochastic depth randomly bypasses entire residual sublayers and preserves the residual path. These regularizers are not interchangeable: they alter gradient routing, effective depth, and residual-signal preservation in different ways.
\par
\vspace{0.35em}
\FieldLabel{Corrective Action Taken}
Implement layer-level stochastic depth and apply it consistently across convolution, self-attention, and FFN sublayers with depth-aware survival probabilities. This supersedes the call-site portion of BUG-N008: the shared encoder route first had to be restored, and this change then redefined that shared route to carry the intended layer-dropout schedule.
\par
\vspace{0.35em}
\begin{Verbatim}[commandchars=\\\{\},fontsize=\footnotesize,breaklines=true,breaksymbol={},frame=leftline,framerule=2pt,rulecolor=\color{gray!40},xleftmargin=6pt,formatcom=\singlespacing]
\textcolor{gray}{# Models/encoder.py}
\textcolor{BrickRed}{- self.conv_drops = nn.ModuleList([Dropout(dropout * (i + 1) / conv_num) for i in range(conv_num)])}
\textcolor{BrickRed}{- self.drop = Dropout(dropout)}
\textcolor{OliveGreen}{+ self.dropout = dropout}
\textcolor{OliveGreen}{+ self.conv_num = conv_num}
\textcolor{OliveGreen}{+}
\textcolor{OliveGreen}{+ def _layer_dropout(self, inputs: torch.Tensor, residual: torch.Tensor, drop_prob: float) -> torch.Tensor:}
\textcolor{OliveGreen}{+  if self.training and drop_prob > 0 and torch.empty(1).uniform_(0, 1).item() < drop_prob:}
\textcolor{OliveGreen}{+   return residual}
\textcolor{OliveGreen}{+  return F.dropout(inputs, p=drop_prob, training=self.training) + residual}
 
\textcolor{BrickRed}{- def forward(self, x: torch.Tensor, mask: torch.Tensor) -> torch.Tensor:}
\textcolor{OliveGreen}{+ def forward(self, x: torch.Tensor, mask: torch.Tensor, l: int = 1, total_layers: int = 0) -> torch.Tensor:}
\textcolor{OliveGreen}{+  drop_scale = self.dropout / total_layers if total_layers > 0 else 0.0}
\textcolor{BrickRed}{-  out = out + res}
\textcolor{BrickRed}{-  if (i + 1) % 2 == 0:}
\textcolor{BrickRed}{-   out = self.conv_drops[i](out)}
\textcolor{OliveGreen}{+  out = self._layer_dropout(out, res, drop_scale * l)}
\textcolor{OliveGreen}{+  l += 1}
\textcolor{BrickRed}{-  out = self.drop(out)}
\textcolor{OliveGreen}{+  out = self._layer_dropout(out, res, drop_scale * l)}
\textcolor{OliveGreen}{+  l += 1}
\textcolor{BrickRed}{-  out = self.drop(out)}
\textcolor{OliveGreen}{+  out = self._layer_dropout(out, res, drop_scale * l)}

\textcolor{gray}{# Models/qanet.py}
\textcolor{BrickRed}{- Ce = self.emb_enc(C, cmask)}
\textcolor{BrickRed}{- Qe = self.emb_enc(Q, qmask)}
\textcolor{OliveGreen}{+ emb_total = self.emb_enc.conv_num + 2}
\textcolor{OliveGreen}{+ Ce = self.emb_enc(C, cmask, l=1, total_layers=emb_total)}
\textcolor{OliveGreen}{+ Qe = self.emb_enc(Q, qmask, l=1, total_layers=emb_total)}
 
\textcolor{BrickRed}{- for enc in self.model_enc_blks:}
\textcolor{BrickRed}{-  M1 = enc(M1, cmask)}
\textcolor{OliveGreen}{+ for i, enc in enumerate(self.model_enc_blks):}
\textcolor{OliveGreen}{+  M1 = enc(M1, cmask, l=i * sublayers_per_blk + 1, total_layers=model_total)}
\textcolor{BrickRed}{- for enc in self.model_enc_blks:}
\textcolor{BrickRed}{-  M2 = enc(M2, cmask)}
\textcolor{OliveGreen}{+ for i, enc in enumerate(self.model_enc_blks):}
\textcolor{OliveGreen}{+  M2 = enc(M2, cmask, l=i * sublayers_per_blk + 1, total_layers=model_total)}
\textcolor{BrickRed}{- for enc in self.model_enc_blks:}
\textcolor{BrickRed}{-  M3 = enc(M3, cmask)}
\textcolor{OliveGreen}{+ for i, enc in enumerate(self.model_enc_blks):}
\textcolor{OliveGreen}{+  M3 = enc(M3, cmask, l=i * sublayers_per_blk + 1, total_layers=model_total)}
\end{Verbatim}
\vspace{0.22em}
\FieldLabel{Impact on Training}
\textbf{Pre-fix:} Ordinary dropout can still regularise activations, but it does not regularise the residual pathway in the same way as stochastic depth. The model therefore trains under a deterministic effective depth even if every other component is correct, so gradient flow, ensemble-like averaging, and the regularisation pressure on deep residual stacks all differ from the intended recipe.\par\smallskip
\textbf{Post-fix:} Whole residual sublayers are now skipped according to the intended depth-dependent survival rule. This restores the regularisation mechanism that controls effective depth rather than merely masking individual activations.
\par
\vspace{0.35em}
\end{bugbox}

\ReportSubsection{3.3 Normalization, Activation \& Dropout}

{\itshape
Normalization behaviour, activation correctness, and dropout mechanics.
\par}

\begin{bugbox}{BUG-004 \checkmark{} LayerNorm keepdim broadcasting fix}
{\color{MetaGray}\small Stage: 1 \;\textbar\; Category: Runtime Failure / Forward Propagation \;\textbar\; Source: Models/Normalizations/layernorm.py L37\&38\par}
\vspace{0.22em}
\FieldLabel{Nature}
Incorrect structural pipeline implementation in LayerNorm broadcasting.
\par
\vspace{0.35em}
\FieldLabel{Symptom}
Layer normalization may produce incorrect subtraction and division shapes, and in some cases trigger runtime shape problems.
\par
\vspace{0.35em}
\FieldLabel{Root Cause}
When means and variances are reduced without keepdim=True, the reduced axes disappear. Under PyTorch broadcasting semantics, singleton dimensions are needed so that subtraction and division align with the original {[}B, C, L{]} tensor. Without them, broadcasting follows the incorrect axes or fails.
\par
\vspace{0.35em}
\FieldLabel{Corrective Action Taken}
Preserve reduced dimensions during mean and variance computation by using keepdim=True.
\par
\vspace{0.35em}
\begin{Verbatim}[commandchars=\\\{\},fontsize=\footnotesize,breaklines=true,breaksymbol={},frame=leftline,framerule=2pt,rulecolor=\color{gray!40},xleftmargin=6pt,formatcom=\singlespacing]
\textcolor{gray}{# Models/Normalizations/layernorm.py L37-L38}
\textcolor{BrickRed}{- mean = x.mean(dim=self.dim, keepdim=False)}
\textcolor{BrickRed}{- var = x.var(dim=self.dim, unbiased=False, keepdim=False)}
\textcolor{OliveGreen}{+ mean = x.mean(dim=self.dim, keepdim=True)}
\textcolor{OliveGreen}{+ var = x.var(dim=self.dim, unbiased=False, keepdim=True)}
x_norm = (x - mean) / torch.sqrt(var + self.eps)
\end{Verbatim}
\vspace{0.22em}
\FieldLabel{Impact on Training}
\textbf{Pre-fix:} Normalization statistics are broadcast incorrectly, so activation scales and residual dynamics become unreliable. The affected path therefore cannot produce a valid loss, gradient, or validation signal.\par\smallskip
\textbf{Post-fix:} LayerNorm now broadcasts its statistics correctly across the original tensor shape, restoring stable normalization behaviour. Once the contract is restored, later loss and metric changes reflect model learning rather than pipeline repair.
\par
\vspace{0.35em}
\end{bugbox}

\begin{bugbox}{BUG-044 \checkmark{} LayerNorm affine order fixed}
{\color{MetaGray}\small Stage: 2 \;\textbar\; Category: Deep Learning Mechanism / Normalization \;\textbar\; Source: Models/Normalizations/layernorm.py L41\par}
\vspace{0.22em}
\FieldLabel{Nature}
Inconsistent mechanism-level implementation flaw in LayerNorm affine ordering.
\par
\vspace{0.35em}
\FieldLabel{Symptom}
The learnable affine parameters of LayerNorm are applied in the incorrect order relative to normalization, so the module no longer implements standard LayerNorm even if it runs.
\par
\vspace{0.35em}
\FieldLabel{Root Cause}
Standard LayerNorm first normalizes the tensor, then applies the learned scale γ and shift beta. If the affine parameters are applied before normalization, or composed in the incorrect algebraic order, the normalization step cancels or alters their intended effect.
\par
\vspace{0.35em}
\FieldLabel{Corrective Action Taken}
Apply the affine transformation only after the normalized activations have been computed.
\par
\vspace{0.35em}
\begin{Verbatim}[commandchars=\\\{\},fontsize=\footnotesize,breaklines=true,breaksymbol={},frame=leftline,framerule=2pt,rulecolor=\color{gray!40},xleftmargin=6pt,formatcom=\singlespacing]
\textcolor{gray}{# Models/Normalizations/layernorm.py L41}
\textcolor{BrickRed}{- return x_norm * self.bias + self.weight}
\textcolor{OliveGreen}{+ return x_norm * self.weight + self.bias}
\end{Verbatim}
\vspace{0.22em}
\FieldLabel{Impact on Training}
\textbf{Pre-fix:} The normalization layer no longer implements the intended mapping, so feature scaling and residual stability are degraded. Activation statistics and residual scaling are miscomputed, which can destabilize optimization. \par\smallskip
\textbf{Post-fix:} LayerNorm now follows the standard normalize-then-affine form, restoring correct learnable rescaling behaviour. Feature statistics are now controlled in the intended way. 
\par
\vspace{0.35em}
\end{bugbox}

\begin{bugbox}{BUG-037 \checkmark{} GroupNorm grouping reshape order fixed}
{\color{MetaGray}\small Stage: 2 \;\textbar\; Category: Deep Learning Mechanism / Normalization \;\textbar\; Source: Models/Normalizations/groupnorm.py L35\par}
\vspace{0.22em}
\FieldLabel{Nature}
Inconsistent mechanism-level implementation flaw in GroupNorm grouping semantics.
\par
\vspace{0.35em}
\FieldLabel{Symptom}
Group normalization may normalize the incorrect subsets of channels, and depending on the reshape sequence may also trigger shape inconsistencies.
\par
\vspace{0.35em}
\FieldLabel{Root Cause}
For channel-first tensors, GroupNorm conceptually reshapes {[}B, C, L{]} into {[}B, G, C/G, L{]}. If the channel and group axes are interchanged during reshape or permutation, channels are assigned to the incorrect groups. Because Python reshaping preserves contiguous order, any missing permutation before view() can silently scramble group membership.
\par
\vspace{0.35em}
\FieldLabel{Corrective Action Taken}
Reshape and order axes so that groups are formed as contiguous channel blocks of size C/G for each group.
\par
\vspace{0.35em}
\begin{Verbatim}[commandchars=\\\{\},fontsize=\footnotesize,breaklines=true,breaksymbol={},frame=leftline,framerule=2pt,rulecolor=\color{gray!40},xleftmargin=6pt,formatcom=\singlespacing]
\textcolor{gray}{# Models/Normalizations/groupnorm.py L35}
\textcolor{BrickRed}{- x = x.view(B, C // self.G, self.G, *spatial)}
\textcolor{OliveGreen}{+ x = x.view(B, self.G, C // self.G, *spatial)}
\end{Verbatim}
\vspace{0.22em}
\FieldLabel{Impact on Training}
\textbf{Pre-fix:} Channels are normalized against the wrong peers, weakening optimization stability and distorting learned feature statistics. Activation statistics and residual scaling are miscomputed, which can destabilize optimization. \par\smallskip
\textbf{Post-fix:} GroupNorm now partitions channels correctly and normalizes over the intended grouped structure. Feature statistics are now controlled in the intended way. 
\par
\vspace{0.35em}
\end{bugbox}

\begin{bugbox}{BUG-036 \checkmark{} ReLU clamp direction fixed}
{\color{MetaGray}\small Stage: 2 \;\textbar\; Category: Deep Learning Mechanism / Activation Functions \;\textbar\; Source: Models/Activations/relu.py L12\par}
\vspace{0.22em}
\FieldLabel{Nature}
Inconsistent mechanism-level implementation flaw in ReLU definition.
\par
\vspace{0.35em}
\FieldLabel{Symptom}
The ReLU implementation clamps the incorrect side of zero, so positive activations may be suppressed or negative activations may survive incorrectly.
\par
\vspace{0.35em}
\FieldLabel{Root Cause}
ReLU is defined as max(x, 0), i.e.~it clips only the negative branch. If the implementation uses the opposite comparison or incorrect clamp direction, it computes a fundamentally different function.
\par
\vspace{0.35em}
\FieldLabel{Corrective Action Taken}
Restore the standard ReLU rule max(x, 0).
\par
\vspace{0.35em}
\begin{Verbatim}[commandchars=\\\{\},fontsize=\footnotesize,breaklines=true,breaksymbol={},frame=leftline,framerule=2pt,rulecolor=\color{gray!40},xleftmargin=6pt,formatcom=\singlespacing]
\textcolor{gray}{# Models/Activations/relu.py L12}
\textcolor{BrickRed}{- return x.clamp(max=0.0)}
\textcolor{OliveGreen}{+ return x.clamp(min=0.0)}
\end{Verbatim}
\vspace{0.22em}
\FieldLabel{Impact on Training}
\textbf{Pre-fix:} Activation flow is distorted throughout the network, changing representational sparsity and gradient propagation. Signal gating and gradient flow are changed throughout the network. \par\smallskip
\textbf{Post-fix:} The activation now behaves as true ReLU, restoring the intended nonlinearity. The nonlinearity now behaves as designed. 
\par
\vspace{0.35em}
\end{bugbox}

\begin{bugbox}{BUG-035 \checkmark{} LeakyReLU branch logic fixed}
{\color{MetaGray}\small Stage: 2 \;\textbar\; Category: Deep Learning Mechanism / Activation Functions \;\textbar\; Source: Models/Activations/leakeyReLU.py L19\par}
\vspace{0.22em}
\FieldLabel{Nature}
Inconsistent mechanism-level implementation flaw in LeakyReLU branch logic.
\par
\vspace{0.35em}
\FieldLabel{Symptom}
LeakyReLU applies the incorrect branch to positive versus negative inputs, so the activation behaves like an inverted piecewise function rather than the intended leak-on-negative rule.
\par
\vspace{0.35em}
\FieldLabel{Root Cause}
LeakyReLU must apply slope 1 on x \textgreater=0 and slope α on x \textless{} 0. If the comparison sign is reversed in a where-style implementation, the positive and negative branches are swapped.
\par
\vspace{0.35em}
\FieldLabel{Corrective Action Taken}
Swap the branch logic so that only the negative side receives the leak coefficient.
\par
\vspace{0.35em}
\begin{Verbatim}[commandchars=\\\{\},fontsize=\footnotesize,breaklines=true,breaksymbol={},frame=leftline,framerule=2pt,rulecolor=\color{gray!40},xleftmargin=6pt,formatcom=\singlespacing]
\textcolor{gray}{# Models/Activations/leakeyReLU.py L19}
\textcolor{BrickRed}{- return torch.where(x < 0, x, self.negative_slope * x)}
\textcolor{OliveGreen}{+ return torch.where(x < 0, self.negative_slope * x, x)}
\end{Verbatim}
\vspace{0.22em}
\FieldLabel{Impact on Training}
\textbf{Pre-fix:} Positive activations are damped incorrectly while negative activations may pass unchanged, weakening gradient flow and altering representation geometry. Signal gating and gradient flow are changed throughout the network. \par\smallskip
\textbf{Post-fix:} LeakyReLU now behaves as intended, preserving the positive branch and leaking only on the negative branch. The nonlinearity now behaves as designed. 
\par
\vspace{0.35em}
\end{bugbox}

\begin{bugbox}{BUG-048 \checkmark{} Inverted dropout scaling fixed}
{\color{MetaGray}\small Stage: 2 \;\textbar\; Category: Deep Learning Mechanism / Regularization Techniques \;\textbar\; Source: Models/dropout.py L17\par}
\vspace{0.22em}
\FieldLabel{Nature}
Inconsistent mechanism-level implementation flaw in inverted-dropout scaling.
\par
\vspace{0.35em}
\FieldLabel{Symptom}
Dropout scales activations incorrectly, so activation magnitudes shift between train and test time.
\par
\vspace{0.35em}
\FieldLabel{Root Cause}
In inverted dropout, activations kept during training must be divided by the keep probability q so that the expected activation magnitude matches test-time inference. If scaling is omitted, reversed, or applied at test time instead, the network sees inconsistent feature scales across modes.
\par
\vspace{0.35em}
\FieldLabel{Corrective Action Taken}
Use correct inverted-dropout scaling on the kept activations during training.
\par
\vspace{0.35em}
\begin{Verbatim}[commandchars=\\\{\},fontsize=\footnotesize,breaklines=true,breaksymbol={},frame=leftline,framerule=2pt,rulecolor=\color{gray!40},xleftmargin=6pt,formatcom=\singlespacing]
\textcolor{gray}{# Models/dropout.py L17}
\textcolor{BrickRed}{- return x * mask / self.p}
\textcolor{OliveGreen}{+ return x * mask / (1.0 - self.p)}
\end{Verbatim}
\vspace{0.22em}
\FieldLabel{Impact on Training and Evaluation}
\textbf{Pre-fix:} Activation magnitude differs systematically between training and inference, weakening optimization stability and calibration. Regularization is applied to the wrong quantity or with the wrong strength, so the model can overfit or suppress the wrong signal. \par\smallskip
\textbf{Post-fix:} Expected activation scale is now consistent across training and evaluation, as required by inverted dropout. 
\par
\vspace{0.35em}
\end{bugbox}

\ReportSubsection{3.4 Weight Initialization}

{\itshape
Parameter initialization formulas and their effect on optimization stability.
\par}

\begin{bugbox}{BUG-041 \checkmark{} Kaiming normal formula fixed}
{\color{MetaGray}\small Stage: 2 \;\textbar\; Category: Deep Learning Mechanism / Initialization \;\textbar\; Source: Models/Initializations/kaiming.py L25\par}
\vspace{0.22em}
\FieldLabel{Nature}
Inconsistent mechanism-level implementation flaw in Kaiming-normal initialization variance.
\par
\vspace{0.35em}
\FieldLabel{Symptom}
Kaiming normal initialization uses the incorrect variance, which may make activations or gradients explode or vanish early in training.
\par
\vspace{0.35em}
\FieldLabel{Root Cause}
Kaiming normal initialization is derived from variance preservation under ReLU-type nonlinearities and depends on the intended fan-in scaling. If the standard deviation formula is incorrect, the intended variance propagation argument no longer holds.
\par
\vspace{0.35em}
\FieldLabel{Corrective Action Taken}
Restore the standard Kaiming normal variance formula using the correct fan-in dependence.
\par
\vspace{0.35em}
\begin{Verbatim}[commandchars=\\\{\},fontsize=\footnotesize,breaklines=true,breaksymbol={},frame=leftline,framerule=2pt,rulecolor=\color{gray!40},xleftmargin=6pt,formatcom=\singlespacing]
\textcolor{gray}{# Models/Initializations/kaiming.py L25}
\textcolor{BrickRed}{- std = math.sqrt(1.0 / fan)}
\textcolor{OliveGreen}{+ std = math.sqrt(2.0 / fan)}
tensor.normal_(0.0, std)
\end{Verbatim}
\vspace{0.22em}
\FieldLabel{Impact on Training}
\textbf{Pre-fix:} Early-layer signal magnitudes can be badly scaled, making optimization unstable from the first updates. The initial parameter scale no longer matches the intended variance-preservation rule, so early optimization is mis-calibrated. \par\smallskip
\textbf{Post-fix:} Kaiming normal now preserves activation scale in the intended way for ReLU-like networks. Initial parameter statistics now follow the intended variance-preservation rule, improving early optimization stability. 
\par
\vspace{0.35em}
\end{bugbox}

\begin{bugbox}{BUG-042 \checkmark{} Kaiming uniform formula fixed}
{\color{MetaGray}\small Stage: 2 \;\textbar\; Category: Deep Learning Mechanism / Initialization \;\textbar\; Source: Models/Initializations/kaiming.py L38\par}
\vspace{0.22em}
\FieldLabel{Nature}
Inconsistent mechanism-level implementation flaw in Kaiming-uniform initialization bound.
\par
\vspace{0.35em}
\FieldLabel{Symptom}
Kaiming uniform initialization samples from the incorrect interval, causing parameter scales that are too large or too small and destabilizing early optimization.
\par
\vspace{0.35em}
\FieldLabel{Root Cause}
Uniform initializers are determined by converting the target variance into a symmetric interval {[}-a, a{]}. If the relationship between a and the desired variance is incorrect, the sampled parameters do not match the intended scale.
\par
\vspace{0.35em}
\FieldLabel{Corrective Action Taken}
Restore the correct bound for Kaiming uniform initialization from the proper variance-preservation derivation.
\par
\vspace{0.35em}
\begin{Verbatim}[commandchars=\\\{\},fontsize=\footnotesize,breaklines=true,breaksymbol={},frame=leftline,framerule=2pt,rulecolor=\color{gray!40},xleftmargin=6pt,formatcom=\singlespacing]
\textcolor{gray}{# Models/Initializations/kaiming.py L38}
\textcolor{BrickRed}{- std = math.sqrt(1.0 / fan)}
\textcolor{OliveGreen}{+ std = math.sqrt(2.0 / fan)}
bound = math.sqrt(3.0) * std
\end{Verbatim}
\vspace{0.22em}
\FieldLabel{Impact on Training}
\textbf{Pre-fix:} The model begins training from a mis-scaled parameter distribution, weakening optimization stability and slowing convergence. The initial parameter scale no longer matches the intended variance-preservation rule, so early optimization is mis-calibrated. \par\smallskip
\textbf{Post-fix:} Kaiming uniform now samples from the correct interval and produces the intended initial variance. Initial parameter statistics now follow the intended variance-preservation rule, improving early optimization stability. 
\par
\vspace{0.35em}
\end{bugbox}

\begin{bugbox}{BUG-043 \checkmark{} Xavier normal denominator fixed}
{\color{MetaGray}\small Stage: 2 \;\textbar\; Category: Deep Learning Mechanism / Initialization \;\textbar\; Source: Models/Initializations/xavier.py L24\par}
\vspace{0.22em}
\FieldLabel{Nature}
Inconsistent mechanism-level implementation flaw in Xavier-normal initialization denominator.
\par
\vspace{0.35em}
\FieldLabel{Symptom}
Xavier normal initialization uses the incorrect denominator, so parameter variance no longer matches the intended forward/backward compromise.
\par
\vspace{0.35em}
\FieldLabel{Root Cause}
Xavier initialization is based on balancing fan-in and fan-out. If the denominator omits one of them or combines them incorrectly, the initializer no longer preserves variance in the intended average sense.
\par
\vspace{0.35em}
\FieldLabel{Corrective Action Taken}
Use the correct Xavier normal denominator involving both fan-in and fan-out.
\par
\vspace{0.35em}
\begin{Verbatim}[commandchars=\\\{\},fontsize=\footnotesize,breaklines=true,breaksymbol={},frame=leftline,framerule=2pt,rulecolor=\color{gray!40},xleftmargin=6pt,formatcom=\singlespacing]
\textcolor{gray}{# Models/Initializations/xavier.py L24}
\textcolor{BrickRed}{- std = math.sqrt(2.0 / (fan_in * fan_out))}
\textcolor{OliveGreen}{+ std = math.sqrt(2.0 / (fan_in + fan_out))}
tensor.normal_(0.0, std)
\end{Verbatim}
\vspace{0.22em}
\FieldLabel{Impact on Training}
\textbf{Pre-fix:} Initial activations and backpropagated gradients can be mis-scaled, reducing training stability and slowing progress. The initial parameter scale no longer matches the intended variance-preservation rule, so early optimization is mis-calibrated. \par\smallskip
\textbf{Post-fix:} Xavier normal now implements the correct fan-in/fan-out balance, improving initial signal propagation. Initial parameter statistics now follow the intended variance-preservation rule, improving early optimization stability. 
\par
\vspace{0.35em}
\end{bugbox}

\begin{bugbox}{BUG-062 \checkmark{} Xavier uniform denominator fixed}
{\color{MetaGray}\small Stage: 2 \;\textbar\; Category: Deep Learning Mechanism / Initialization \;\textbar\; Source: Models/Initializations/xavier.py L36\par}
\vspace{0.22em}
\FieldLabel{Nature}
Inconsistent mechanism-level implementation flaw in Xavier-uniform initialization denominator.
\par
\vspace{0.35em}
\FieldLabel{Symptom}
Xavier uniform initialization samples from the incorrect interval because the underlying denominator is incorrect.
\par
\vspace{0.35em}
\FieldLabel{Root Cause}
As with Xavier normal, the required bound for Xavier uniform must be derived from the combined fan-in/fan-out variance target. Using the incorrect denominator produces the incorrect interval width.
\par
\vspace{0.35em}
\FieldLabel{Corrective Action Taken}
Restore the correct Xavier uniform denominator and resulting symmetric sampling bound.
\par
\vspace{0.35em}
\begin{Verbatim}[commandchars=\\\{\},fontsize=\footnotesize,breaklines=true,breaksymbol={},frame=leftline,framerule=2pt,rulecolor=\color{gray!40},xleftmargin=6pt,formatcom=\singlespacing]
\textcolor{gray}{# Models/Initializations/xavier.py L36}
\textcolor{BrickRed}{- std = math.sqrt(2.0 / (fan_in * fan_out))}
\textcolor{OliveGreen}{+ std = math.sqrt(2.0 / (fan_in + fan_out))}
bound = math.sqrt(3.0) * std
\end{Verbatim}
\vspace{0.22em}
\FieldLabel{Impact on Training}
\textbf{Pre-fix:} Parameters start from a distribution that does not preserve signal scale as intended, weakening early optimization behaviour. The initial parameter scale no longer matches the intended variance-preservation rule, so early optimization is mis-calibrated. \par\smallskip
\textbf{Post-fix:} Xavier uniform now samples from the correct interval and restores the intended variance compromise. Initial parameter statistics now follow the intended variance-preservation rule, improving early optimization stability. 
\par
\vspace{0.35em}
\end{bugbox}

\ReportSection{4 Context-Query Attention: Bug Repairs}

\ReportSubsection{4.1 Attention Mechanisms}

{\itshape
Multi-head self-attention and context-query attention.
\par}

\begin{bugbox}{BUG-005 \checkmark{} CQAttention batched-matmul order fixed}
{\color{MetaGray}\small Stage: 1 \;\textbar\; Category: Runtime Failure / Forward Propagation \;\textbar\; Source: Models/attention.py L38\par}
\vspace{0.22em}
\FieldLabel{Nature}
Incorrect structural pipeline implementation in context-query attention batched multiplication.
\par
\vspace{0.35em}
\FieldLabel{Symptom}
torch.bmm is called with incompatible operand order, producing a size-mismatch runtime failure in context-query attention.
\par
\vspace{0.35em}
\FieldLabel{Root Cause}
For batched matrix multiplication, inner dimensions must match. Here S1 has shape {[}B, Lc, Lq{]} and Q has shape {[}B, Lq, C{]}, so S1 @ Q is valid and yields {[}B, Lc, C{]}. Reversing the order tries to multiply {[}B, Lq, C{]} by {[}B, Lc, Lq{]}, whose inner dimensions do not align.
\par
\vspace{0.35em}
\FieldLabel{Corrective Action Taken}
Change the multiplication to torch.bmm(S1, Q).
\par
\vspace{0.35em}
\begin{Verbatim}[commandchars=\\\{\},fontsize=\footnotesize,breaklines=true,breaksymbol={},frame=leftline,framerule=2pt,rulecolor=\color{gray!40},xleftmargin=6pt,formatcom=\singlespacing]
\textcolor{gray}{# Models/attention.py L36 class CQAttention(nn.Module):}
\textcolor{BrickRed}{- A = torch.bmm(Q, S1)}
\textcolor{OliveGreen}{+ A = torch.bmm(S1, Q)}
\end{Verbatim}
\vspace{0.22em}
\FieldLabel{Impact on Training}
\textbf{Pre-fix:} Context-query attention fails at runtime because batched matrix multiplication dimensions do not align, so fused attention features cannot be constructed. The affected path therefore cannot produce a valid loss, gradient, or validation signal.\par\smallskip
\textbf{Post-fix:} The dimensions now align correctly, allowing context-query attention to compute valid aligned representations and continue execution. Once the contract is restored, later loss and metric changes reflect model learning rather than pipeline repair.
\par
\vspace{0.35em}
\end{bugbox}

\begin{bugbox}{BUG-051 \checkmark{} Context/query masks passed in correct order}
{\color{MetaGray}\small Stage: 1 \;\textbar\; Category: Runtime Failure / Forward Propagation \;\textbar\; Source: Models/qanet.py L82\par}
\vspace{0.22em}
\FieldLabel{Nature}
Incorrect structural pipeline implementation in attention-mask routing.
\par
\vspace{0.35em}
\FieldLabel{Symptom}
Context-query attention receives the incorrect masks: context padding tokens are treated as query padding tokens and vice versa. Under standard limits, this also creates a runtime shape mismatch.
\par
\vspace{0.35em}
\FieldLabel{Root Cause}
The CQAttention interface expects arguments in the order (C, Q, cmask, qmask). Passing them in the opposite order means the mask shapes are no longer aligned with the similarity matrix axes they are supposed to cover.
\par
\vspace{0.35em}
\FieldLabel{Corrective Action Taken}
Pass cmask and qmask in the correct order.
\par
\vspace{0.35em}
\begin{Verbatim}[commandchars=\\\{\},fontsize=\footnotesize,breaklines=true,breaksymbol={},frame=leftline,framerule=2pt,rulecolor=\color{gray!40},xleftmargin=6pt,formatcom=\singlespacing]
\textcolor{gray}{# Models/qanet.py L73 class QANet(nn.Module):}
\textcolor{BrickRed}{- X = self.cq_att(Ce, Qe, qmask, cmask)}
\textcolor{OliveGreen}{+ X = self.cq_att(Ce, Qe, cmask, qmask)}
\end{Verbatim}
\vspace{0.22em}
\FieldLabel{Impact on Training}
\textbf{Pre-fix:} Masked attention is computed against the wrong padding structure, and under the standard para\_limit/ques\_limit setting this also creates a shape-mismatch crash. If the mismatch becomes explicit, execution stops before a valid loss or metric can be produced. If it remains latent because dimensions stay compatible, downstream attention, optimization, and evaluation are driven by the incorrect semantic signal.\par\smallskip
\textbf{Post-fix:} Context-query attention now receives correctly aligned context and question masks, restoring valid masked-attention computation. Later layers now receive the intended signal, so downstream loss and span metrics become meaningful evidence of learning.
\par
\vspace{0.35em}
\end{bugbox}

\begin{bugbox}{BUG-038/050 \checkmark{} Self-attention output no longer discarded}
{\color{MetaGray}\small Stage: 2 \;\textbar\; Category: Deep Learning Mechanism / Attention-Based Modules \;\textbar\; Source: Models/encoder.py L117\par}
\vspace{0.22em}
\FieldLabel{Nature}
Inconsistent mechanism-level implementation flaw in self-attention residual integration.
\par
\vspace{0.35em}
\FieldLabel{Symptom}
Self-attention output is completely discarded, so the encoder block's attention sublayer contributes nothing.
\par
\vspace{0.35em}
\FieldLabel{Root Cause}
Residual connections must add the transformed sublayer output back to the residual stream. Instead, the attention output is overwritten by the residual tensor itself, erasing the attention computation entirely.
\par
\vspace{0.35em}
\FieldLabel{Corrective Action Taken}
Use a residual addition out = out + res rather than overwriting with res.
\par
\vspace{0.35em}
\begin{Verbatim}[commandchars=\\\{\},fontsize=\footnotesize,breaklines=true,breaksymbol={},frame=leftline,framerule=2pt,rulecolor=\color{gray!40},xleftmargin=6pt,formatcom=\singlespacing]
\textcolor{gray}{# Models/encoder.py L122 class EncoderBlock(nn.Module):}
\textcolor{BrickRed}{- out = res}
\textcolor{OliveGreen}{+ out = out + res}
\end{Verbatim}
\vspace{0.22em}
\FieldLabel{Impact on Training}
\textbf{Pre-fix:} The attention sublayer is effectively dead, so the encoder behaves like a conv-only network and loses its ability to model long-range dependencies. Information routing across tokens is no longer what the architecture intends, so the model may over-sharpen, ignore, or miscombine contextual evidence. \par\smallskip
\textbf{Post-fix:} Self-attention output is now correctly combined with the residual path, restoring the encoder's global interaction mechanism. Token-to-token interaction now follows the intended attention rule, restoring a defensible explanation for contextual learning and answer extraction quality. 
\par
\vspace{0.35em}
\end{bugbox}

\begin{bugbox}{BUG-N003 \checkmark{} Missing 1/√dk attention scaling restored}
{\color{MetaGray}\small Stage: 2 \;\textbar\; Category: Deep Learning Mechanism / Attention-Based Modules \;\textbar\; Source: Models/encoder.py L78\par}
\vspace{0.22em}
\FieldLabel{Nature}
Inconsistent mechanism-level implementation flaw in scaled dot-product attention.
\par
\vspace{0.35em}
\FieldLabel{Symptom}
Self-attention logits become excessively large, yielding overly peaked attention distributions and weak gradient flow through the softmax.
\par
\vspace{0.35em}
\FieldLabel{Root Cause}
In scaled dot-product attention, the $1/\sqrt{d_k}$ factor normalizes logit variance as key/query dimensionality grows. Omitting it makes the softmax saturate, so token-to-token weighting becomes brittle and the attention module no longer behaves as the intended differentiable weighting operator.
\par
\vspace{0.35em}
\FieldLabel{Corrective Action Taken}
Divide attention logits by √dk before applying softmax.
\par
\vspace{0.35em}
\begin{Verbatim}[commandchars=\\\{\},fontsize=\footnotesize,breaklines=true,breaksymbol={},frame=leftline,framerule=2pt,rulecolor=\color{gray!40},xleftmargin=6pt,formatcom=\singlespacing]
\textcolor{gray}{# Models/encoder.py L78 class MultiHeadAttention(nn.Module):}
q = self.wq(x).view(batch_size, length, self.num_heads, self.d_k).permute(2, 0, 1, 3).contiguous().view(-1, length, self.d_k)
k = self.wk(x).view(batch_size, length, self.num_heads, self.d_k).permute(2, 0, 1, 3).contiguous().view(-1, length, self.d_k)
\textcolor{BrickRed}{- attn = torch.bmm(q, k.transpose(1, 2))}
\textcolor{OliveGreen}{+ attn = torch.bmm(q, k.transpose(1, 2)) * self.scale}
attn = mask_logits(attn, mask)
\end{Verbatim}
\vspace{0.22em}
\FieldLabel{Impact on Training}
\textbf{Pre-fix:} Even if masking, projection, and tensor shapes are otherwise correct, the attention scores become too large as \(d_k\) grows. The softmax then saturates, low-probability alternatives receive almost no gradient, and later layers inherit brittle attention patterns that look confident but are poorly calibrated.\par\smallskip
\textbf{Post-fix:} The \(1/\sqrt{d_k}\) factor restores the variance control assumed by scaled dot-product attention, keeping the softmax in a trainable regime and making long-range weighting behaviour far more stable.
\par
\vspace{0.35em}
\end{bugbox}

\begin{bugbox}{BUG-049 \checkmark{} Multi-head head/batch reassembly fixed}
{\color{MetaGray}\small Stage: 2 \;\textbar\; Category: Deep Learning Mechanism / Attention-Based Modules \;\textbar\; Source: Models/encoder.py L84\par}
\vspace{0.22em}
\FieldLabel{Nature}
Incorrect structural pipeline implementation in multi-head attention reassembly.
\par
\vspace{0.35em}
\FieldLabel{Symptom}
After per-head attention is computed, the heads are reassembled incorrectly, mixing head and batch structure or producing the incorrect output layout.
\par
\vspace{0.35em}
\FieldLabel{Root Cause}
When splitting into multiple heads, tensors are typically reshaped and permuted between {[}B, L, H, d{]}, {[}B, H, L, d{]}, and flattened forms. In PyTorch, view() alone does not reorder memory; if the necessary permutation is omitted before reshaping back, head and batch information become interleaved incorrectly.
\par
\vspace{0.35em}
\FieldLabel{Corrective Action Taken}
Permute and reshape the multi-head tensor in the correct order so the final output is reassembled as the intended channel dimension.
\par
\vspace{0.35em}
\begin{Verbatim}[commandchars=\\\{\},fontsize=\footnotesize,breaklines=true,breaksymbol={},frame=leftline,framerule=2pt,rulecolor=\color{gray!40},xleftmargin=6pt,formatcom=\singlespacing]
\textcolor{gray}{# Models/encoder.py L84 class MultiHeadAttention(nn.Module):}
attn = torch.bmm(p_attn, v)
\textcolor{BrickRed}{- attn = attn.view(batch_size, self.num_heads, length, self.d_k)}
\textcolor{OliveGreen}{+ attn = attn.view(self.num_heads, batch_size, length, self.d_k)}
attn = attn.permute(1, 2, 0, 3).contiguous().view(batch_size, length, -1)
\end{Verbatim}
\vspace{0.22em}
\FieldLabel{Impact on Training}
\textbf{Pre-fix:} Head-specific information is mixed incorrectly during reassembly, so the attended representation no longer matches the intended multi-head computation. If the reshaped tensor becomes incompatible, attention execution fails outright; if it remains shape-compatible, the model still reassembles heads incorrectly and propagates semantically wrong contextual features. That makes both optimization and later evaluation hard to trust, because the model may appear to run while learning from a broken attention representation.\par\smallskip
\textbf{Post-fix:} The attention output is now reconstructed from heads in the correct order, restoring proper multi-head behaviour. With head and batch structure reassembled correctly, attention outputs once again correspond to the intended multi-head computation, which makes downstream learning dynamics and validation behaviour interpretable.
\par
\vspace{0.35em}
\end{bugbox}

\ReportSection{5 Model Encoder: Bug Repairs}

\ReportSubsection{5.1 Model Encoder}

The model encoder reuses the same encoder-block machinery across three passes. The bugs most directly affecting this stage are already accounted for above: BUG-N006 (FFN structure), BUG-N010 (model-encoder input projection), BUG-N012 (stochastic depth), BUG-038/050 (self-attention residual), BUG-N003 (attention scaling), and the shared convolution / normalization fixes in the embedding-encoder section. No additional distinct bug IDs are introduced only for the model encoder beyond those already counted in the 56-bug set.

\ReportSection{6 Output Layer: Bug Repairs}

\ReportSubsection{6.1 Output Layer \& Loss Function}

{\itshape
Pointer head, NLL/CE loss interface, and loss scaling.
\par}

\begin{bugbox}{BUG-003 \checkmark{} NLL loss argument order fixed}
{\color{MetaGray}\small Stage: 1 \;\textbar\; Category: Runtime Failure / Gradient Computation \;\textbar\; Source: Losses/loss.py L6\par}
\vspace{0.22em}
\FieldLabel{Nature}
Incorrect structural pipeline implementation in loss-function input/target ordering.
\par
\vspace{0.35em}
\FieldLabel{Symptom}
F.nll\_loss receives the target tensor as input and the log-probability tensor as target, causing a shape/dtype mismatch runtime failure.
\par
\vspace{0.35em}
\FieldLabel{Root Cause}
PyTorch loss functions have a fixed signature: the first argument is the model output and the second is the supervision target. Swapping them violates both the expected tensor shapes and the expected dtypes.
\par
\vspace{0.35em}
\FieldLabel{Corrective Action Taken}
Pass the log-probabilities first and the target indices second.
\par
\vspace{0.35em}
\begin{Verbatim}[commandchars=\\\{\},fontsize=\footnotesize,breaklines=true,breaksymbol={},frame=leftline,framerule=2pt,rulecolor=\color{gray!40},xleftmargin=6pt,formatcom=\singlespacing]
\textcolor{gray}{# Losses/loss.py L5 def qa_nll_loss(p1, p2, y1, y2):}
\textcolor{BrickRed}{- return 0.5 * (F.nll_loss(y1, p1) + F.nll_loss(p2, y2))}
\textcolor{OliveGreen}{+ return 0.5 * (F.nll_loss(p1, y1) + F.nll_loss(p2, y2))}
\end{Verbatim}
\vspace{0.22em}
\FieldLabel{Impact on Training}
\textbf{Pre-fix:} Training loss computation fails because F.nll\_loss receives invalid input/target types and shapes, so optimization cannot begin. Without a valid loss/backward contract, gradients cannot be computed in the mathematically intended way, so the optimizer never receives a usable training signal and the model cannot demonstrate decreasing loss for the right reason. Empirically, this is a direct training blocker rather than a subtle quality issue.\par\smallskip
\textbf{Post-fix:} The NLL loss now receives valid (log\_probs, target\_indices) arguments, enabling correct span-loss computation and stable backpropagation. Autograd now receives the correct objects and can propagate the intended supervision signal through the network, so subsequent optimization behaviour becomes meaningful rather than vacuous.
\par
\vspace{0.35em}
\end{bugbox}

\begin{bugbox}{BUG-020 \checkmark{} Pointer-head concatenation dimension fixed}
{\color{MetaGray}\small Stage: 1 \;\textbar\; Category: Runtime Failure / Forward Propagation \;\textbar\; Source: Models/heads.py L23\par}
\vspace{0.22em}
\FieldLabel{Nature}
Incorrect structural pipeline implementation in pointer-head feature concatenation.
\par
\vspace{0.35em}
\FieldLabel{Symptom}
torch.cat({[}M1, M2{]}, dim=0) concatenates along batch dimension, producing {[}2B, C, L{]} instead of {[}B, 2C, L{]}.
\par
\vspace{0.35em}
\FieldLabel{Root Cause}
The pointer head needs to concatenate features along the channel axis so that each token position receives the combined representation. In a channel-first tensor {[}B, C, L{]}, that axis is dimension 1, not dimension 0. Concatenating on dimension 0 doubles the batch size instead.
\par
\vspace{0.35em}
\FieldLabel{Corrective Action Taken}
Concatenate along dimension 1 rather than dimension 0.
\par
\vspace{0.35em}
\begin{Verbatim}[commandchars=\\\{\},fontsize=\footnotesize,breaklines=true,breaksymbol={},frame=leftline,framerule=2pt,rulecolor=\color{gray!40},xleftmargin=6pt,formatcom=\singlespacing]
\textcolor{gray}{# Models/heads.py L21 class Pointer(nn.Module):}
\textcolor{BrickRed}{- X1 = torch.cat([M1, M2], dim=0)}
\textcolor{OliveGreen}{+ X1 = torch.cat([M1, M2], dim=1)}
\end{Verbatim}
\vspace{0.22em}
\FieldLabel{Impact on Training}
\textbf{Pre-fix:} Pointer-head logits cannot be computed because the batch and channel axes are mixed, so start/end span prediction fails at runtime. The affected path therefore cannot produce a valid loss, gradient, or validation signal.\par\smallskip
\textbf{Post-fix:} Concatenation now preserves the batch axis and doubles the channel axis as intended, allowing valid pointer logits to be produced. Once the contract is restored, later loss and metric changes reflect model learning rather than pipeline repair.
\par
\vspace{0.35em}
\end{bugbox}

\begin{bugbox}{BUG-N007 \checkmark{} Extra 0.5 loss scaling removed}
{\color{MetaGray}\small Stage: 2 \;\textbar\; Category: Deep Learning Mechanism / Loss Function \;\textbar\; Source: Losses/loss.py L7\par}
\vspace{0.22em}
\FieldLabel{Nature}
Inconsistent mechanism-level implementation flaw in span-loss scaling.
\par
\vspace{0.35em}
\FieldLabel{Symptom}
The span objective is scaled by an additional factor of 0.5, so gradients reflect a different objective magnitude from the intended start-plus-end loss.
\par
\vspace{0.35em}
\FieldLabel{Root Cause}
The QANet span loss is the direct sum of the start and end negative log-likelihood terms. Multiplying that sum by 0.5 rescales every gradient and therefore changes the optimizer's effective step size, especially in non-adaptive settings.
\par
\vspace{0.35em}
\FieldLabel{Corrective Action Taken}
Remove the extra factor of 0.5 from the loss definition.
\par
\vspace{0.35em}
\begin{Verbatim}[commandchars=\\\{\},fontsize=\footnotesize,breaklines=true,breaksymbol={},frame=leftline,framerule=2pt,rulecolor=\color{gray!40},xleftmargin=6pt,formatcom=\singlespacing]
\textcolor{gray}{# Losses/loss.py L5 def qa_nll_loss(p1, p2, y1, y2):}
\textcolor{BrickRed}{- return 0.5 * (F.nll_loss(p1, y1) + F.nll_loss(p2, y2))}

\textcolor{OliveGreen}{+ return F.nll_loss(p1, y1) + F.nll_loss(p2, y2)}
\end{Verbatim}
\vspace{0.22em}
\FieldLabel{Impact on Training}
\textbf{Pre-fix:} Every optimization step is driven by a systematically weakened supervision signal. Under SGD-like settings this behaves like an unintended learning-rate reduction, slowing loss descent and making early stopping or scheduler comparisons less trustworthy.\par\smallskip
\textbf{Post-fix:} The optimizer now receives the full intended span-loss gradient. Loss values and convergence rates therefore reflect the actual objective rather than an artificially down-scaled surrogate.
\par
\vspace{0.35em}
\end{bugbox}

\begin{bugbox}{BUG-N016 \checkmark{} Output/loss interface aligned around raw logits}
{\color{MetaGray}\small Stage: 2 \;\textbar\; Category: Deep Learning Mechanism / Loss Function \;\textbar\; Source: Models/heads.py L29-30 \& Losses/loss.py L6-7 \& EvaluateTools/eval\_utils.py L107-113\par}
\vspace{0.22em}
\FieldLabel{Nature}
Incorrect structural interface and inconsistent mechanism-level implementation flaw in output/loss interface consistency.
\par
\vspace{0.35em}
\FieldLabel{Symptom}
The Pointer head emits log-probabilities, whereas qa\_ce\_loss expects raw logits; the CE path therefore applies \texttt{log\_softmax} twice and optimizes the wrong gradient.
\par
\vspace{0.35em}
\FieldLabel{Root Cause}
Probability normalization must occur in exactly one place. If the head already applies \texttt{log\_softmax}, then losses expecting raw logits will normalize again internally, which breaks the mathematical contract between the output layer, the selected loss, and the evaluation decoder.
\par
\vspace{0.35em}
\FieldLabel{Corrective Action Taken}
Make the pointer head return masked raw logits, move log\_softmax into the NLL loss, and apply log-prob conversion explicitly in evaluation decoding. This supersedes BUG-N007 inside the loss path: the earlier fix restored the correct loss magnitude, and this later interface repair restored the correct raw-logit versus log-probability contract.
\par
\vspace{0.35em}
\begin{Verbatim}[commandchars=\\\{\},fontsize=\footnotesize,breaklines=true,breaksymbol={},frame=leftline,framerule=2pt,rulecolor=\color{gray!40},xleftmargin=6pt,formatcom=\singlespacing]
\textcolor{gray}{# EvaluateTools/eval_utils.py L107-113}
\textcolor{BrickRed}{- outer = p1.unsqueeze(2) + p2.unsqueeze(1)}
\textcolor{OliveGreen}{+ p1_log = torch.nn.functional.log_softmax(p1, dim=1)}
\textcolor{OliveGreen}{+ p2_log = torch.nn.functional.log_softmax(p2, dim=1)}
\textcolor{OliveGreen}{+ outer = p1_log.unsqueeze(2) + p2_log.unsqueeze(1)}
\textcolor{gray}{# Losses/loss.py L6-7}
\textcolor{BrickRed}{- return F.nll_loss(p1, y1) + F.nll_loss(p2, y2)}
\textcolor{OliveGreen}{+ return F.nll_loss(F.log_softmax(p1, dim=1), y1) + F.nll_loss(F.log_softmax(p2, dim=1), y2)}
\textcolor{gray}{# Models/heads.py L29-30}
\textcolor{BrickRed}{- p1 = F.log_softmax(Y1, dim=1)}
\textcolor{BrickRed}{- p2 = F.log_softmax(Y2, dim=1)}
\textcolor{OliveGreen}{+ p1 = mask_logits(Y1, mask)}
\textcolor{OliveGreen}{+ p2 = mask_logits(Y2, mask)}
\end{Verbatim}
\vspace{0.22em}
\FieldLabel{Impact on Training and Evaluation}
\textbf{Pre-fix:} One registered loss path no longer optimizes the intended objective because the same normalization is applied once in the head and again in the loss. Training may still show decreasing numbers, but those gradients are not interpretable as raw-logit cross-entropy updates, and evaluation decoding is tied to inconsistent tensor semantics.\par\smallskip
\textbf{Post-fix:} The head now emits masked raw logits, while each loss and decoder applies normalization exactly where required. Gradient flow, decoded spans, and reported metrics therefore reflect a single consistent output contract.
\par
\vspace{0.35em}
\end{bugbox}

\ReportSection{7 Training Details: Bug Repairs}

\ReportSubsection{7.1 Evaluation \& Inference}

{\itshape
Checkpoint loading, span decoding, evaluation semantics, and checkpoint-configuration consistency.
\par}

\begin{bugbox}{BUG-001/002 \checkmark{} Checkpoint load key fixed to model\_state}
{\color{MetaGray}\small Stage: 1 \;\textbar\; Category: Runtime Failure / Checkpoint Management \;\textbar\; Source: EvaluateTools/evaluate.py L119\par}
\vspace{0.22em}
\FieldLabel{Nature}
Incorrect structural pipeline implementation in checkpoint state restoration.
\par
\vspace{0.35em}
\FieldLabel{Symptom}
KeyError: `model' occurs when loading a checkpoint because the saved key is `model\_state'.
\par
\vspace{0.35em}
\FieldLabel{Root Cause}
The save path and load path disagree about the dictionary key used for model parameters. Since checkpoint dictionaries are plain Python mappings, key names must match exactly across write and restore code.
\par
\vspace{0.35em}
\FieldLabel{Corrective Action Taken}
Restore the model from ckpt{[}``model\_state''{]} rather than ckpt{[}``model''{]}.
\par
\vspace{0.35em}
\begin{Verbatim}[commandchars=\\\{\},fontsize=\footnotesize,breaklines=true,breaksymbol={},frame=leftline,framerule=2pt,rulecolor=\color{gray!40},xleftmargin=6pt,formatcom=\singlespacing]
\textcolor{gray}{# EvaluateTools/evaluate.py L119}
\textcolor{BrickRed}{- model.load_state_dict(ckpt["model"])}
\textcolor{OliveGreen}{+ model.load_state_dict(ckpt["model_state"])}
\end{Verbatim}
\vspace{0.22em}
\FieldLabel{Impact on Evaluation}
\textbf{Pre-fix:} Evaluation cannot restore trained model weights from saved checkpoints, so the evaluation pipeline fails before any metrics can be computed. This is critical because the model may have learned useful parameters during training, but evaluation cannot access them, so no trustworthy validation or test conclusion can be drawn from the saved artifact.\par\smallskip
\textbf{Post-fix:} Checkpoint key names are now consistent across save/load paths, enabling successful state restoration and normal evaluation metric computation. Checkpoint restoration now refers to the intended model state, so reported metrics once again correspond to the trained parameters rather than to a load-time failure.
\par
\vspace{0.35em}
\end{bugbox}

\begin{notebugbox}{BUG-N002 $\triangle$ torch.load(..., weights\_only=False) compatibility note}
{\color{MetaGray}\small Stage: 1 \;\textbar\; Category: Runtime Failure / Checkpoint Management \;\textbar\; Source: EvaluateTools/evaluate.py L118\par}
\vspace{0.22em}
\FieldLabel{Nature}
Documented structural pipeline warning in checkpoint loading compatibility.
\par
\vspace{0.35em}
\FieldLabel{Symptom}
On PyTorch 2.6+, checkpoint loading may fail because the default behaviour of torch.load changed.
\par
\vspace{0.35em}
\FieldLabel{Root Cause}
PyTorch changed the default of weights\_only. When external library defaults change, code that previously relied on implicit behaviour can fail even though the local program logic has not changed.
\par
\vspace{0.35em}
\FieldLabel{Corrective Action Taken}
Set weights\_only=False explicitly for trusted internal checkpoints in the target checkpointing workflow.
\par
\vspace{0.35em}
\begin{Verbatim}[commandchars=\\\{\},fontsize=\footnotesize,breaklines=true,breaksymbol={},frame=leftline,framerule=2pt,rulecolor=\color{gray!40},xleftmargin=6pt,formatcom=\singlespacing]
\textcolor{gray}{# EvaluateTools/evaluate.py L118}
\textcolor{BrickRed}{- ckpt = torch.load(ckpt_path, map_location=DEVICE)}
\textcolor{OliveGreen}{+ ckpt = torch.load(ckpt_path, map_location=DEVICE, weights_only=False)}
\end{Verbatim}
\vspace{0.22em}
\FieldLabel{Impact on Evaluation}
\textbf{Pre-fix:} If unhandled on newer environments, evaluation or checkpoint restoration can fail even when training artifacts are otherwise valid. The main risk is evaluation fragility rather than model mathematics: a valid experiment artifact can become unreadable solely because the environment changed underneath the code. That undermines reproducibility and makes cross-machine validation brittle.\par\smallskip
\textbf{Post-fix:} The load path becomes robust to the PyTorch default change. This is an environment compatibility issue rather than a core model-logic bug. The explicit flag removes that environmental ambiguity, so checkpoint loading behaviour is stable across the target software versions without changing the model itself.
\par
\vspace{0.35em}
\end{notebugbox}

\begin{bugbox}{BUG-034 \checkmark{} Evaluation argmax dimension fixed}
{\color{MetaGray}\small Stage: 1 \;\textbar\; Category: Runtime Failure / Performance Evaluation \;\textbar\; Source: EvaluateTools/eval\_utils.py L107-108\par}
\vspace{0.22em}
\FieldLabel{Nature}
Incorrect structural pipeline implementation in prediction-axis reduction during evaluation.
\par
\vspace{0.35em}
\FieldLabel{Symptom}
Span prediction indices are reduced over the batch axis, producing incorrect per-sample answers and unreliable F1/EM reporting.
\par
\vspace{0.35em}
\FieldLabel{Root Cause}
For tensors shaped {[}B, L{]}, per-example span prediction requires reducing over token dimension L, not batch dimension B. Choosing dim=0 aggregates across different examples instead of across different positions within one example.
\par
\vspace{0.35em}
\FieldLabel{Corrective Action Taken}
Reduce over dim=1 rather than dim=0.
\par
\vspace{0.35em}
\begin{Verbatim}[commandchars=\\\{\},fontsize=\footnotesize,breaklines=true,breaksymbol={},frame=leftline,framerule=2pt,rulecolor=\color{gray!40},xleftmargin=6pt,formatcom=\singlespacing]
\textcolor{gray}{# EvaluateTools/eval_utils.py L107-108}
\textcolor{BrickRed}{- yp1 = torch.argmax(p1, dim=0)}
\textcolor{BrickRed}{- yp2 = torch.argmax(p2, dim=0)}
\textcolor{OliveGreen}{+ yp1 = torch.argmax(p1, dim=1)}
\textcolor{OliveGreen}{+ yp2 = torch.argmax(p2, dim=1)}
\end{Verbatim}
\vspace{0.22em}
\FieldLabel{Impact on Evaluation}
\textbf{Pre-fix:} Predicted start and end positions do not correspond to per-example sequence maxima, so extracted spans and reported metrics become semantically invalid. The model may still produce scores, but the evaluation route no longer measures per-example behaviour. As a result, F1/EM can look precise while actually reflecting a reduction over the wrong axis, which makes model selection and debugging conclusions invalid.\par\smallskip
\textbf{Post-fix:} Predicted spans are now computed per sample, restoring valid token extraction and making F1/EM trustworthy for model validation. Per-example span extraction and metric computation are now aligned with the intended evaluation contract, so metric changes can be interpreted as real changes in answer quality.
\par
\vspace{0.35em}
\end{bugbox}

\begin{bugbox}{BUG-N011 \checkmark{} Joint span decoding restored}
{\color{MetaGray}\small Stage: 1 \;\textbar\; Category: Performance Evaluation / Inference Decoding \;\textbar\; Source: EvaluateTools/eval\_utils.py L107-113\par}
\vspace{0.22em}
\FieldLabel{Nature}
Inconsistent mechanism-level implementation flaw in joint span decoding.
\par
\vspace{0.35em}
\FieldLabel{Symptom}
Inference selects start and end positions independently, so the reported span need not maximize the legal joint start-end objective.
\par
\vspace{0.35em}
\FieldLabel{Root Cause}
The decoding objective is joint under the constraint $s \le e$. Independent argmax followed by swapping is not equivalent to maximizing the pairwise score matrix, because the individually best start and end positions can form a jointly suboptimal span.
\par
\vspace{0.35em}
\FieldLabel{Corrective Action Taken}
Replace independent argmax plus min/max swap with joint outer-product decoding under an upper-triangular validity mask. This supersedes BUG-034: the earlier axis fix restored per-sample reduction, and this later repair then replaced the whole decoder with the correct joint span objective.
\par
\vspace{0.35em}
\begin{Verbatim}[commandchars=\\\{\},fontsize=\footnotesize,breaklines=true,breaksymbol={},frame=leftline,framerule=2pt,rulecolor=\color{gray!40},xleftmargin=6pt,formatcom=\singlespacing]
\textcolor{gray}{# EvaluateTools/eval_utils.py L107-113}
\textcolor{BrickRed}{- yp1 = torch.argmax(p1, dim=1)}
\textcolor{BrickRed}{- yp2 = torch.argmax(p2, dim=1)}
\textcolor{BrickRed}{- yps = torch.stack([yp1, yp2], dim=1)}
\textcolor{BrickRed}{- ymin, _ = torch.min(yps, dim=1)}
\textcolor{BrickRed}{- ymax, _ = torch.max(yps, dim=1)}
\textcolor{OliveGreen}{+ outer = p1.unsqueeze(2) + p2.unsqueeze(1)}
\textcolor{OliveGreen}{+ mask = torch.triu(torch.ones(outer.size(-2), outer.size(-1), device=outer.device, dtype=torch.bool))}
\textcolor{OliveGreen}{+ outer = outer.masked_fill(~mask, float('-inf'))}
\textcolor{OliveGreen}{+ flat = outer.view(outer.size(0), -1)}
\textcolor{OliveGreen}{+ idx = torch.argmax(flat, dim=1)}
\textcolor{OliveGreen}{+ L = p1.size(1)}
\textcolor{OliveGreen}{+ ymin = idx // L}
\textcolor{OliveGreen}{+ ymax = idx % L}
\end{Verbatim}
\vspace{0.22em}
\FieldLabel{Impact on Inference}
\textbf{Pre-fix:} The scorer can be partly right while the decoder is still wrong. In particular, start and end marginals may each rank sensible positions, but independent argmax can combine them into an illegal or jointly suboptimal pair, so reported EM/F1 understate what the span scorer actually knows.\par\smallskip
\textbf{Post-fix:} Decoding now maximizes the legal joint span objective. Evaluation therefore attributes failure to the scoring model only when the joint span distribution is genuinely wrong, not when the decoder discarded a better legal pair.
\par
\vspace{0.35em}
\end{bugbox}

\begin{rednotebugbox}{BUG-N018 $\triangle$ Checkpoint config is not restored during evaluation}
{\color{MetaGray}\small Stage: 1 \;\textbar\; Category: Runtime Failure / Checkpoint Management \;\textbar\; Source: EvaluateTools/evaluate.py\par}
\vspace{0.22em}
\FieldLabel{Nature}
Documented structural pipeline warning in evaluation-time configuration restoration.
\par
\vspace{0.35em}
\FieldLabel{Symptom}
evaluate() rebuilds the model from function defaults and \texttt{getattr} fallbacks rather than from the configuration saved in the checkpoint. If training used non-default architecture parameters, evaluation may silently construct a different model.
\par
\vspace{0.35em}
\FieldLabel{Root Cause}
The evaluation path reconstructs an \texttt{args} object that omits several architecture fields. Because the constructor then fills them with defaults, the inconsistency is hidden rather than surfaced explicitly, turning a hard configuration contract into a silent fallback path.
\par
\vspace{0.35em}
\FieldLabel{Corrective Action Taken}
Recommended patch (documented, not applied): load ckpt{[}``config''{]} before model construction and override architecture-related fields accordingly.
\par
\vspace{0.35em}
\begin{Verbatim}[commandchars=\\\{\},fontsize=\footnotesize,breaklines=true,breaksymbol={},frame=leftline,framerule=2pt,rulecolor=\color{gray!40},xleftmargin=6pt,formatcom=\singlespacing]
\textcolor{gray}{# EvaluateTools/evaluate.py (recommended patch, not applied)}
ckpt = torch.load(ckpt_path, map_location=DEVICE, weights_only=False)
\textcolor{OliveGreen}{+ ckpt_cfg = ckpt.get("config", \{\})}
\textcolor{OliveGreen}{+ for key in ["norm_name", "activation", "init_name", "norm_groups"]:}
\textcolor{OliveGreen}{+  if key in ckpt_cfg:}
\textcolor{OliveGreen}{+   setattr(args, key, ckpt_cfg[key])}
model = QANet(word_mat, char_mat, args)
\end{Verbatim}
\vspace{0.22em}
\FieldLabel{Impact on Evaluation}
\textbf{Pre-fix:} Current default-configuration runs are low-risk only because they use the constructor defaults. Once architecture knobs are varied, evaluation can report F1/EM for the incorrect model family even though no obvious runtime failure occurs, which makes mechanism-level comparisons unsafe.\par\smallskip
\textbf{Post-fix:} No code fix is currently applied in the submitted branch; this remains a documented warning. If the recommended patch is applied later, checkpoint evaluation becomes architecture-consistent and comparison claims remain defensible.
\par
\vspace{0.35em}
\end{rednotebugbox}

\ReportSubsection{7.2 Optimizers}

{\itshape
Adam state, momentum state, weight decay, squared-gradient accumulation, and bias correction.
\par}

\begin{bugbox}{BUG-022 \checkmark{} Adam state-key mismatch}
{\color{MetaGray}\small Stage: 1 \;\textbar\; Category: Runtime Failure / Parameter Update \;\textbar\; Source: Optimizers/adam.py L63\par}
\vspace{0.22em}
\FieldLabel{Nature}
Incorrect structural pipeline implementation in Adam state bookkeeping.
\par
\vspace{0.35em}
\FieldLabel{Symptom}
Adam optimizer state is read and written under inconsistent keys, which may induce a runtime lookup failure on the first update or silently destroy state persistence across steps and resumes.
\par
\vspace{0.35em}
\FieldLabel{Root Cause}
The state dictionary uses different key names at initialization and update time. Since Python dictionaries do not provide aliasing automatically, a mismatch means one part of the optimizer is reading buffers that the other part never wrote.
\par
\vspace{0.35em}
\FieldLabel{Corrective Action Taken}
Unify the Adam state dictionary so both moment buffers are initialized and read under the same names.
\par
\vspace{0.35em}
\begin{Verbatim}[commandchars=\\\{\},fontsize=\footnotesize,breaklines=true,breaksymbol={},frame=leftline,framerule=2pt,rulecolor=\color{gray!40},xleftmargin=6pt,formatcom=\singlespacing]
\textcolor{gray}{# Optimizers/adam.py L63}
\textcolor{BrickRed}{- m = state["m"]}
\textcolor{BrickRed}{- v = state["v"]}
\textcolor{OliveGreen}{+ m = state["exp_avg"]}
\textcolor{OliveGreen}{+ v = state["exp_avg_sq"]}
\end{Verbatim}
\vspace{0.22em}
\FieldLabel{Impact on Training}
\textbf{Pre-fix:} In the hard-failure case, Adam crashes before the first real parameter update. In the softer case, state persistence is effectively broken, so Adam behaves like a badly implemented noisy optimizer rather than an adaptive method with memory. Qualitatively this prevents stable curvature adaptation; quantitatively it shows up as unstable loss trajectories, poor reproducibility after resume, and slower or stalled F1/EM gains. Broken optimizer state either prevents valid updates or discards the temporal information the rule is meant to reuse. That prevents stable, interpretable learning behaviour because gradients are not translated into the intended parameter trajectory across steps or resumes.\par\smallskip
\textbf{Post-fix:} First- and second-moment history now persists across steps and checkpoints, so Adam's adaptive step sizing is coherent from warmup through late training. State reuse is now consistent across iterations.
\par
\vspace{0.35em}
\end{bugbox}

\begin{bugbox}{BUG-023 \checkmark{} SGD-momentum state-key mismatch}
{\color{MetaGray}\small Stage: 1 \;\textbar\; Category: Runtime Failure / Parameter Update \;\textbar\; Source: Optimizers/sgd\_momentum.py L49\par}
\vspace{0.22em}
\FieldLabel{Nature}
Incorrect structural pipeline implementation in SGD momentum-state bookkeeping.
\par
\vspace{0.35em}
\FieldLabel{Symptom}
SGD with momentum may fail to retrieve or correctly reuse its velocity buffer, producing incorrect updates or runtime failures in momentum-based training.
\par
\vspace{0.35em}
\FieldLabel{Root Cause}
As in Adam, the implementation uses inconsistent state-key naming between buffer initialization and buffer access. A momentum buffer cannot persist if later lookups ask for it under a different name.
\par
\vspace{0.35em}
\FieldLabel{Corrective Action Taken}
Align the momentum-buffer key names across initialization and update logic.
\par
\vspace{0.35em}
\begin{Verbatim}[commandchars=\\\{\},fontsize=\footnotesize,breaklines=true,breaksymbol={},frame=leftline,framerule=2pt,rulecolor=\color{gray!40},xleftmargin=6pt,formatcom=\singlespacing]
\textcolor{gray}{# Optimizers/sgd_momentum.py L49}
\textcolor{BrickRed}{- state["vel"] = torch.zeros_like(p)}
\textcolor{BrickRed}{- v = state["velocity"]}
\textcolor{OliveGreen}{+ state["velocity"] = torch.zeros_like(p)}
\textcolor{OliveGreen}{+ v = state["velocity"]}
\end{Verbatim}
\vspace{0.22em}
\FieldLabel{Impact on Training}
\textbf{Pre-fix:} Momentum is either unavailable due to a direct KeyError or silently reset, so the optimizer loses the very mechanism that smooths noisy gradients and accelerates motion along consistent descent directions. That raises gradient noise, slows convergence, and makes F1/EM curves far less stable than they should be. Broken optimizer state either prevents valid updates or discards the temporal information the rule is meant to reuse. That prevents stable, interpretable learning behaviour because gradients are not translated into the intended parameter trajectory across steps or resumes.\par\smallskip
\textbf{Post-fix:} Velocity is preserved across steps, so SGD-momentum again accumulates useful directional history instead of acting like broken plain SGD. State reuse is now consistent across iterations.
\par
\vspace{0.35em}
\end{bugbox}

\begin{bugbox}{BUG-052 \checkmark{} Adam weight-decay sign fixed}
{\color{MetaGray}\small Stage: 2 \;\textbar\; Category: Deep Learning Mechanism / Optimization Strategies \;\textbar\; Source: Optimizers/adam.py L53\par}
\vspace{0.22em}
\FieldLabel{Nature}
Inconsistent mechanism-level implementation flaw in Adam weight decay.
\par
\vspace{0.35em}
\FieldLabel{Symptom}
Weight decay is applied with the incorrect sign in Adam, so parameters are pushed away from regularized shrinkage rather than toward it.
\par
\vspace{0.35em}
\FieldLabel{Root Cause}
Weight decay is a penalty that must add +wd * p to the gradient path (equivalently shrink parameters in the update). Flipping the sign turns a regularizer into an anti-regularizer that actively grows parameter norms.
\par
\vspace{0.35em}
\FieldLabel{Corrective Action Taken}
Reverse the sign of the decay term so it performs shrinkage rather than expansion.
\par
\vspace{0.35em}
\begin{Verbatim}[commandchars=\\\{\},fontsize=\footnotesize,breaklines=true,breaksymbol={},frame=leftline,framerule=2pt,rulecolor=\color{gray!40},xleftmargin=6pt,formatcom=\singlespacing]
\textcolor{gray}{# Optimizers/adam.py L53}
\textcolor{BrickRed}{- grad = grad.add(p, alpha=-wd)}
\textcolor{OliveGreen}{+ grad = grad.add(p, alpha=wd)}
\end{Verbatim}
\vspace{0.22em}
\FieldLabel{Impact on Training}
\textbf{Pre-fix:} Instead of discouraging large weights, the optimizer rewards them. Qualitatively this makes attention projections and feed-forward layers increasingly overconfident and brittle; quantitatively it tends to inflate parameter norms, worsen overfitting, and can produce oscillatory or NaN-prone loss behaviour under long runs. The optimiser no longer follows its intended update rule, so stability and convergence claims become unreliable. \par\smallskip
\textbf{Post-fix:} Weight decay again constrains parameter growth, so Adam's update rule regularizes rather than amplifies the model. Optimizer updates again match the intended rule. 
\par
\vspace{0.35em}
\end{bugbox}

\begin{bugbox}{BUG-053 \checkmark{} Adam second moment uses squared gradients}
{\color{MetaGray}\small Stage: 2 \;\textbar\; Category: Deep Learning Mechanism / Optimization Strategies \;\textbar\; Source: Optimizers/adam.py L69\par}
\vspace{0.22em}
\FieldLabel{Nature}
Inconsistent mechanism-level implementation flaw in Adam second-moment accumulation.
\par
\vspace{0.35em}
\FieldLabel{Symptom}
Adam's variance estimate is incorrect if gradients are not squared before accumulation.
\par
\vspace{0.35em}
\FieldLabel{Root Cause}
The second moment in Adam is an exponential average of g\_t\^{}2, not of g\_t. Without the square, the denominator in the Adam step no longer tracks gradient magnitude and therefore cannot normalize updates correctly.
\par
\vspace{0.35em}
\FieldLabel{Corrective Action Taken}
Accumulate squared gradients in the second-moment buffer.
\par
\vspace{0.35em}
\begin{Verbatim}[commandchars=\\\{\},fontsize=\footnotesize,breaklines=true,breaksymbol={},frame=leftline,framerule=2pt,rulecolor=\color{gray!40},xleftmargin=6pt,formatcom=\singlespacing]
\textcolor{gray}{# Optimizers/adam.py L69}
\textcolor{BrickRed}{- v.mul_(beta2).add_(grad, alpha=1.0 - beta2)}
\textcolor{OliveGreen}{+ v.mul_(beta2).addcmul_(grad, grad, value=1.0 - beta2)}
\end{Verbatim}
\vspace{0.22em}
\FieldLabel{Impact on Training}
\textbf{Pre-fix:} Adam loses its central mechanism for per-parameter step normalization. When gradients are large, updates are not damped enough; when gradients are small or sign-changing, updates can be mis-scaled in the opposite direction. The result is poorer loss conditioning, more erratic late-stage training, and lower ceiling performance on F1/EM. The optimiser no longer follows its intended update rule, so stability and convergence claims become unreliable. \par\smallskip
\textbf{Post-fix:} The denominator now reflects gradient magnitude as intended, so adaptive step sizes become meaningful again and optimization is far more stable. Optimizer updates again match the intended rule. 
\par
\vspace{0.35em}
\end{bugbox}

\begin{bugbox}{BUG-039/040 \checkmark{} Adam bias-correction exponentiation fixed}
{\color{MetaGray}\small Stage: 2 \;\textbar\; Category: Deep Learning Mechanism / Optimization Strategies \;\textbar\; Source: Optimizers/adam.py L72-73\par}
\vspace{0.22em}
\FieldLabel{Nature}
Inconsistent mechanism-level implementation flaw in Adam bias correction.
\par
\vspace{0.35em}
\FieldLabel{Symptom}
Bias correction for Adam moments is computed incorrectly, which skews early-step update magnitudes and may destabilize optimization during the most sensitive phase of training.
\par
\vspace{0.35em}
\FieldLabel{Root Cause}
Bias correction depends on powers beta1\^{}t and beta2\^{}t. Using multiplication instead of exponentiation means the early moving averages are not debiased correctly, so the effective step size is systematically incorrect.
\par
\vspace{0.35em}
\FieldLabel{Corrective Action Taken}
Use the standard Adam bias-correction formulas with the correct coefficient powers at step t.
\par
\vspace{0.35em}
\begin{Verbatim}[commandchars=\\\{\},fontsize=\footnotesize,breaklines=true,breaksymbol={},frame=leftline,framerule=2pt,rulecolor=\color{gray!40},xleftmargin=6pt,formatcom=\singlespacing]
\textcolor{gray}{# Optimizers/adam.py L72-73}
\textcolor{BrickRed}{- bias_correction1 = 1.0 - beta1 * t}
\textcolor{BrickRed}{- bias_correction2 = 1.0 - beta2 * t}
\textcolor{OliveGreen}{+ bias_correction1 = 1.0 - beta1 ** t}
\textcolor{OliveGreen}{+ bias_correction2 = 1.0 - beta2 ** t}
\end{Verbatim}
\vspace{0.22em}
\FieldLabel{Impact on Training}
\textbf{Pre-fix:} The most fragile part of Adam training is the beginning, when moment estimates are still heavily biased. Wrong correction distorts the effective learning rate exactly there. Qualitatively this can make the model either lurch too far before representations settle or move too slowly to escape poor initial regimes; quantitatively it shows up as unstable warmup loss, delayed F1/EM growth, or outright divergence. The optimiser no longer follows its intended update rule, so stability and convergence claims become unreliable. \par\smallskip
\textbf{Post-fix:} Early-step scaling now matches the actual Adam algorithm, so warmup and scheduler behaviour are interpretable rather than accidentally counteracting optimizer math. Optimizer updates again match the intended rule. 
\par
\vspace{0.35em}
\end{bugbox}

\begin{bugbox}{BUG-054 \checkmark{} SGD weight-decay sign fixed}
{\color{MetaGray}\small Stage: 2 \;\textbar\; Category: Deep Learning Mechanism / Optimization Strategies \;\textbar\; Source: Optimizers/sgd.py L38\par}
\vspace{0.22em}
\FieldLabel{Nature}
Inconsistent mechanism-level implementation flaw in SGD weight decay.
\par
\vspace{0.35em}
\FieldLabel{Symptom}
SGD weight decay is applied with the incorrect sign, so regularization becomes anti-regularization.
\par
\vspace{0.35em}
\FieldLabel{Root Cause}
As in Adam, the decay term must reduce parameter magnitude. A sign error changes the geometry of the update from shrinkage to growth.
\par
\vspace{0.35em}
\FieldLabel{Corrective Action Taken}
Correct the sign of the decay term in the SGD update.
\par
\vspace{0.35em}
\begin{Verbatim}[commandchars=\\\{\},fontsize=\footnotesize,breaklines=true,breaksymbol={},frame=leftline,framerule=2pt,rulecolor=\color{gray!40},xleftmargin=6pt,formatcom=\singlespacing]
\textcolor{gray}{# Optimizers/sgd.py L38}
\textcolor{BrickRed}{- grad = grad.add(p, alpha=-wd)}
\textcolor{OliveGreen}{+ grad = grad.add(p, alpha=wd)}
\end{Verbatim}
\vspace{0.22em}
\FieldLabel{Impact on Training}
\textbf{Pre-fix:} Parameter norms can grow faster than intended, weakening any intended L2-style control. Even if training loss decreases for a while, the model is being driven toward larger, less stable weights, so dev-set metrics can plateau or regress earlier than expected. The optimiser no longer follows its intended update rule, so stability and convergence claims become unreliable. \par\smallskip
\textbf{Post-fix:} The optimizer again penalizes large weights rather than rewarding them, restoring proper regularization geometry. Optimizer updates again match the intended rule. 
\par
\vspace{0.35em}
\end{bugbox}

\begin{bugbox}{BUG-055 \checkmark{} SGD momentum update direction fixed}
{\color{MetaGray}\small Stage: 2 \;\textbar\; Category: Deep Learning Mechanism / Optimization Strategies \;\textbar\; Source: Optimizers/sgd\_momentum.py L54\par}
\vspace{0.22em}
\FieldLabel{Nature}
Inconsistent mechanism-level implementation flaw in SGD momentum update direction.
\par
\vspace{0.35em}
\FieldLabel{Symptom}
The momentum update may be applied in the incorrect direction, producing systematically incorrect parameter motion and distorting training dynamics.
\par
\vspace{0.35em}
\FieldLabel{Root Cause}
Momentum must accumulate a descent-oriented velocity vector. If the sign in the recurrence or the parameter update is reversed, momentum reinforces the incorrect direction in parameter space instead of smoothing descent.
\par
\vspace{0.35em}
\FieldLabel{Corrective Action Taken}
Restore the standard momentum-update direction so velocity reinforces descent rather than opposing it.
\par
\vspace{0.35em}
\begin{Verbatim}[commandchars=\\\{\},fontsize=\footnotesize,breaklines=true,breaksymbol={},frame=leftline,framerule=2pt,rulecolor=\color{gray!40},xleftmargin=6pt,formatcom=\singlespacing]
\textcolor{gray}{# Optimizers/sgd_momentum.py L54}
\textcolor{BrickRed}{- v.mul_(mu).sub_(grad)}
\textcolor{OliveGreen}{+ v.mul_(mu).add_(grad)}
\end{Verbatim}
\vspace{0.22em}
\FieldLabel{Impact on Training}
\textbf{Pre-fix:} Momentum no longer acts as a smoother on coherent descent directions; it can instead amplify the wrong motion and push the optimizer away from minima. Qualitatively this produces oscillation or divergence; quantitatively it slows convergence and can keep loss and span metrics from improving even when the rest of the pipeline is correct. The optimiser no longer follows its intended update rule, so stability and convergence claims become unreliable. \par\smallskip
\textbf{Post-fix:} Velocity now reinforces descent-consistent updates, so SGD-momentum regains its intended acceleration and stabilization role. Optimizer updates again match the intended rule. 
\par
\vspace{0.35em}
\end{bugbox}

\ReportSubsection{7.3 Learning Rate Schedulers}

{\itshape
Cosine, lambda, and step-based schedules.
\par}

\begin{bugbox}{BUG-024 \checkmark{} Cosine scheduler math.pi fix}
{\color{MetaGray}\small Stage: 2 \;\textbar\; Category: Deep Learning Mechanism / Learning Rate Scheduling \;\textbar\; Source: Schedulers/cosine\_scheduler.py L28\par}
\vspace{0.22em}
\FieldLabel{Nature}
Incorrect structural implementation inside cosine-scheduler execution.
\par
\vspace{0.35em}
\FieldLabel{Symptom}
The cosine scheduler uses an invalid constant, producing a direct AttributeError instead of a usable learning-rate value.
\par
\vspace{0.35em}
\FieldLabel{Root Cause}
Python's math module defines math.pi, not math.PI. A scheduler that cannot compute its scalar coefficient cannot advance training at all.
\par
\vspace{0.35em}
\FieldLabel{Corrective Action Taken}
Restore the standard cosine expression using math.pi.
\par
\vspace{0.35em}
\begin{Verbatim}[commandchars=\\\{\},fontsize=\footnotesize,breaklines=true,breaksymbol={},frame=leftline,framerule=2pt,rulecolor=\color{gray!40},xleftmargin=6pt,formatcom=\singlespacing]
\textcolor{gray}{# Schedulers/cosine_scheduler.py L28}
\textcolor{BrickRed}{- math.cos(math.PI * t / self.T_max)}
\textcolor{OliveGreen}{+ math.cos(math.pi * t / self.T_max)}
\end{Verbatim}
\vspace{0.22em}
\FieldLabel{Impact on Training}
\textbf{Pre-fix:} Any run selecting the cosine scheduler fails immediately at the first LR computation, so that entire training route is dead. Even though this bug lives in the scheduler rather than in the forward graph, it still kills the training route because the optimizer never receives a valid learning-rate value. The result is not just inconvenience: the entire schedule-dependent experiment becomes unexecutable.\par\smallskip
\textbf{Post-fix:} Cosine scheduling becomes executable again, restoring a valid alternative optimizer path for experimentation. The scheduler can now produce valid step sizes.
\par
\vspace{0.35em}
\end{bugbox}

\begin{bugbox}{BUG-056 \checkmark{} Missing 0.5 factor in cosine scheduler fixed}
{\color{MetaGray}\small Stage: 2 \;\textbar\; Category: Deep Learning Mechanism / Learning Rate Scheduling \;\textbar\; Source: Schedulers/cosine\_scheduler.py L28\par}
\vspace{0.22em}
\FieldLabel{Nature}
Inconsistent mechanism-level implementation flaw in cosine-scheduler amplitude.
\par
\vspace{0.35em}
\FieldLabel{Symptom}
The cosine schedule is scaled incorrectly if the standard 0.5 factor is omitted.
\par
\vspace{0.35em}
\FieldLabel{Root Cause}
The standard cosine schedule uses 1/2 (1 + cos(.)) to map cosine values from {[}-1, 1{]} into {[}0, 1{]}. Omitting the factor changes the amplitude of the schedule and therefore the actual learning-rate magnitude.
\par
\vspace{0.35em}
\FieldLabel{Corrective Action Taken}
Restore the missing 0.5 factor.
\par
\vspace{0.35em}
\begin{Verbatim}[commandchars=\\\{\},fontsize=\footnotesize,breaklines=true,breaksymbol={},frame=leftline,framerule=2pt,rulecolor=\color{gray!40},xleftmargin=6pt,formatcom=\singlespacing]
\textcolor{gray}{# Schedulers/cosine_scheduler.py L28}
\textcolor{BrickRed}{- lr = self.eta_min + (base_lr - self.eta_min) * (1 + math.cos(math.pi * t / self.T_max))}
\textcolor{OliveGreen}{+ lr = self.eta_min + 0.5 * (base_lr - self.eta_min) * (1 + math.cos(math.pi * t / self.T_max))}
\end{Verbatim}
\vspace{0.22em}
\FieldLabel{Impact on Training}
\textbf{Pre-fix:} At t = 0 the schedule starts near 2 * base\_lr instead of base\_lr. That means the optimizer enters training more aggressively than intended, so early loss can spike and the whole decay curve becomes miscalibrated. The learning-rate path no longer matches the intended schedule. \par\smallskip
\textbf{Post-fix:} The cosine trajectory now has the correct range and phase, so schedule-based comparisons are mathematically meaningful. The learning-rate path again follows the intended schedule. 
\par
\vspace{0.35em}
\end{bugbox}

\begin{bugbox}{BUG-057 \checkmark{} Lambda scheduler multiplicative fix}
{\color{MetaGray}\small Stage: 2 \;\textbar\; Category: Deep Learning Mechanism / Learning Rate Scheduling \;\textbar\; Source: Schedulers/lambda\_scheduler.py L23\par}
\vspace{0.22em}
\FieldLabel{Nature}
Inconsistent mechanism-level implementation flaw in lambda-scheduler multiplicative update.
\par
\vspace{0.35em}
\FieldLabel{Symptom}
The lambda scheduler applies its factor incorrectly, so the resulting learning rate does not reflect the intended multiplicative schedule. In the broken path, a nominal factor of 1.0 increases lr by 1 instead of leaving it unchanged.
\par
\vspace{0.35em}
\FieldLabel{Root Cause}
A lambda scheduler is meant to return a multiplier on the base learning rate. If that factor is combined additively rather than multiplicatively, the schedule no longer implements the intended policy.
\par
\vspace{0.35em}
\FieldLabel{Corrective Action Taken}
Apply the lambda function as a multiplicative scale on the base learning rate.
\par
\vspace{0.35em}
\begin{Verbatim}[commandchars=\\\{\},fontsize=\footnotesize,breaklines=true,breaksymbol={},frame=leftline,framerule=2pt,rulecolor=\color{gray!40},xleftmargin=6pt,formatcom=\singlespacing]
\textcolor{gray}{# Schedulers/lambda_scheduler.py L23}
\textcolor{BrickRed}{- return [base_lr + factor for base_lr in self.base_lrs]}
\textcolor{OliveGreen}{+ return [base_lr * factor for base_lr in self.base_lrs]}
\end{Verbatim}
\vspace{0.22em}
\FieldLabel{Impact on Training}
\textbf{Pre-fix:} Learning-rate values can jump from 0.001 to roughly 1.001 under what should have been a no-op factor, which is catastrophic for Adam and readily explains exploding loss or NaN routes. This is not a minor schedule drift; it changes the optimizer regime entirely. The learning-rate path no longer matches the intended schedule. \par\smallskip
\textbf{Post-fix:} The scheduler again scales the base learning rate instead of overriding it with an effectively huge additive offset. The learning-rate path again follows the intended schedule. 
\par
\vspace{0.35em}
\end{bugbox}

\begin{bugbox}{BUG-058 \checkmark{} Step scheduler exponentiation fix}
{\color{MetaGray}\small Stage: 2 \;\textbar\; Category: Deep Learning Mechanism / Learning Rate Scheduling \;\textbar\; Source: Schedulers/step\_scheduler.py L25\par}
\vspace{0.22em}
\FieldLabel{Nature}
Inconsistent mechanism-level implementation flaw in step-scheduler decay exponent.
\par
\vspace{0.35em}
\FieldLabel{Symptom}
The step scheduler decays the learning rate by the incorrect amount when exponentiation is handled incorrectly.
\par
\vspace{0.35em}
\FieldLabel{Root Cause}
Step decay must apply a factor such as gamma\^{}k, where k is the number of completed decay intervals. If k is computed correctly but the decay is applied linearly rather than exponentially, the schedule no longer follows the intended staircase geometry.
\par
\vspace{0.35em}
\FieldLabel{Corrective Action Taken}
Use the correct exponent corresponding to the number of completed decay intervals.
\par
\vspace{0.35em}
\begin{Verbatim}[commandchars=\\\{\},fontsize=\footnotesize,breaklines=true,breaksymbol={},frame=leftline,framerule=2pt,rulecolor=\color{gray!40},xleftmargin=6pt,formatcom=\singlespacing]
\textcolor{gray}{# Schedulers/step_scheduler.py L25}
\textcolor{BrickRed}{- return [base_lr * self.gamma * (t // self.step_size) for base_lr in self.base_lrs]}
\textcolor{OliveGreen}{+ return [base_lr * (self.gamma ** (t // self.step_size)) for base_lr in self.base_lrs]}
\end{Verbatim}
\vspace{0.22em}
\FieldLabel{Impact on Training}
\textbf{Pre-fix:} The learning rate can collapse too fast, decay too slowly, or even become zero at the wrong time, all of which make optimizer comparisons misleading. Quantitatively this appears as stalled loss reduction or needlessly late convergence. The learning-rate path no longer matches the intended schedule. \par\smallskip
\textbf{Post-fix:} The schedule again follows the intended stepwise exponential decay, so its optimisation behaviour is interpretable and reproducible. The learning-rate path again follows the intended schedule. 
\par
\vspace{0.35em}
\end{bugbox}

\begin{bugbox}{BUG-N001 \checkmark{} Lambda warmup scheduler redesigned into a picklable object}
{\color{MetaGray}\small Stage: 2 \;\textbar\; Category: Deep Learning Mechanism / Learning Rate Scheduling \;\textbar\; Source: Schedulers/scheduler.py L29\par}
\vspace{0.22em}
\FieldLabel{Nature}
Incorrect engineering interface and inconsistent mechanism-level implementation flaw in lambda warmup scheduler design and serialization.
\par
\vspace{0.35em}
\FieldLabel{Symptom}
The warmup scheduler is encoded as an anonymous lambda, which is neither safely serializable nor expressive enough to represent the intended warmup law.
\par
\vspace{0.35em}
\FieldLabel{Root Cause}
Checkpoint-safe schedulers need explicit state, and the warmup rule itself should define a time-dependent multiplier. A constant lambda closure satisfies neither requirement: it is brittle under serialization and it cannot implement the intended warmup behaviour.
\par
\vspace{0.35em}
\FieldLabel{Corrective Action Taken}
Redesign the warmup scheduler as a picklable callable object with explicit state and log-warmup update logic.
\par
\vspace{0.35em}
\begin{Verbatim}[commandchars=\\\{\},fontsize=\footnotesize,breaklines=true,breaksymbol={},frame=leftline,framerule=2pt,rulecolor=\color{gray!40},xleftmargin=6pt,formatcom=\singlespacing]
\textcolor{gray}{# Schedulers/scheduler.py L29}
\textcolor{OliveGreen}{+ class _WarmupFactor:}
\textcolor{OliveGreen}{+  def __init__(self, warmup_steps: int):}
\textcolor{OliveGreen}{+   self.warmup_steps = max(1, warmup_steps)}
\textcolor{OliveGreen}{+   self.c = 1.0 / math.log(self.warmup_steps)}
\textcolor{OliveGreen}{+  def __call__(self, step: int) -> float:}
\textcolor{OliveGreen}{+   return self.c * math.log(step + 1) if step < self.warmup_steps else 1.0}
\textcolor{BrickRed}{- return LambdaScheduler(optimizer, lr_lambda=lambda _: 1.0)}
\textcolor{OliveGreen}{+ return LambdaScheduler(optimizer, lr_lambda=_WarmupFactor(warmup_steps))}
\end{Verbatim}
\vspace{0.22em}
\FieldLabel{Impact on Training}
\textbf{Pre-fix:} Warmup behaviour cannot be resumed reliably, and the scheduler logic does not encode the intended time-varying scaling law. As a result, early Adam steps may be mis-scaled both numerically and procedurally, which weakens stability and undermines schedule-dependent comparisons.\par\smallskip
\textbf{Post-fix:} The warmup rule is now explicit, stateful, and serializable. Resume behaviour and early-step scaling are therefore both reproducible and mathematically aligned with the intended schedule.
\par
\vspace{0.35em}
\end{bugbox}

\ReportSubsection{7.4 Training Pipeline}

{\itshape
Argument construction, backpropagation, gradient clipping, best-checkpoint retention, and EMA integration.
\par}

\begin{bugbox}{BUG-025 \checkmark{} argparse.Namespace construction fixed}
{\color{MetaGray}\small Stage: 1 \;\textbar\; Category: Runtime Failure / Training Loop Logic \;\textbar\; Source: TrainTools/train.py L112\par}
\vspace{0.22em}
\FieldLabel{Nature}
Incorrect structural pipeline implementation in args object construction.
\par
\vspace{0.35em}
\FieldLabel{Symptom}
Training or evaluation argument objects are constructed incorrectly, which may induce immediate runtime failures or, more dangerously, silent misconfiguration in downstream components.
\par
\vspace{0.35em}
\FieldLabel{Root Cause}
The program relies on a namespace object whose fields are later read by the model, data, scheduler, and checkpoint code. Passing a dictionary positionally creates an invalid Namespace contract.
\par
\vspace{0.35em}
\FieldLabel{Corrective Action Taken}
Rebuild the Namespace object using keyword arguments instead of a single positional dictionary.
\par
\vspace{0.35em}
\begin{Verbatim}[commandchars=\\\{\},fontsize=\footnotesize,breaklines=true,breaksymbol={},frame=leftline,framerule=2pt,rulecolor=\color{gray!40},xleftmargin=6pt,formatcom=\singlespacing]
\textcolor{gray}{# TrainTools/train.py L112}
\textcolor{BrickRed}{- args = argparse.Namespace(\{k: v for k, v in locals().items()\})}
\textcolor{OliveGreen}{+ args = argparse.Namespace(**\{k: v for k, v in locals().items()\})}
\end{Verbatim}
\vspace{0.22em}
\FieldLabel{Impact on Training and Evaluation}
\textbf{Pre-fix:} In the explicit failure mode the run aborts at startup. In the silent failure mode, later components can read wrong defaults, making any reported metrics or checkpoints configuration-inconsistent. The pipeline fails before the intended train/eval route is even instantiated, so there is no meaningful learning curve or metric trace to analyse. Empirically, this belongs to the class of setup blockers that had to be removed before any substantive debugging of model behaviour was possible.\par\smallskip
\textbf{Post-fix:} Runtime arguments again form a valid contract across the pipeline, so later behaviour reflects the actual notebook configuration rather than accidental fallbacks. The training loop can now be constructed and executed as intended, which turns later loss curves and evaluation traces into meaningful evidence rather than unreachable code paths.
\par
\vspace{0.35em}
\end{bugbox}

\begin{bugbox}{BUG-026 \checkmark{} loss.item().backward() fixed to loss.backward()}
{\color{MetaGray}\small Stage: 1 \;\textbar\; Category: Runtime Failure / Gradient Computation \;\textbar\; Source: TrainTools/train\_utils.py L34\par}
\vspace{0.22em}
\FieldLabel{Nature}
Incorrect structural pipeline implementation in loss backward invocation.
\par
\vspace{0.35em}
\FieldLabel{Symptom}
Calling backward() on loss.item() triggers a runtime failure because loss.item() is a Python scalar rather than a differentiable tensor.
\par
\vspace{0.35em}
\FieldLabel{Root Cause}
loss.item() detaches the scalar numerical value from the computation graph. Once converted to a Python float, it carries no gradient history, so backpropagation is impossible.
\par
\vspace{0.35em}
\FieldLabel{Corrective Action Taken}
Call backward() on the tensor-valued loss directly.
\par
\vspace{0.35em}
\begin{Verbatim}[commandchars=\\\{\},fontsize=\footnotesize,breaklines=true,breaksymbol={},frame=leftline,framerule=2pt,rulecolor=\color{gray!40},xleftmargin=6pt,formatcom=\singlespacing]
\textcolor{gray}{# TrainTools/train_utils.py L34}
\textcolor{BrickRed}{- loss.item().backward()}
\textcolor{OliveGreen}{+ loss.backward()}
\end{Verbatim}
\vspace{0.22em}
\FieldLabel{Impact on Training}
\textbf{Pre-fix:} No gradient can flow through the model at all, so the network never learns and the first training step terminates with a runtime error. Without a valid loss/backward contract, gradients cannot be computed in the mathematically intended way, so the optimizer never receives a usable training signal and the model cannot demonstrate decreasing loss for the right reason. Empirically, this is a direct training blocker rather than a subtle quality issue.\par\smallskip
\textbf{Post-fix:} Gradients are now computed from the actual computation graph, reactivating the entire optimization path. Autograd now receives the correct objects and can propagate the intended supervision signal through the network, so subsequent optimization behaviour becomes meaningful rather than vacuous.
\par
\vspace{0.35em}
\end{bugbox}

\begin{bugbox}{BUG-059 \checkmark{} Gradient clipping moved before optimizer step}
{\color{MetaGray}\small Stage: 2 \;\textbar\; Category: Deep Learning Mechanism / Optimization Strategies \;\textbar\; Source: TrainTools/train\_utils.py L35\par}
\vspace{0.22em}
\FieldLabel{Nature}
Inconsistent mechanism-level implementation flaw in gradient-clipping order.
\par
\vspace{0.35em}
\FieldLabel{Symptom}
Gradient clipping is applied after the optimizer step instead of before it, so the actually executed update is not clipped.
\par
\vspace{0.35em}
\FieldLabel{Root Cause}
Gradient clipping only affects parameter updates if it is applied while gradients are still stored and before those gradients are consumed by the optimizer. Once the step has already been taken, clipping can no longer change that update.
\par
\vspace{0.35em}
\FieldLabel{Corrective Action Taken}
Move gradient clipping to the standard position between backpropagation and the optimizer step.
\par
\vspace{0.35em}
\begin{Verbatim}[commandchars=\\\{\},fontsize=\footnotesize,breaklines=true,breaksymbol={},frame=leftline,framerule=2pt,rulecolor=\color{gray!40},xleftmargin=6pt,formatcom=\singlespacing]
\textcolor{gray}{# TrainTools/train_utils.py L35}
\textcolor{BrickRed}{- optimizer.step()}
\textcolor{BrickRed}{- torch.nn.utils.clip_grad_norm_(model.parameters(), grad_clip)}
\textcolor{OliveGreen}{+ torch.nn.utils.clip_grad_norm_(model.parameters(), grad_clip)}
\textcolor{OliveGreen}{+ optimizer.step()}
\end{Verbatim}
\vspace{0.22em}
\FieldLabel{Impact on Training}
\textbf{Pre-fix:} Rare but large gradient spikes still produce full-size destructive updates, so even a pipeline that looks stable on most batches can be knocked off course by a single pathological step. This hurts loss smoothness and makes later checkpoints less reliable. The optimiser no longer follows its intended update rule, so stability and convergence claims become unreliable. \par\smallskip
\textbf{Post-fix:} The executed update is now the clipped update, so gradient clipping finally performs its intended stabilisation role. Optimizer updates again match the intended rule. 
\par
\vspace{0.35em}
\end{bugbox}

\begin{bugbox}{BUG-N015 \checkmark{} Best checkpoint retained explicitly}
{\color{MetaGray}\small Stage: 1 \;\textbar\; Category: Checkpoint Management / Training-Loop Semantics \;\textbar\; Source: TrainTools/train.py L206-223 \& TrainTools/train\_utils.py L44-60\par}
\vspace{0.22em}
\FieldLabel{Nature}
Incorrect checkpoint-management interface and inconsistent training-loop semantics in best-checkpoint retention.
\par
\vspace{0.35em}
\FieldLabel{Symptom}
The training loop may overwrite or fail to preserve the best-performing checkpoint, so later evaluation may use a suboptimal model even when a better one was observed during training.
\par
\vspace{0.35em}
\FieldLabel{Root Cause}
The checkpoint-saving logic did not properly distinguish between ``latest state'' persistence and ``best validation model'' retention. These are different objects with different purposes and must not share a single overwrite policy.
\par
\vspace{0.35em}
\FieldLabel{Corrective Action Taken}
Save and retain the best-performing checkpoint separately based on validation metrics.
\par
\vspace{0.35em}
\begin{Verbatim}[commandchars=\\\{\},fontsize=\footnotesize,breaklines=true,breaksymbol={},frame=leftline,framerule=2pt,rulecolor=\color{gray!40},xleftmargin=6pt,formatcom=\singlespacing]
\textcolor{gray}{# TrainTools/train.py L206-223 &}
\textcolor{OliveGreen}{+ is_best = dev_f1 > best_f1 or dev_em > best_em}
\textcolor{BrickRed}{- save_checkpoint(model, optimizer, scheduler, step, save_dir, config)}
\textcolor{OliveGreen}{+ save_checkpoint(model, optimizer, scheduler, step, save_dir, config, is_best=is_best)}

\textcolor{gray}{# TrainTools/train_utils.py L44-60}
\textcolor{BrickRed}{- def save_checkpoint(model, optimizer, scheduler, step, save_dir, config):}
\textcolor{OliveGreen}{+ def save_checkpoint(model, optimizer, scheduler, step, save_dir, config, is_best=False):}
\textcolor{OliveGreen}{+  if not is_best:}
\textcolor{OliveGreen}{+   return}
\end{Verbatim}
\vspace{0.22em}
\FieldLabel{Impact on Training and Evaluation}
\textbf{Pre-fix:} Final evaluation can underreport the repaired system simply because a later worse checkpoint happened to overwrite the real best model. That contaminates the quantitative story of the report itself. The model may genuinely learn during training, but the artifact kept for later evaluation can still correspond to a worse state than the true best one. That disconnects reported final metrics from the actual best learning behaviour achieved during the run and weakens ablation conclusions.\par\smallskip
\textbf{Post-fix:} The reported model now corresponds to the best observed validation state rather than the last saved state, so final F1/EM are representative rather than accidental. The checkpoint selected for later evaluation now corresponds to the best observed validation state, so final reported metrics are better aligned with the strongest model behaviour reached during training.
\par
\vspace{0.35em}
\end{bugbox}

\begin{bugbox}{BUG-N017 \checkmark{} EMA support added and integrated}
{\color{MetaGray}\small Stage: 2 \;\textbar\; Category: Deep Learning Mechanism / Regularization Techniques \;\textbar\; Source: TrainTools/train.py \& TrainTools/train\_utils.py \& TrainTools/ema.py\par}
\vspace{0.22em}
\FieldLabel{Nature}
Inconsistent mechanism-level implementation flaw in EMA parameter averaging.
\par
\vspace{0.35em}
\FieldLabel{Symptom}
Training and evaluation operate only on instantaneous parameters; no exponential moving-average state is maintained or swapped in for evaluation.
\par
\vspace{0.35em}
\FieldLabel{Root Cause}
EMA requires shadow-parameter storage, post-update accumulation, and evaluation-time substitution. Without all three elements, the training recipe omits the smoothing mechanism that stabilizes evaluation against short-term parameter noise.
\par
\vspace{0.35em}
\FieldLabel{Corrective Action Taken}
Add EMA state management and integrate its update and usage into the train/eval pipeline.
\par
\vspace{0.35em}
\begin{Verbatim}[commandchars=\\\{\},fontsize=\footnotesize,breaklines=true,breaksymbol={},frame=leftline,framerule=2pt,rulecolor=\color{gray!40},xleftmargin=6pt,formatcom=\singlespacing]
\textcolor{gray}{# TrainTools/train.py L206-223}
\textcolor{OliveGreen}{+ from TrainTools.ema import EMA}
 
\textcolor{BrickRed}{-  early_stop=50,}
\textcolor{OliveGreen}{+  early_stop=50,}
\textcolor{OliveGreen}{+  ema_decay=0.9999,}
 
\textcolor{OliveGreen}{+ ema = EMA(ema_decay)}
\textcolor{OliveGreen}{+ ema.register(model)}
 optimizer.step()
\textcolor{OliveGreen}{+ ema.update(model)}
\textcolor{OliveGreen}{+ ema.assign(model)}
 metrics, ans = run_eval(model, dev_dataset, dev_eval_file, dev_batches, batch_size, loss_fn)
\textcolor{OliveGreen}{+ ema.resume(model)}
 
\textcolor{gray}{# TrainTools/train_utils.py L44-60}
\textcolor{BrickRed}{- def train_single_epoch(model, dataset, optimizer, scheduler, loss_fn, batch_size):}
\textcolor{OliveGreen}{+ def train_single_epoch(model, dataset, optimizer, scheduler, loss_fn, batch_size, ema=None):}
 optimizer.step()
\textcolor{OliveGreen}{+ if ema is not None:}
\textcolor{OliveGreen}{+  ema.update(model)}
 
\textcolor{gray}{# TrainTools/ema.py}
\textcolor{OliveGreen}{+ class EMA:}
\textcolor{OliveGreen}{+  def __init__(self, decay: float):}
\textcolor{OliveGreen}{+   self.decay = decay}
\textcolor{OliveGreen}{+   self.shadow = \{\}}
\textcolor{OliveGreen}{+   self.backup = \{\}}
\textcolor{OliveGreen}{+}
\textcolor{OliveGreen}{+  def register(self, model):}
\textcolor{OliveGreen}{+   for name, param in model.named_parameters():}
\textcolor{OliveGreen}{+    if param.requires_grad:}
\textcolor{OliveGreen}{+     self.shadow[name] = param.data.clone()}
\textcolor{OliveGreen}{+}
\textcolor{OliveGreen}{+  def update(self, model):}
\textcolor{OliveGreen}{+   for name, param in model.named_parameters():}
\textcolor{OliveGreen}{+    if param.requires_grad:}
\textcolor{OliveGreen}{+     self.shadow[name] = (1.0 - self.decay) * param.data + self.decay * self.shadow[name]}
\textcolor{OliveGreen}{+}
\textcolor{OliveGreen}{+  def assign(self, model):}
\textcolor{OliveGreen}{+   for name, param in model.named_parameters():}
\textcolor{OliveGreen}{+    if param.requires_grad:}
\textcolor{OliveGreen}{+     self.backup[name] = param.data.clone()}
\textcolor{OliveGreen}{+     param.data.copy_(self.shadow[name])}
\textcolor{OliveGreen}{+}
\textcolor{OliveGreen}{+  def resume(self, model):}
\textcolor{OliveGreen}{+   for name, param in model.named_parameters():}
\textcolor{OliveGreen}{+    if param.requires_grad:}
\textcolor{OliveGreen}{+     param.data.copy_(self.backup[name])}
\end{Verbatim}
\vspace{0.22em}
\FieldLabel{Impact on Training and Evaluation}
\textbf{Pre-fix:} Validation is performed on noisier instantaneous weights rather than on the smoothed parameter trajectory. That increases checkpoint-to-checkpoint variance and can understate the model's true generalization behaviour even when the training objective is otherwise correct.\par\smallskip
\textbf{Post-fix:} Evaluation now uses the moving-average parameters, so validation and final test scores reflect a more stable estimate of the learned model.
\par
\vspace{0.35em}
\end{bugbox}




\ReportSection{8 Notebook Implementations, Results, and Bridge to Later Mechanism Studies}

\ReportSubsection{8.1 Actual notebook implementations}

This notebook study does not attempt a paper-scale reproduction. The dataset regime is materially smaller, so the final settings were selected to expose stable learning behaviour, interpretable monitoring, and reproducible checkpoint selection under the repaired implementation rather than to mirror the original training budget exactly. The original notebook defaults explicitly carried into this comparison are \texttt{num\_steps=60000}, \texttt{batch\_size=8}, and \texttt{seed=42}; the final runs keep the 60000-step budget and the same seed while changing the optimizer, scheduler, batch size, early-stopping patience, and regularisation bundle.

Throughout this section, metric names follow the notebook output exactly: \texttt{STEP loss} denotes the current optimisation-step objective, \texttt{VALID(train)} denotes the training-side monitor printed at each checkpoint, and \texttt{TEST} denotes the checkpoint evaluation monitor and the standalone evaluation output.

Adam + lambda is the primary route because the repaired notebook is already learnable under the two SGD references, so the remaining question is optimisation quality rather than basic executability. It is also the optimizer/scheduler family used in the QANet paper and, more generally, the more stable adaptive choice for this architecture. The two SGD routes are therefore retained only as learnability baselines and comparison points for the later Adam analysis.

\scriptsize
\begin{center}
\setlength{\tabcolsep}{2pt}
\begin{tabularx}{\linewidth}{@{}>{\raggedright\arraybackslash}p{0.18\linewidth}>{\raggedright\arraybackslash}p{0.22\linewidth}>{\raggedright\arraybackslash}p{0.12\linewidth}>{\raggedright\arraybackslash}p{0.15\linewidth}>{\raggedright\arraybackslash}X@{}}
\toprule
\textbf{Run} & \textbf{Actual notebook settings} & \textbf{Budget} & \textbf{Notebook monitoring} & \textbf{Role in this section} \\
\midrule
Adam + lambda & optimizer=adam;\newline scheduler=lambda;\newline loss=qa\_nll;\newline dropout=0.3;\newline dropout\_char=0.15;\newline ema\_decay=0.9999;\newline batch\_size=32 & num\_steps=\newline 60000;\newline seed=42 & checkpoint=200;\newline early\_stop=50;\newline labels=\texttt{VALID(train)}\newline / \texttt{TEST} & Primary notebook route; strongest \texttt{TEST} performance; closest to the paper-style optimisation recipe. \\
SGD + step & optimizer=sgd;\newline scheduler=step;\newline learning\_rate=0.02;\newline lr\_step\_size=10000;\newline lr\_gamma=0.8;\newline ema\_decay=0;\newline batch\_size=16 & num\_steps=\newline 60000;\newline seed=42 & checkpoint=200;\newline early\_stop=100;\newline labels=\texttt{VALID(train)}\newline / \texttt{TEST} & Reference run showing that the repaired notebook is learnable even under plain SGD. \\
SGD-momentum + cosine & optimizer=sgd\_momentum;\newline momentum=0.9;\newline scheduler=cosine;\newline learning\_rate=0.01;\newline ema\_decay=0;\newline batch\_size=16 & num\_steps=\newline 60000;\newline seed=42 & checkpoint=200;\newline early\_stop=100;\newline labels=\texttt{VALID(train)}\newline / \texttt{TEST} & Stronger reference run showing how a non-adaptive optimizer improves once momentum and smoother decay are added. \\
\bottomrule
\end{tabularx}
\end{center}
\setlength{\tabcolsep}{4.5pt}
\normalsize


\ReportSubsection{8.2 SGD reference runs as learnability baselines}

The two SGD routes are retained only as reference baselines. Their purpose is to show that, after the implementation repairs, the notebook already learns under comparatively simple non-adaptive optimisation and that changing the optimizer/scheduler bundle still materially changes the final \texttt{TEST} behaviour.

\small
\setlength{\LTleft}{0pt}
\setlength{\LTright}{0pt}
\begin{longtable}{>{\raggedright\arraybackslash}p{0.16\linewidth}>{\raggedright\arraybackslash}p{0.18\linewidth}>{\raggedright\arraybackslash}p{0.20\linewidth}>{\raggedright\arraybackslash}p{0.18\linewidth}>{\raggedright\arraybackslash}p{0.22\linewidth}}
\toprule
\textbf{Run} & \textbf{Initial logged \texttt{TEST} checkpoint} & \textbf{Best logged \texttt{TEST} checkpoint} & \textbf{Standalone \texttt{TEST} evaluation of \texttt{model\_best.pt}} & \textbf{Analytical significance} \\
\midrule
\endfirsthead
\toprule
\textbf{Run} & \textbf{Initial logged \texttt{TEST} checkpoint} & \textbf{Best logged \texttt{TEST} checkpoint} & \textbf{Standalone \texttt{TEST} evaluation of \texttt{model\_best.pt}} & \textbf{Analytical significance} \\
\midrule
\endhead
\bottomrule
\endfoot
SGD + step & Step 200: loss 11.224925; F1 7.486044; EM 0.208333 & Step 60000: loss 6.493038; F1 26.370738; EM 18.750000 & loss 6.791959; F1 23.862558; EM 15.747731 & Even the simplest route moves from near-random extraction to non-trivial span selection. The repaired notebook therefore no longer depends on adaptive optimisation merely to accumulate usable answer signal. \\
SGD-momentum + cosine & Step 200: loss 11.224555; F1 6.445970; EM 1.125000 & Step 54400: loss 6.620702; F1 41.075035; EM 33.250000 & loss 6.586266; F1 39.028693; EM 30.186335 & Adding momentum and smooth annealing materially raises final \texttt{TEST} F1/EM over plain SGD. Once the code path is repaired, optimisation design still changes how much of the recovered signal is converted into usable span quality. \\
\end{longtable}
\normalsize

These baselines establish two facts. First, learnability is already present without the full paper-style training bundle. Second, optimizer and scheduler choice still strongly influence the final checkpoint. The final configuration therefore moves to Adam + lambda, which is both the paper-aligned optimizer family and the more stable adaptive continuation of the baseline study.

\ReportSubsection{8.3 Adam + lambda: parameterisation relative to the paper-style recipe and the SGD baselines}

The remaining question is therefore not whether the notebook can learn, but why the final Adam route differs from the simpler SGD baselines.

\scriptsize
\begin{center}
\begin{tabularx}{\linewidth}{>{\raggedright\arraybackslash}p{0.11\linewidth}>{\raggedright\arraybackslash}p{0.15\linewidth}>{\raggedright\arraybackslash}p{0.15\linewidth}>{\raggedright\arraybackslash}p{0.22\linewidth}>{\raggedright\arraybackslash}X}
\toprule
\textbf{Parameter family} & \textbf{Adam + lambda} & \textbf{SGD / SGD-momentum references} & \textbf{Real notebook motivation} & \textbf{Paper / theoretical justification} \\
\midrule
Objective & \texttt{qa\_nll} & \texttt{qa\_nll} in both reference runs & Keep the training target fixed so that the comparison isolates the optimisation and regularisation bundle rather than changing the learning objective itself. & Holding the objective constant makes later Adam gains interpretable as optimisation effects rather than loss-definition effects. \\
Optimizer family & Adam & SGD with step decay; SGD with momentum 0.9 and cosine decay & After the baselines already show learnability, the final route should be less sensitive to a single manually tuned global step size. Adam adapts parameter-wise updates, which is useful when gradients from convolution, attention, fusion, and pointer layers have different scales. & This is the optimizer family used in the paper. The adaptive first/second-moment preconditioner is a better match to heterogeneous curvature than plain SGD. \\
Scheduler family & Lambda warmup to a constant rate & Step decay or cosine annealing & The final route needs a scheduler that protects the unstable early phase without progressively choking learning once span extraction has become meaningful. & Inverse-log warmup prevents overly large early updates before Adam's moment estimates are reliable; the constant post-warmup rate then leaves most late adaptation to Adam itself rather than forcing repeated global decay. \\
Batch size & \texttt{batch\_size=32} & \texttt{batch\_size=16}; original notebook default was 8 & The original batch size 8 gives noisier checkpoint traces. The two SGD baselines therefore move first to 16 so that their update direction is less dominated by mini-batch noise and the effect of step versus cosine decay remains readable. Adam then moves to 32 because that larger batch is still feasible in this notebook and further stabilises both its moment estimates and the monitored \texttt{TEST} trajectory. & Gradient variance scales roughly inversely with batch size. For plain SGD this mainly stabilises the update direction; for Adam it stabilises both \(m_t\) and \(v_t\), so effective per-parameter step sizes become less erratic. The paper also uses batch size 32. \\
Regularisation\newline bundle & \texttt{dropout=0.3},\newline \texttt{dropout\_char=0.15} & Not separately tuned; omitted from the baseline recipes & The notebook dataset is materially smaller than the paper-scale regime, so the final route adds explicit dropout once baseline learnability is already established. The baselines remain lighter because their job is only to show reasonable above-random behaviour, not to carry the full final regularisation bundle. & Dropout reduces co-adaptation in the word-level stack, and character dropout separately discourages over-reliance on narrow subword cues. In a smaller-data regime, stronger dropout is a principled way to prevent faster Adam optimisation from turning into narrow training-set fit. \\
EMA & \texttt{ema\_decay=}\newline \texttt{0.9999} & \texttt{ema\_decay=0} & EMA is reserved for the final route because checkpoint quality matters there; the baselines are not tuned beyond demonstrating learnability. & EMA averages nearby parameter states, suppressing high-frequency fluctuations in late Adam updates and making evaluation closer to a local ensemble / Polyak-style average. This is particularly useful once training loss can still decrease while external checkpoint metrics have begun to plateau. \\
Early-stopping patience & \texttt{early\_stop=50} with \texttt{checkpoint=200} & \texttt{early\_stop=100} with \texttt{checkpoint=200} & With checkpointing every 200 steps, patience 50 already allows 10{,}000 extra updates after the last improvement. That is sufficient for the faster Adam route: once its \texttt{TEST} curve has flattened, doubling the patience mostly extends training further into a region where \texttt{VALID(train)} can still improve without yielding a better external checkpoint. & Shorter patience acts as a model-selection regulariser. A fast adaptive optimizer reaches the flat generalisation regime earlier, so overlong patience mainly increases exposure to late-stage over-specialisation rather than adding useful optimisation time. \\
\bottomrule
\end{tabularx}
\end{center}
\normalsize

The lambda schedule is the clearest point of direct paper alignment. In the paper-aligned implementation, the effective learning rate follows
\[
\mathrm{lr}(t)=
\begin{cases}
\eta \, \dfrac{\log(t+1)}{\log W}, & t < W,\\[0.6em]
\eta, & t \ge W,
\end{cases}
\]
where \(W\) is the warmup horizon and \(\eta\) is the target post-warmup rate. In the logged Adam run, \(W=1000\) and \(\eta = 10^{-3}\), giving the observed sequence \(7.677\times10^{-4}\) at step 200, \(8.677\times10^{-4}\) at step 400, \(9.263\times10^{-4}\) at step 600, \(9.679\times10^{-4}\) at step 800, and \(1.000\times10^{-3}\) from step 1000 onward. The practical effect is that the first 1000 steps are damped while Adam's moment estimates are still forming, after which the run keeps a constant global rate and lets the adaptive preconditioner carry most of the later refinement. Relative to step decay and cosine decay, this makes the final route a deliberately different training bundle rather than a single optimizer swap.

\ReportSubsection{8.4 Adam + lambda: logged \texttt{VALID(train)} / \texttt{TEST} routine}

The Adam route is best read from the logged batch-32 checkpoints. In this subsection, step 40800 is the single selected checkpoint highlighted in the figures because it is the strongest logged \texttt{TEST} F1 point and sits on the same late plateau as the strongest \texttt{TEST} EM. The plotted curves therefore refer to the in-loop batch-32 training record, while the standalone evaluation of \texttt{model\_best.pt} is reported separately afterward.

\begin{center}
\begin{minipage}{0.32\linewidth}
\centering
\SafeIncludeGraphics[width=\linewidth]{prism-uploads/adam_lambda_loss_shared_step_40800_exact_style.png}
\end{minipage}\hfill
\begin{minipage}{0.32\linewidth}
\centering
\SafeIncludeGraphics[width=\linewidth]{prism-uploads/adam_lambda_valid_test_f1_shared_step_40800_exact_style.png}
\end{minipage}\hfill
\begin{minipage}{0.32\linewidth}
\centering
\SafeIncludeGraphics[width=\linewidth]{prism-uploads/adam_lambda_valid_test_em_shared_step_40800_exact_style.png}
\end{minipage}

\vspace{0.35em}
{\small\textbf{Figure.} \texttt{VALID(train)} and \texttt{TEST} F1/EM/Loss trajectories for the batch-32 Adam + lambda run. The vertical marker denotes the selected checkpoint at step 40800.}
\end{center}

\small
\begin{center}
\begin{tabularx}{\linewidth}{>{\raggedright\arraybackslash}p{0.16\linewidth}>{\centering\arraybackslash}p{0.08\linewidth}>{\centering\arraybackslash}p{0.09\linewidth}>{\raggedright\arraybackslash}p{0.16\linewidth}>{\raggedright\arraybackslash}p{0.16\linewidth}>{\raggedright\arraybackslash}X}
\toprule
\textbf{Checkpoint / event} & \textbf{LR} & \textbf{STEP loss} & \textbf{\texttt{VALID(train)}} & \textbf{\texttt{TEST}} & \textbf{Analytical reading} \\
\midrule
Step 200 & 7.68e-4 & 29.701719 & loss 8.212687; F1 8.617219; EM 3.000000 & loss 8.569967; F1 9.610054; EM 3.354167 & Even during warmup, both notebook monitors move off the near-random regime. The run is already extracting answer signal rather than merely producing finite numbers. \\
Step 1000 & 1.00e-3 & 7.278081 & loss 6.769837; F1 15.975063; EM 8.520833 & loss 7.442681; F1 16.006891; EM 8.375000 & Warmup is complete and both printed monitors continue improving together. This is the expected signature of a healthy transition into the main optimisation phase. \\
Step 40800 (selected checkpoint) & 1.00e-3 & 1.544758 & loss 0.696244; F1 93.275025; EM 85.416667 & loss 4.593508; F1 69.634521; EM 59.145833 & This is the selected reporting checkpoint because it gives the strongest logged \texttt{TEST} F1 while remaining on the same late plateau as the best \texttt{TEST} EM. \\
Step 51200 (early stop) & 1.00e-3 & 1.408177 & loss 0.526814; F1 94.782374; EM 88.666667 & loss 4.901060; F1 68.830518; EM 58.500000 & Training-side fit keeps improving after the selected checkpoint, but external checkpoint quality has already flattened. Early stopping truncates that late tail rather than letting Adam continue fitting the training-side monitor indefinitely. \\
Standalone evaluation of \texttt{model\_best.pt} & --- & --- & --- & loss 4.414317; F1 70.243489; EM 59.216436 & The saved checkpoint remains consistent outside the training loop and is slightly stronger in standalone evaluation, confirming that the final model is not an artefact of transient logging. \\
\bottomrule
\end{tabularx}
\end{center}
\normalsize

Three conclusions follow from this routine. First, the warmup phase is behaving correctly: the learning rate rises conservatively to its target value while both printed monitors improve rather than diverge. Second, the loss trajectory (left panel) shows that the lowest \texttt{TEST} loss occurs much earlier than the selected step-40800 checkpoint, whereas \texttt{TEST} F1/EM continue improving into the late plateau. The later part of training is therefore still refining span ranking and joint boundary selection even after token-level negative log-likelihood has stopped identifying the best checkpoint. Third, \texttt{VALID(train)} loss continues to decrease after step 40800 while \texttt{TEST} has already flattened. That is exactly the regime in which the stronger dropout, EMA, and shorter patience become useful: the model can still fit the training-side signal, but that no longer guarantees a better external checkpoint.




%% ====================================================================
%%  SECTION 9 — H1
%% ====================================================================

\ReportSection{9 Hypothesis 1: Mechanistic Divergence Between Architecturally Identical Convolution Layers in QANet's Encoder}

\ReportSubsection{9.1 Motivation and research question}

Each encoder block in QANet's Model Encoder contains four sub-layers---Conv\_0, Conv\_1, Self-Attention, and FFN. Each block follows this pipeline:

\begin{tcolorbox}[enhanced,breakable,colback=CodeBg,colframe=CodeFrame,boxrule=0.5pt,arc=1.5mm,left=2.6mm,right=2.6mm,top=2.0mm,bottom=1.8mm]
\ttfamily
Position Encoding $\rightarrow$ Conv\_0 $\rightarrow$ Conv\_1 $\rightarrow$ Self-Attention $\rightarrow$ FFN\\
\rmfamily
\smallskip
Each sub-layer is followed by residual addition, so later components operate on the accumulated residual stream rather than on isolated features.
\end{tcolorbox}

The QANet paper's ablation study (Table~5, Yu et al.\ 2018) shows that convolution contributes more to performance than Self-Attention ($-2.7$~F1 vs $-1.3$~F1 upon removal, measured after retraining). However, the paper treats both convolution layers as a single unit and does not examine Conv\_0 and Conv\_1 separately.

By separating the two through eval-time ablation and causal tracing, both methods agree that Conv\_1 is the most important component by a large margin. Ablation (measuring necessity) yields:

\begin{center}
\textbf{Conv\_1 $\gg$ Self-Attention $>$ FFN $>$ Conv\_0}
\end{center}

Conv\_1 is the \textbf{most important} of all four component types ($-32.54$~F1 upon removal), while Conv\_0 is the \textbf{least important} ($-1.09$~F1)---a $\Delta$F1 ratio of $32.54 / 1.09 \approx 30\times$. Yet Conv\_0 and Conv\_1 are \textbf{architecturally identical}: both are depthwise separable convolutions with kernel size $k=5$ and the same hidden dimension, trained end-to-end under the same objective.

\textbf{A note on the 30$\times$ ratio.} This is the ratio of eval-time ablation $\Delta$F1 values. The denominator ($-1.09$~F1) is small, so the precise multiplier is sensitive to measurement granularity---if Conv\_0's drop were 0.5~F1 instead, the ratio would be $\sim$65$\times$. Causal tracing gives a different ratio (AIE share: $71.3\% / 15.3\% \approx 4.7\times$), confirming that the precise multiplier is method-dependent. The \textbf{qualitative ordering} (Conv\_1 $\gg$ Conv\_0) is robust across both methods, but the ``30$\times$'' should be read as ``order-of-magnitude difference'' rather than a precise measurement. We retain it as a concise shorthand for the large, reproducible gap.

The same architecture produces both the most important and the least important component. Architecture alone cannot explain this gap; the divergence must relate to their different positions in the pipeline and the different functions they consequently learn during training.

\begin{tcolorbox}[enhanced,breakable,colback=white,colframe=BugBlueBorder,boxrule=0.5pt,arc=1.5mm,left=2.8mm,right=2.8mm,top=1.8mm,bottom=1.8mm]
\textbf{H1 research question.} Conv\_0 and Conv\_1 share identical architecture yet exhibit an order-of-magnitude difference in importance ($\Delta$F1 ratio $\approx 30\times$). What mechanistic differences between them explain this gap?
\end{tcolorbox}

\ReportSubsection{9.2 Experimental design overview}

We address this question through three experiments, each building on the previous:

\textbf{Experiment~1 (Is the gap real?)} We quantify component importance using two independent methods---eval-time ablation and causal tracing---and check for output-magnitude confounds. The goal is to establish that the 30$\times$ gap is reproducible and not an artefact.

\textbf{Experiment~2 (What is the functional difference?)} Experiment~1 confirms that Conv\_1 is far more important than Conv\_0 ($\Delta$F1 ratio $\approx 30\times$). So how do their outputs differ functionally? We first test with linear probes and find that both encode nearly identical task-relevant information (AUC gap only 0.025)---information difference cannot explain an order-of-magnitude importance gap. We therefore examine whether the two layers apply different transformations to their inputs (representational structure).

\textbf{Experiment~3 (How does the model fail?)} Experiment~2 characterises the output differences. What happens inside the model when Conv\_1 is removed? Since Self-Attention is Conv\_1's direct downstream consumer, we capture attention weights and observe how attention patterns change under Conv\_1 removal.

\FieldLabel{9.2.1 Methodological note: eval-time ablation}

All ablations are performed at eval-time (zeroing out component outputs without retraining). Ideally, retraining after ablation would control for compensatory adaptation, but each retraining run requires a full GPU training cycle, infeasible under Colab quota constraints. Since the core objective is to determine the \textbf{relative importance ordering} among components (rather than precise absolute differences), eval-time ablation is sufficient for qualitative conclusions. The QANet paper's retrain-based ablation (Yu et al., 2018) also ranks convolution above Self-Attention, consistent with our ordering.

Under residual architectures, zero-out ablation is the cleanest intervention. Mean or noise replacement injects actively incorrect signals that propagate forward and confound interpretation, whereas zero-out asks the narrower question: what remains once the model is denied this component's contribution?

\ReportSubsection{9.3 Experiment 1: Establishing the gap}

\FieldLabel{9.3.1 Global ablation}

Zero out all 21 instances of a component type (eval-time, full dev set, 10{,}465 samples). Zero-out = sub-layer output replaced with zeros before residual addition.

\small
\begin{center}
{\bfseries Table 9.1. Global component ablation}\\[0.25em]
\begin{tabularx}{0.78\linewidth}{>{\raggedright\arraybackslash}p{0.20\linewidth}>{\centering\arraybackslash}p{0.12\linewidth}>{\centering\arraybackslash}p{0.16\linewidth}>{\centering\arraybackslash}X}
\toprule
\textbf{Config} & \textbf{F1} & \textbf{$\Delta$F1} & \textbf{EM}\\
\midrule
Baseline & 70.24 & --- & 59.22\\
$-$Conv\_1 & 37.70 & \textbf{$-32.54$} & 23.59\\
$-$Conv\_0 & 69.16 & \textbf{$-1.09$} & 57.85\\
$-$Attn & 65.44 & $-4.81$ & 55.18\\
$-$FFN & 66.92 & $-3.32$ & 54.56\\
$-$Conv (both) & 27.14 & $-43.11$ & 13.60\\
\bottomrule
\end{tabularx}
\end{center}
\normalsize

\textbf{Finding~1}: Conv\_1 is the most important component; Conv\_0 is the least. The $\Delta$F1 ratio between architecturally identical layers is $32.54 / 1.09 \approx$ \textbf{30$\times$} (see note on ratio sensitivity in Section~9.1).

\textbf{Finding~2}: Joint removal ($-43.11$) exceeds the sum of individuals ($-33.63$) by 9.48~F1---a super-additive interaction. Conv\_0 contributes non-redundant preprocessing that only becomes critical when Conv\_1 is also absent.

\FieldLabel{9.3.2 Per-block ablation}

Remove conv (both) or self\_attn in one (pass, block) at a time: $3 \text{ passes} \times 7 \text{ blocks} \times 2 \text{ component types} = 42$ configs, full dev set (10{,}465 samples).

\textbf{Finding~3}: Max single-block impact = $-1.23$~F1 (p0\_b0\_conv); collective removal = $-43.11$. A 35$\times$ gap reveals \textbf{distributed redundancy}: no single block is essential, but together they are indispensable.

\textbf{Finding~4}: Early passes matter more. Averaging per-block $\Delta$F1 across all component types within each pass: $M_1$ ($-0.16$ avg) $\gg$ $M_2$ ($-0.06$) $\gg$ $M_3$ ($+0.01$).

\FieldLabel{9.3.3 Causal tracing (Meng et al., 2022)}

Per sample: (1)~clean run $\to$ collect all 84 sub-layer activations ($3 \text{ passes} \times 7 \text{ blocks} \times 4 \text{ components}$); (2)~corrupt context embeddings (Gaussian, $\sigma = 3 \times \text{std}$); (3)~restore ONE sub-layer's clean activation into the corrupted run $\to$ measure recovery. The Average Indirect Effect (AIE) for each sub-layer $i$ is:
\[
\text{AIE}_i = P(\text{answer} \mid \text{corrupted} + \text{restore}_i) - P(\text{answer} \mid \text{corrupted})
\]
A higher AIE means that sub-layer $i$ carries more of the causal information needed for the correct answer. 200 samples, 3 noise repeats.

\small
\begin{center}
{\bfseries Table 9.2. Causal tracing results (AIE)}\\[0.25em]
\begin{tabularx}{0.72\linewidth}{>{\raggedright\arraybackslash}p{0.20\linewidth}>{\centering\arraybackslash}p{0.16\linewidth}>{\centering\arraybackslash}p{0.16\linewidth}>{\centering\arraybackslash}X}
\toprule
\textbf{Component} & \textbf{Mean AIE} & \textbf{Sum AIE} & \textbf{Share}\\
\midrule
Conv\_1 & 0.0147 & 0.3089 & 71.3\%\\
Conv\_0 & 0.0032 & 0.0664 & 15.3\%\\
Self-Attn & 0.0016 & 0.0330 & 7.6\%\\
FFN & 0.0012 & 0.0246 & 5.7\%\\
\bottomrule
\end{tabularx}
\end{center}
\normalsize

\textbf{Finding~5}: Causal tracing independently confirms the gap---Conv\_1 accounts for 71.3\% of total AIE. Conv\_0 ranks second (15.3\%), above Attn and FFN, suggesting it does carry some causal weight even though ablation damage is minimal. This asymmetry (small ablation damage, moderate causal recovery) is consistent with Conv\_0 providing preprocessing that is useful but compensable via residual connections.

\begin{center}
\SafeIncludeGraphics[width=0.92\linewidth]{prism-uploads/h1_fig1_causal_tracing_heatmap.png}\\[0.3em]
{\small\textbf{Figure 9.1.} Causal tracing heatmap: AIE per sub-layer across 3 passes $\times$ 7 blocks. Conv\_1 (top-right) dominates all other panels under a shared colour scale. Early passes contribute more AIE than later passes, consistent with Finding~4.}
\end{center}

\FieldLabel{9.3.4 Confound check: output magnitude}

Sub-layer L2 norms: Conv\_1 (250) $>$ Conv\_0 (113) $>$ FFN (104) $>$ Attn (82). Pearson $r(\text{norm}, \text{AIE}) = 0.893$.

Could the gap simply reflect Conv\_1 having larger outputs? If magnitude were the driver, the norm ordering Conv\_0 (113) $>$ FFN (104) $>$ Attn (82) should predict ablation ordering Conv\_0 $>$ FFN $>$ Attn. The actual ablation ordering is exactly reversed: Attn ($-4.81$) $>$ FFN ($-3.32$) $>$ Conv\_0 ($-1.09$)---\textbf{inversely} correlated. This rules out output magnitude as the driver for these three components. Conv\_1's dominance direction is robust regardless of magnitude; the precise ratio over Attn (2--7$\times$) carries uncertainty from norm co-correlation.

\FieldLabel{9.3.5 Experiment 1 summary}

The order-of-magnitude importance gap between identical architectures is real and reproducible across two independent methods. It cannot be explained by output magnitude alone. Something about what Conv\_1 \emph{learns to do} is fundamentally different from Conv\_0.

\begin{center}
\SafeIncludeGraphics[width=0.72\linewidth]{prism-uploads/h1_fig2_dual_method.png}\\[0.3em]
{\small\textbf{Figure 9.2.} Cross-validation: both ablation (|$\Delta$F1| share, red) and causal tracing (AIE share, blue) independently identify Conv\_1 as the dominant component. Rankings agree at the top ($>$70\% for Conv\_1) while differing on the remaining three---Conv\_0 is lowest in ablation but second in causal tracing, reflecting its compensable-but-causal role.}
\end{center}

\ReportSubsection{9.4 Experiment 2: Functional differences between Conv\_1 and Conv\_0}

\textbf{Question}: Conv\_1 is far more important than Conv\_0 ($\Delta$F1 ratio $\approx 30\times$). How do their outputs differ functionally?

\FieldLabel{9.4.1 Information content (linear probe)}

Linear probing is a standard technique in representation analysis (Alain \& Bengio, 2017; Belinkov et al., 2017). We train logistic regression on per-token sub-layer outputs \textbf{before residual addition}---i.e., the sub-layer's own transformation, isolated from the skip-connection input. This is a deliberate choice: probing the post-residual stream would mix the sub-layer's contribution with accumulated upstream information, making all components appear similar regardless of their actual transformation. Label: inside answer span. 500 samples, 80/20 split, 5 (pass, block) positions.

\small
\begin{center}
{\bfseries Table 9.3. Linear probe AUC (answer-position)}\\[0.25em]
\begin{tabularx}{0.52\linewidth}{>{\raggedright\arraybackslash}p{0.24\linewidth}>{\centering\arraybackslash}X}
\toprule
\textbf{Component} & \textbf{Mean AUC}\\
\midrule
Conv\_1 & 0.901\\
FFN & 0.898\\
Self-Attn & 0.894\\
Conv\_0 & 0.876\\
\bottomrule
\end{tabularx}
\end{center}
\normalsize

\textbf{Finding~6}: On the answer-position property we tested, all components encode similar information \emph{even in their isolated transformations} (before residual addition). The Conv\_1 vs Conv\_0 gap is only 0.025~AUC, negligible compared to their order-of-magnitude ablation gap. Since we probed pre-residual outputs, this similarity cannot be attributed to the skip-connection carrying accumulated information---it reflects the sub-layers' own transformations encoding comparable answer-position signal.

\textbf{Limitation}: We only probed one property (answer position). Conv\_1 may encode critical information that Conv\_0 does not on other task-relevant properties (e.g., syntactic boundaries, entity types)---these remain untested. The ``similar information content'' conclusion is strictly limited to the tested property---\textbf{the gap's root cause may partly lie in untested information differences, partly in representational structure}. Additionally, linear probes can only capture linearly separable features; non-linearly encoded differences may be missed.

\FieldLabel{9.4.2 Representational structure (discriminativity)}

Compare Conv\_1 vs Conv\_0 output properties directly. 200 samples, 5 (pass, block) positions. Metrics: local contrast ($1 - \text{lag-1 cosine similarity}$), token norm variance (CoV), effective rank ($\exp$ of singular value entropy).

\small
\begin{center}
{\bfseries Table 9.4. Representational structure comparison}\\[0.25em]
\begin{tabularx}{0.96\linewidth}{>{\raggedright\arraybackslash}p{0.24\linewidth}>{\centering\arraybackslash}p{0.12\linewidth}>{\centering\arraybackslash}p{0.12\linewidth}>{\centering\arraybackslash}p{0.18\linewidth}>{\centering\arraybackslash}p{0.10\linewidth}>{\centering\arraybackslash}X}
\toprule
\textbf{Metric} & \textbf{Conv\_0} & \textbf{Conv\_1} & \textbf{Conv\_1/Conv\_0} & \textbf{Attn} & \textbf{FFN}\\
\midrule
Token Norm Variance & 4.57 & 14.48 & \textbf{3.2$\times$} & 1.62 & 14.79\\
Norm CoV & 0.345 & 0.216 & 0.6$\times$ & 0.176 & 0.410\\
Local Contrast ($1-\cos$) & 0.323 & 0.305 & 0.9$\times$ & 0.103 & 0.341\\
Effective Rank & 32.5 & 43.2 & \textbf{1.3$\times$} & 8.6 & 26.5\\
\bottomrule
\end{tabularx}
\end{center}
\normalsize

\textbf{Finding~7}: Conv\_1 is significantly higher than Conv\_0 on two dimensions: token-level amplitude variation (Norm Variance 3.2$\times$) and output subspace dimensionality (Effective Rank 1.3$\times$). However, in directional distinctiveness (Local Contrast) and normalised variation (Norm CoV), Conv\_0 is actually slightly higher. Conv\_1's transformation is primarily characterised by \textbf{amplitude modulation and dimensional utilisation}, rather than directional separation of adjacent tokens. Conv\_0's outputs, while more directionally diverse, have minimal overall impact (ablation only $-1.09$~F1), suggesting this directional diversity does not critically contribute to downstream computation.

\FieldLabel{9.4.3 Spatial structure (local coherence)}

Mean cosine similarity between output vectors at positions $t$ and $t+k$, at p0\_b0, 500 samples.

\small
\begin{center}
{\bfseries Table 9.5. Spatial coherence (cosine similarity by lag)}\\[0.25em]
\begin{tabularx}{0.86\linewidth}{>{\raggedright\arraybackslash}p{0.16\linewidth}>{\centering\arraybackslash}p{0.10\linewidth}>{\centering\arraybackslash}p{0.10\linewidth}>{\centering\arraybackslash}p{0.10\linewidth}>{\centering\arraybackslash}p{0.10\linewidth}>{\centering\arraybackslash}X}
\toprule
\textbf{Component} & \textbf{lag=1} & \textbf{lag=2} & \textbf{lag=3} & \textbf{lag=5} & \textbf{decay(1$\to$5)}\\
\midrule
Conv\_0 & 0.679 & 0.602 & 0.597 & 0.571 & 0.109\\
Conv\_1 & 0.667 & 0.585 & 0.610 & 0.608 & 0.058\\
Self-Attn & 0.841 & 0.819 & 0.810 & 0.804 & 0.037\\
FFN & 0.595 & 0.564 & 0.562 & 0.556 & 0.039\\
\bottomrule
\end{tabularx}
\end{center}
\normalsize

\textbf{Finding~8}: Conv\_0's decay (0.109) is nearly double Conv\_1's (0.058), indicating Conv\_0's output has stronger locality. Conv\_1's flatter decay means its output is more spatially uniform. Self-Attention has the highest overall similarity (0.84) with near-zero decay (0.037), consistent with its global mixing role.

\FieldLabel{9.4.4 Experiment 2 summary}

The functional difference between Conv\_1 and Conv\_0---at least on the answer-position property---is not in information content (AUC gap only 0.025) but in \textbf{specific dimensions of representational structure}. Conv\_1's output has significantly higher token-level amplitude variation (Norm Variance 3.2$\times$) and subspace dimensionality (Effective Rank 1.3$\times$), with more spatially uniform structure (decay 0.058 vs 0.109).

The divergence across metrics has diagnostic value. Conv\_0 scores slightly higher on directional distinctiveness yet is nearly dispensable ($-1.09$~F1)---ruling out directional distinctiveness as a sufficient condition. The two dimensions where Conv\_1 holds a clear advantage---amplitude modulation and subspace dimensionality---are candidate mechanisms. This is consistent with dot-product attention mechanics: $QK^{\!\top}$ scores are sensitive to input magnitude differences, and higher-dimensional inputs allow $Q$, $K$ projections to carry richer discriminative signals. Notably, FFN provides a natural control: its Norm Variance (14.79) is nearly identical to Conv\_1's (14.48), so if Norm Variance alone were sufficient for importance, FFN should match Conv\_1 in ablation impact. It does not ($-3.32$ vs $-32.54$~F1), precisely because FFN sits \emph{downstream} of Self-Attention---its output properties cannot shape $QK^{\!\top}$ scores. This confirms that amplitude modulation's causal impact depends critically on pipeline position: the same output property matters when it feeds into attention, and is irrelevant when it does not.

\ReportSubsection{9.5 Experiment 3: How does the model fail without Conv\_1?}

\textbf{Question}: Removing Conv\_1 causes a $-32.54$~F1 collapse. What happens inside the model?

Self-Attention is Conv\_1's direct downstream consumer. We use forward hooks to capture attention weights under three conditions (50 samples; sufficient because the per-block JSD ratio is tightly concentrated at 6.9--10.4$\times$ across all 7 blocks, indicating low variance): clean, $-$Conv\_1, $-$Conv\_0 (control).

\FieldLabel{9.5.1 JS divergence (clean vs ablated)}

\small
\begin{center}
{\bfseries Table 9.6. Attention JSD under component removal}\\[0.25em]
\begin{tabularx}{0.72\linewidth}{>{\centering\arraybackslash}p{0.12\linewidth}>{\centering\arraybackslash}p{0.22\linewidth}>{\centering\arraybackslash}p{0.22\linewidth}>{\centering\arraybackslash}X}
\toprule
\textbf{Block} & \textbf{JSD($-$Conv\_1)} & \textbf{JSD($-$Conv\_0)} & \textbf{Ratio}\\
\midrule
0 & 0.0331 & 0.0041 & 8.1$\times$\\
1 & 0.0870 & 0.0110 & 7.9$\times$\\
2 & 0.0749 & 0.0108 & 7.0$\times$\\
3 & 0.1222 & 0.0118 & 10.3$\times$\\
4 & 0.0946 & 0.0118 & 8.0$\times$\\
5 & 0.0888 & 0.0110 & 8.1$\times$\\
6 & 0.0898 & 0.0112 & 8.0$\times$\\
\textbf{Avg} & \textbf{0.0843} & \textbf{0.0102} & \textbf{8.2$\times$}\\
\bottomrule
\end{tabularx}
\end{center}
\normalsize

\textbf{Finding~9}: Attention distortion under Conv\_1 removal is \textbf{8.2$\times$ greater} than under Conv\_0 removal. The ratio is consistent across all 7 blocks (6.9--10.4$\times$).

\begin{center}
\SafeIncludeGraphics[width=0.76\linewidth]{prism-uploads/h1_fig3_attention_jsd.png}\\[0.3em]
{\small\textbf{Figure 9.3.} JS divergence of attention weights (clean vs ablated) across 7 encoder blocks. Removing Conv\_1 (red) causes 7--10$\times$ more distortion than removing Conv\_0 (blue) in every block, demonstrating Conv\_1's role as the critical upstream supplier of Self-Attention features.}
\end{center}

\FieldLabel{9.5.2 Attention entropy and answer focus}

\small
\begin{center}
{\bfseries Table 9.7. Entropy and answer-focus change under component removal}\\[0.25em]
\begin{tabularx}{0.86\linewidth}{>{\raggedright\arraybackslash}p{0.16\linewidth}>{\centering\arraybackslash}p{0.14\linewidth}>{\centering\arraybackslash}p{0.16\linewidth}>{\centering\arraybackslash}p{0.16\linewidth}>{\centering\arraybackslash}X}
\toprule
\textbf{Condition} & \textbf{Entropy} & \textbf{$\Delta$Entropy} & \textbf{Answer Mass} & \textbf{$\Delta$Mass}\\
\midrule
Clean & 4.621 & --- & 0.0463 & ---\\
$-$Conv\_1 & 4.362 & $-0.259$ & 0.0373 & $-19\%$\\
$-$Conv\_0 & 4.554 & $-0.067$ & 0.0489 & $+6\%$\\
\bottomrule
\end{tabularx}
\end{center}
\normalsize

\textbf{Finding~10}: Without Conv\_1, attention does not diffuse---it \textbf{collapses} into over-concentrated patterns on wrong positions (entropy \emph{decreases}, answer mass drops 19\%). The model does not lose track of information; instead, attention becomes erroneously focused on non-answer positions.

\FieldLabel{9.5.3 On pipeline distance}

An intuitive objection: does Conv\_1 removal affect attention more simply because Conv\_1 is Self-Attention's immediate upstream neighbour? Our existing data addresses this. Conv\_0 is also an immediate upstream neighbour---of Conv\_1 (0 layers apart). Yet removing Conv\_0 causes only $-1.09$~F1. Experiment~2 shows that Conv\_0 does apply a non-trivial transformation, but its contribution is compensable: the residual connection ensures Conv\_1 still receives the core information flow.

Therefore, \textbf{pipeline proximity is necessary but not sufficient for large impact}. What matters is whether the downstream layer \emph{irreplaceably depends} on the upstream layer's specific output. Conv\_1 $\to$ Attention: attention critically relies on Conv\_1's high-amplitude, high-dimensional features; removal causes pattern collapse. Conv\_0 $\to$ Conv\_1: Conv\_1 does not critically depend on Conv\_0's specific output; its contribution is compensable via the residual pathway.

\ReportSubsection{9.6 H1 synthesis}

\textbf{Direct answer to H1.} The directly observable cause lies in the different transformations Conv\_1 and Conv\_0 apply, and Self-Attention's irreplaceable dependence on Conv\_1's specific output. On the answer-position property, both layers encode similar information (AUC gap only 0.025), but Conv\_1 produces outputs with significantly higher token-level amplitude variation (Norm Variance 3.2$\times$) and subspace dimensionality (Effective Rank 1.3$\times$). Conv\_0's transformation, while non-trivial, is compensable ($-1.09$~F1). When Conv\_1 is removed, Self-Attention's attention patterns collapse onto wrong positions (JSD 8.2$\times$ vs Conv\_0 control, answer focus drops 19\%)---this is the proximate mechanism of the performance breakdown.

\small
\begin{center}
{\bfseries Table 9.8. Chain of evidence for H1}\\[0.25em]
\begin{tabularx}{0.96\linewidth}{>{\raggedright\arraybackslash}p{0.24\linewidth}>{\raggedright\arraybackslash}p{0.30\linewidth}>{\raggedright\arraybackslash}X}
\toprule
\textbf{Evidence layer} & \textbf{Main result} & \textbf{Contribution to the final answer}\\
\midrule
Global ablation (Exp~1a) & Conv\_1 $-32.54$~F1, Conv\_0 $-1.09$~F1; joint removal super-additive. & Establishes the 30$\times$ gap across architecturally identical layers.\\
Per-block ablation (Exp~1b) & No single block dominates; early passes matter more. & Indicates distributed redundancy, not one brittle bottleneck.\\
Causal tracing (Exp~1c) & Conv\_1 accounts for 71.3\% of total AIE. & Confirms the ordering with a restoration-based method.\\
Linear probes (Exp~2a) & AUC gap only 0.025 on answer-position property. & Rules out information difference on this property as the primary explanation; other properties may still differ.\\
Representational analysis (Exp~2b/c) & Conv\_1: Norm Variance 3.2$\times$, Effective Rank 1.3$\times$, flatter spatial decay. & Points to the real mechanism: amplitude modulation and subspace utilisation.\\
Downstream attention (Exp~3) & JSD 8.2$\times$; entropy $-0.259$; answer mass $-19\%$. & Identifies the proximate failure mode: attention collapses onto wrong positions.\\
\bottomrule
\end{tabularx}
\end{center}
\normalsize

Three limitations should be kept explicit. First, \textbf{position and learned function are confounded}: Conv\_1's behaviour may arise from its pipeline position, its learned weights, or their interaction; current eval-time interventions cannot fully separate those factors. Second, \textbf{limited probe coverage}: we only probed answer position; Conv\_1 and Conv\_0 may differ in encoding of other task-relevant properties. The ``similar information'' conclusion applies strictly to the tested property and does not rule out meaningful differences on untested attributes. Third, \textbf{propagation effect of global ablation}: globally removing Conv\_1 alters the entire residual stream; the observed attention degradation includes both direct and indirect effects.

%% ====================================================================
%%  SECTION 10 — H2
%% ====================================================================

\ReportSection{10 Hypothesis 2: Directional Asymmetry in CQ Attention --- Cross-Architecture Validation from BiDAF to QANet}

\ReportSubsection{10.1 Motivation and research question}

QANet inherits its bidirectional CQ Attention design from BiDAF (Seo et al., 2017). This layer fuses the independently encoded Context and Question streams into a single representation. The output is a concatenation of four sub-components:

\begin{tcolorbox}[enhanced,breakable,colback=CodeBg,colframe=CodeFrame,boxrule=0.5pt,arc=1.5mm,left=2.6mm,right=2.6mm,top=2.0mm,bottom=1.8mm]
\ttfamily
Output = [C;\; A;\; C$\odot$A;\; C$\odot$B] \quad--- total 4$d$ dimensions\\[0.3em]
\hphantom{Output = [}raw\quad C$\to$Q\quad gated\quad Q$\to$C$\to$C
\rmfamily
\end{tcolorbox}

where $A = S_1 Q$ (C$\to$Q direction), $C \odot B$ derives from $B = S_1 S_2^{\!\top} C$ (Q$\to$C direction), $S_1$ is softmax over the Q dimension, and $S_2$ is softmax over the C dimension. Both directions share the same similarity matrix $S$ and trainable weight $w$; they differ only in the softmax normalisation direction.

BiDAF (Seo et al., 2017, Table~1b) measured the relative importance of the two directions via \emph{training-time} ablation (removing a component and retraining the model):

\small
\begin{center}
{\bfseries Table 10.1. BiDAF direction-level ablation (Seo et al., 2017)}\\[0.25em]
\begin{tabularx}{0.70\linewidth}{>{\raggedright\arraybackslash}p{0.35\linewidth}>{\centering\arraybackslash}p{0.15\linewidth}>{\centering\arraybackslash}X}
\toprule
\textbf{Condition} & \textbf{F1} & \textbf{$\Delta$F1}\\
\midrule
BiDAF (single) & 77.3 & ---\\
Remove C2Q ($A$) & 67.7 & $-9.6$\\
Remove Q2C ($B$) & 73.7 & $-3.6$\\
\bottomrule
\end{tabularx}
\end{center}
\normalsize

The C2Q direction is approximately \textbf{2.7$\times$} more important than Q2C. However, BiDAF uses LSTMs as its encoder, while QANet replaces them entirely with convolutions and self-attention. \emph{Does the asymmetry still hold after the architecture change?}

This question is non-trivial: QANet's Model Encoder ($7 \text{ blocks} \times 3 \text{ passes} = 21$ layers) has far greater post-fusion processing capacity than BiDAF's two-layer BiLSTM, and could potentially compensate from the weaker direction's coarser signal. If QANet's stronger encoder ``flattens'' the directional gap, the asymmetry is an artefact of LSTM capacity bottleneck; if it persists or intensifies, it is an \emph{inherent property} of the attention fusion mechanism.

Furthermore, BiDAF's ablation granularity was direction-level (C2Q vs Q2C), without separating the independent contributions of $[C,\; A,\; C \odot A,\; C \odot B]$. We revisit this question at finer granularity.

\begin{tcolorbox}[enhanced,breakable,colback=white,colframe=BugBlueBorder,boxrule=0.5pt,arc=1.5mm,left=2.8mm,right=2.8mm,top=1.8mm,bottom=1.8mm]
\textbf{H2 research question.} Does BiDAF's C2Q $>$ Q2C directional asymmetry (2.7$\times$) hold in QANet's Conv+Attention architecture? If so, what mechanism maintains it across architectures?
\end{tcolorbox}

\ReportSubsection{10.2 Methodological note: eval-time intervention}

BiDAF's ablation experiments used \textbf{training-time ablation}---removing a component and retraining the entire model. This answers an architectural design question: \emph{Is this component architecturally replaceable?}

We choose \textbf{eval-time ablation}---zeroing sub-component outputs directly on a fully trained model. This is a deliberate methodological choice, not a fallback due to resource constraints. Our research goal is \textbf{mechanistic analysis}---understanding the functional roles that components \emph{currently serve} in the trained model---not architecture search. This aligns with the causal tracing methodology of Meng et al.\ (2022) for localising factual associations in GPT: intervening on a trained model (corruption + restoration) to localise functional roles. Training-time ablation allows the model to learn compensatory strategies during retraining, erasing the functional traces of the removed component---making it \emph{unsuitable} for mechanistic analysis.

\FieldLabel{10.2.1 How the comparison with BiDAF holds}

Due to methodological differences, the absolute $\Delta$F1 values are \textbf{not directly comparable} (eval-time ablation tends to produce larger drops due to distribution shift). However, the methodological bias affects both C2Q and Q2C directions approximately equally---both are attention sub-components, and zeroing produces comparable distribution shift. Therefore, the \textbf{C2Q/Q2C direction ratio} is a robust cross-method comparison metric. Our comparison is strictly limited to the ratio level.

\ReportSubsection{10.3 Pre-validation: CQ Attention is a critical bottleneck}

Before dissecting CQ Attention internally, we first validate its importance as a whole. CQ Attention is the sole fusion point of the two streams---architecturally positioned as a bottleneck. We confirm this by zeroing its entire output ($X = 0$) at eval time and measuring the F1/EM drop.

\small
\begin{center}
{\bfseries Table 10.2. CQ Attention bottleneck validation}\\[0.25em]
\begin{tabularx}{0.75\linewidth}{>{\raggedright\arraybackslash}p{0.30\linewidth}>{\centering\arraybackslash}p{0.12\linewidth}>{\centering\arraybackslash}p{0.12\linewidth}>{\centering\arraybackslash}X}
\toprule
\textbf{Condition} & \textbf{F1} & \textbf{EM} & \textbf{$\Delta$F1}\\
\midrule
Baseline & 70.24 & 59.22 & ---\\
CQ output zeroed & 7.11 & --- & $-63.13$ ($-89.9\%$)\\
\bottomrule
\end{tabularx}
\end{center}
\normalsize

F1 drops by 89.9\%. CQ Attention carries nearly all information the model's performance depends on. Dissecting its internal sub-components is well-motivated.

\ReportSubsection{10.4 Experiment 1: Does directional asymmetry hold?}

\textbf{Goal}: Measure the causal importance of four output sub-components, aggregate to direction-level, and cross-compare with BiDAF.

\FieldLabel{10.4.1 Sub-component ablation}

CQ Attention output $X$ has shape $[B, 4d, L]$. The four sub-components are arranged by dimension: $X_{[:,:d]}$ is $C$, $X_{[:,d:2d]}$ is $A$, $X_{[:,2d:3d]}$ is $C \odot A$, $X_{[:,3d:4d]}$ is $C \odot B$. Before $X$ enters \texttt{cq\_resizer}, we zero the target range while keeping the rest unchanged. Evaluation uses the full dev set.

\small
\begin{center}
{\bfseries Table 10.3. Sub-component ablation results}\\[0.25em]
\begin{tabularx}{0.96\linewidth}{>{\raggedright\arraybackslash}p{0.12\linewidth}>{\raggedright\arraybackslash}p{0.18\linewidth}>{\centering\arraybackslash}p{0.12\linewidth}>{\centering\arraybackslash}p{0.10\linewidth}>{\centering\arraybackslash}X}
\toprule
\textbf{Condition} & \textbf{Zeroed content} & \textbf{Direction} & \textbf{F1} & \textbf{$\Delta$F1}\\
\midrule
Baseline & None & --- & 70.24 & ---\\
$-C$ & Raw context & No attention & 61.73 & $-8.51$\\
$-A$ & $S_1 Q$ & C$\to$Q & 66.20 & $-4.04$\\
$-(C \odot A)$ & Gated & C$\to$Q gated & 54.41 & $-15.83$\\
$-(C \odot B)$ & $S_1 S_2^{\!\top} C$ & Q$\to$C$\to$C & 69.21 & $-1.03$\\
Only $C$ & $[C;\, 0;\, 0;\, 0]$ & No attention & 8.35 & $-61.90$\\
\bottomrule
\end{tabularx}
\end{center}
\normalsize

\FieldLabel{10.4.2 Direction-level aggregation and BiDAF comparison}

Aggregate sub-component results to direction level: C2Q = simultaneously zero $A$ and $C \odot A$; Q2C = zero $C \odot B$.

\small
\begin{center}
{\bfseries Table 10.4. Cross-architecture direction-level comparison}\\[0.25em]
\begin{tabularx}{0.82\linewidth}{>{\raggedright\arraybackslash}p{0.34\linewidth}>{\centering\arraybackslash}p{0.22\linewidth}>{\centering\arraybackslash}X}
\toprule
\textbf{Comparison} & \textbf{BiDAF (LSTM)} & \textbf{QANet (Conv+Attn)}\\
\midrule
$-$C2Q (remove $A$ + $C \odot A$) & $-9.6$ F1 & $-62.78$ F1\\
$-$Q2C (remove $C \odot B$) & $-3.6$ F1 & $-1.03$ F1\\
C2Q / Q2C ratio & 2.7$\times$ & \textbf{61.0$\times$}\\
\bottomrule
\end{tabularx}
\end{center}
\normalsize

\FieldLabel{10.4.3 cq\_resizer weight analysis}

\texttt{cq\_resizer} is \texttt{Conv1d(4d, d, kernel\_size=1)}, with weight $W$ of shape $[d, 4d]$. We partition $W$ by columns into four $[d, d]$ quadrants and compute each quadrant's Frobenius norm. Ablation measures \emph{performance loss} from removing a component (output side); weight norm measures the model's \emph{learned dependency} (parameter side). If ablation ranking and weight norm ranking agree, conclusion credibility is significantly strengthened.

\small
\begin{center}
{\bfseries Table 10.5. cq\_resizer weight quadrant Frobenius norms}\\[0.25em]
\begin{tabularx}{0.82\linewidth}{>{\raggedright\arraybackslash}p{0.18\linewidth}>{\raggedright\arraybackslash}p{0.22\linewidth}>{\centering\arraybackslash}p{0.20\linewidth}>{\centering\arraybackslash}X}
\toprule
\textbf{Quadrant} & \textbf{Sub-component} & \textbf{Frobenius norm} & \textbf{norm\,/\,$\|W\|_F$}\\
\midrule
$W_{[:,:d]}$ & $C$ (raw) & 12.42 & 44.6\%\\
$W_{[:,d:2d]}$ & $A$ (C$\to$Q) & \textbf{18.44} & \textbf{66.2\%}\\
$W_{[:,2d:3d]}$ & $C \odot A$ (gated) & 11.90 & 42.8\%\\
$W_{[:,3d:4d]}$ & $C \odot B$ (Q$\to$C) & 11.78 & 42.3\%\\
\bottomrule
\end{tabularx}
\end{center}
\normalsize

\FieldLabel{10.4.4 Experiment 1 summary}

\textbf{Finding 1 (Sub-component ranking).} The causal importance ranking is $C \odot A$ ($-15.83$) $\gg$ $C$ ($-8.51$) $>$ $A$ ($-4.04$) $\gg$ $C \odot B$ ($-1.03$). The gated term $C \odot A$ is the single most important sub-component; the Q2C direction ($C \odot B$) is nearly negligible.

\textbf{Finding 2 (Cross-architecture comparison).} QANet's direction-level C2Q/Q2C ratio is \textbf{61.0$\times$} (BiDAF: 2.7$\times$). This headline ratio involves unequal intervention scales (C2Q zeros 2/4 dimensions, Q2C zeros 1/4); the controlled single-component ratio is 15.4$\times$ (see Intervention size control below). Both ratios confirm that the asymmetry not only \emph{persists} after the architecture change but \emph{dramatically intensifies}. Removing C2Q drops F1 to 7.46 (approaching 7.11 from full-layer removal); Q2C ($C \odot B$) contributes only $-1.03$ F1, nearly redundant in QANet.

\begin{center}
\SafeIncludeGraphics[width=0.88\linewidth]{prism-uploads/h2_fig1_cross_arch.png}\\[0.3em]
{\small\textbf{Figure 10.1.} Cross-architecture comparison. Left: absolute $\Delta$F1 by direction. Right: C2Q/Q2C ratio jumps from 2.7$\times$ (BiDAF) to 61$\times$ (QANet)---a 23$\times$ amplification. The asymmetry is not an artefact of LSTM capacity; it is an inherent property of the fusion mechanism that \emph{intensifies} with stronger downstream encoding.}
\end{center}

\textbf{Finding 3 (Weight validation).} $A$ (C$\to$Q main term) has the highest Frobenius norm (18.44), significantly above the other three quadrants (11.78--12.42), \textbf{supporting} C2Q dominance from the parameter side. However, $C \odot A$ and $C \odot B$ have similar norms (11.90 vs 11.78), failing to predict their $15\times$ ablation difference. Weight norms agree with ablation at the direction level ($A$ largest $\to$ C2Q dominant), but at sub-component level, input signal quality equally determines causal effect.

\textbf{Finding 4 ($C$ component role).} Raw context passthrough ($C$) is the second most important sub-component ($-8.51$ F1), indicating the model relies on both attention-processed signals and unmodified raw context.

\textbf{Intervention size control.} In the direction-level comparison, C2Q ablation zeros $2/4$ dimensions ($A$ + $C \odot A$) while Q2C zeros only $1/4$ dimension ($C \odot B$)---different intervention scales. To rule out this confound, we check single-component comparisons (each zeroing exactly $1/4$ of dimensions): $C \odot A$ ($-15.83$) vs $C \odot B$ ($-1.03$) = \textbf{15.4$\times$}, $A$ ($-4.04$) vs $C \odot B$ ($-1.03$) = \textbf{3.9$\times$}. Even with strictly controlled intervention size, every C2Q sub-component significantly outweighs Q2C. The asymmetry holds at any granularity.

\ReportSubsection{10.5 Experiment 2: What mechanism maintains the asymmetry?}

\textbf{Question}: C2Q is more important, but C2Q and Q2C share the same similarity matrix $S$, differing only in softmax direction. Why does merely changing the normalisation direction produce such a massive causal difference?

Starting from the mathematical definitions, the two directions perform fundamentally different operations:

\begin{itemize}
\item \textbf{C2Q ($A = S_1 Q$)}: \emph{Injects} Question embeddings into every Context position---introducing vector content that \textbf{did not previously exist} in Context.
\item \textbf{Q2C ($C \odot B$, where $B = S_1 S_2^{\!\top} C$)}: \emph{Redistributes} Context vectors among Context positions, mediated by Question---does not introduce Question's vector content, only reorganises existing Context information using Question-guided attention patterns. Q2C injects Question-guided \emph{structure}, not Question \emph{content}.
\end{itemize}

If answer-relevant information in Context is spatially concentrated, then ``redistributing existing Context information'' has low value---useful information is already at the correct positions. Meanwhile, ``injecting Question information'' is irreplaceable---the Question signal has no other pathway into CQ Attention's output.

We test the information concentration premise through selective corruption.

\FieldLabel{10.5.1 Method}

Apply noise ($\sigma = 3 \times \text{std}$) only to specific Context positions, keeping others clean. 200 samples, 3 noise repeats per sample.

Conditions: \textbf{ans\_only} corrupts only answer-span positions ($y_1$ to $y_2$); \textbf{non\_ans\_only} corrupts only non-answer, non-padding positions; \textbf{full\_context} corrupts all Context positions (baseline reference).

\small
\begin{center}
{\bfseries Table 10.6. Selective corruption results}\\[0.25em]
\begin{tabularx}{0.92\linewidth}{>{\raggedright\arraybackslash}p{0.18\linewidth}>{\centering\arraybackslash}p{0.13\linewidth}>{\centering\arraybackslash}p{0.13\linewidth}>{\centering\arraybackslash}p{0.15\linewidth}>{\centering\arraybackslash}X}
\toprule
\textbf{Condition} & \textbf{Mean TE} & \textbf{$\pm$95\% CI} & \textbf{Avg \#tokens} & \textbf{TE/token}\\
\midrule
ans\_only & 0.170 & $\pm$0.034 & 3.1 & $55.7 \times 10^{-3}$\\
non\_ans\_only & 0.266 & $\pm$0.041 & 140.5 & $1.89 \times 10^{-3}$\\
full\_context & 0.361 & $\pm$0.045 & 143.6 & $2.52 \times 10^{-3}$\\
\bottomrule
\end{tabularx}
\end{center}
\normalsize

\textbf{Information density ratio}: TE/token(answer) $\div$ TE/token(non-answer) = \textbf{29.5$\times$}.

\begin{center}
\SafeIncludeGraphics[width=0.82\linewidth]{prism-uploads/h2_fig2_info_density.png}\\[0.3em]
{\small\textbf{Figure 10.2.} Left: Total Effect (TE) by corruption region---answer-span tokens (3.1 avg) account for nearly half the total effect. Right: per-token information density is 29.5$\times$ higher at answer positions, confirming the spatial concentration that makes Q2C's redistribution redundant.}
\end{center}

Supplementary observation: $\text{TE(ans)} + \text{TE(non)} = 0.436 > \text{TE(full)} = 0.361$, interaction $= -0.075$ (sub-additive), indicating partial overlap between answer-region and non-answer-region corruption effects.

\FieldLabel{10.5.2 Reasoning}

Density ratio = 29.5$\times$. Context's causal information is \textbf{highly concentrated} in the answer span (on average, 3.1 tokens account for 47\% of total effect). This validates the mechanism argument's premise:

\begin{enumerate}
\item \textbf{Q2C's redistribution function is nearly redundant}: Useful Context information is already concentrated at the correct positions; redistributing via Question mediation adds almost no new value. Furthermore, QANet's Model Encoder contains 21 layers of Self-Attention whose function is precisely to redistribute information among Context positions---highly overlapping with Q2C. Q2C's function is structurally subsumed by downstream processing.
\item \textbf{C2Q's injection function is irreplaceable}: Question information has no other pathway into the Context representation stream after CQ Attention. No matter how powerful the downstream encoder, it cannot generate Question signal from nothing---only C2Q can provide this.
\item \textbf{Information concentration amplifies the asymmetry}: The more concentrated the information, the lower the marginal value of ``redistribution'' (already at correct positions), while the value of ``injecting new signal to correct positions'' remains unchanged.
\end{enumerate}

\FieldLabel{10.5.3 Limitations}

Corrupting answer-span tokens directly removes the target signal; high TE is expected. The truly informative observation is the non-answer token effect and the specific density ratio value---29.5$\times$ far exceeds baseline. Information concentration is measured at the Context input side. The ``novel information injection vs existing information redistribution'' functional distinction derives from CQ Attention's mathematical definition ($A = S_1 Q$ introduces $Q$, $B = S_1 S_2^{\!\top} C$ reorganises $C$)---it is deductive, not directly measured experimentally.

\ReportSubsection{10.6 H2 synthesis}

\textbf{Direct answer to H2.} The directional asymmetry is not only cross-architecture robust, but \textbf{dramatically intensifies} in QANet---from BiDAF's 2.7$\times$ to \textbf{61.0$\times$}.

Under QANet's Conv+Attention architecture, the joint causal effect of C2Q sub-components ($A$, $C \odot A$) is 61.0 times that of Q2C ($C \odot B$) at direction level (unequal intervention scale: C2Q zeros 2/4 dimensions, Q2C zeros 1/4). Under controlled single-component comparison (each zeroing exactly 1/4 of dimensions), the asymmetry remains stark: \textbf{$C \odot A$ vs $C \odot B$ = 15.4$\times$}, $A$ vs $C \odot B$ = 3.9$\times$. The 15.4$\times$ is the core ratio under comparable intervention scale. Removing C2Q drops F1 to 7.46; C2Q carries nearly all of CQ Attention's value. Q2C contributes only $-1.03$ F1, nearly redundant in QANet. \texttt{cq\_resizer} weight norms validate C2Q dominance at the direction level.

\textbf{Mechanism.} The asymmetry's root cause is that the two directions perform \textbf{fundamentally different operations}:

\begin{itemize}
\item \textbf{C2Q} ($A = S_1 Q$) injects Question embeddings into Context representations---the \emph{sole source of new information}. Question signal has no other pathway into the Context representation stream after CQ Attention. C2Q is irreplaceable.
\item \textbf{Q2C} ($C \odot B$, where $B = S_1 S_2^{\!\top} C$) redistributes Context vectors among Context positions---does not introduce Question's vector content. This Context-internal redistribution function highly overlaps with QANet's Model Encoder's 21 layers of Self-Attention. Q2C is replaceable.
\end{itemize}

Experiment~2 validates this mechanism's precondition: Context's causal information is highly concentrated at the answer span (density ratio 29.5$\times$, with 3.1 answer tokens accounting for 47\% of total effect). Information is already at the correct positions---redistributing it has almost no value, but injecting Question signal to ``mark'' these positions is indispensable.

\textbf{Amplification effect} (2.7$\times$ $\to$ 61$\times$): BiDAF's post-fusion processing consists of only 2 LSTM layers, insufficient to fully substitute Q2C's Context redistribution function, so Q2C retains $-3.6$ F1 of residual value. QANet has 21 layers including Self-Attention, whose Context redistribution capacity far exceeds Q2C's, compressing Q2C's residual value to $-1.03$ F1 (near zero). C2Q's irreplaceability is unchanged across architectures; Q2C's redundancy increases with encoder strength---hence the ratio amplifies from 2.7$\times$ to 61$\times$. This explanation is consistent with the data but has not been independently tested (e.g., by varying the number of Model Encoder layers).

\small
\begin{center}
{\bfseries Table 10.7. Chain of evidence for H2}\\[0.25em]
\begin{tabularx}{0.96\linewidth}{>{\raggedright\arraybackslash}p{0.24\linewidth}>{\raggedright\arraybackslash}p{0.30\linewidth}>{\raggedright\arraybackslash}X}
\toprule
\textbf{Evidence layer} & \textbf{Main result} & \textbf{Contribution to the final answer}\\
\midrule
Pre-validation & CQ Attention full ablation: F1 drops 89.9\% & Confirms information bottleneck; justifies internal dissection.\\
Core fact (Exp~1) & Sub-component ranking; controlled C2Q/Q2C = 15.4$\times$ & Establishes the asymmetry at every granularity with matched intervention scale.\\
Weight analysis (Exp~1c) & $A$-quadrant Frobenius norm highest (18.44) & Validates C2Q dominance from the parameter side.\\
Mechanism (Exp~2) & Answer token info density 29.5$\times$ & Validates the precondition: Context info is highly localised, making Q2C's redistribution redundant.\\
Architecture analysis & Model Encoder has 21 Self-Attention layers & Explains the amplification: downstream encoder subsumes Q2C's function, compressing its residual value.\\
\bottomrule
\end{tabularx}
\end{center}
\normalsize

Three limitations should be kept explicit. First, $A$ and $C \odot A$ have algebraic dependency ($C \odot A = C \times A$); zeroing one individually may not reflect independent contribution. Direction-level ablation (simultaneously zeroing $A$ + $C \odot A$) mitigates this and is the correct granularity for BiDAF comparison. Second, we argue that Q2C's Context redistribution function is subsumed by the Model Encoder's Self-Attention, but we did not directly measure their functional overlap (e.g., whether Q2C becomes more important when Model Encoder Self-Attention is removed). Third, both directions' attention matrices derive from the same similarity matrix $S$ under different normalisation directions, sharing trainable weight $w$---they are not truly independent modules but two views of the same computation.

%% ====================================================================
%%  SECTION 11 — H3
%% ====================================================================

\ReportSection{11 Hypothesis 3: Necessity of Pointer Asymmetric Wiring --- Testing Functional Specialisation in Iterative Encoding}

\ReportSubsection{11.1 Motivation and research question}

QANet's Model Encoder applies the same 7 encoder blocks iteratively three times (weight-shared), producing three intermediate representations:

\begin{tcolorbox}[enhanced,breakable,colback=CodeBg,colframe=CodeFrame,boxrule=0.5pt,arc=1.5mm,left=2.6mm,right=2.6mm,top=2.0mm,bottom=1.8mm]
\ttfamily
M0 $\xrightarrow{\text{7 blocks}}$ M1 $\xrightarrow{\text{same 7 blocks}}$ M2 $\xrightarrow{\text{same 7 blocks}}$ M3
\rmfamily
\end{tcolorbox}

The Pointer layer uses these three products \textbf{asymmetrically}:

\begin{tcolorbox}[enhanced,breakable,colback=CodeBg,colframe=CodeFrame,boxrule=0.5pt,arc=1.5mm,left=2.6mm,right=2.6mm,top=2.0mm,bottom=1.8mm]
\ttfamily
P(start) = softmax($w_1 \cdot [\text{M1};\, \text{M2}]$) \quad$\leftarrow$ pass 1 + pass 2\\
P(end)\hphantom{t} = softmax($w_2 \cdot [\text{M1};\, \text{M3}]$) \quad$\leftarrow$ pass 1 + pass 3
\rmfamily
\end{tcolorbox}

\FieldLabel{11.1.1 Design provenance}

The QANet paper merely states ``We adopt the strategy of Seo et al.\ (2016)'' for this design, providing no theoretical justification. Tracing back to BiDAF, the original form is: start prediction uses the modelling layer output $M$, while end prediction passes $M$ through an additional BiLSTM with \emph{independent} parameters to obtain $M_2$. In BiDAF, this asymmetry has at least structural justification---the extra BiLSTM with independent parameters can learn end-specific features. However, QANet adapted this pattern into \emph{weight-shared} encoder iterations: the third iteration applies the exact same transformation function as the first two, significantly weakening the premise that ``additional processing produces specialised information.'' This never-formally-tested design choice implies three assumptions:

\begin{enumerate}
\item \textbf{Depth specialisation}: $M_2$ (14 layers) and $M_3$ (21 layers) carry different information suited for different prediction tasks---$M_2$'s moderate depth suits start localisation, $M_3$'s additional processing provides specialised refinement for end prediction.
\item \textbf{Shared anchor}: $M_1$ (7 layers) provides the base localisation signal needed by both predictions.
\item \textbf{Iterative gain}: More encoder iterations produce better or at least different representations.
\end{enumerate}

\begin{tcolorbox}[enhanced,breakable,colback=white,colframe=BugBlueBorder,boxrule=0.5pt,arc=1.5mm,left=2.8mm,right=2.8mm,top=1.8mm,bottom=1.8mm]
\textbf{H3 research question.} Does the functional specialisation hypothesis implied by the Pointer's asymmetric wiring---$M_1$ serving as a shared anchor providing base signal, $M_2$ and $M_3$ each developing distinct expertise due to different processing depths---hold in the trained model?
\end{tcolorbox}

\ReportSubsection{11.2 Methodological note: eval-time intervention}

All experiments in this study are performed at \textbf{eval time}---wiring replacement, ablation, and information probing are all executed on a fully trained model without retraining. Our research goal is \textbf{mechanistic analysis}---understanding the functional roles that components currently serve in the trained model. This aligns with the causal tracing methodology of Meng et al.\ (2022). Training-time ablation allows compensatory learning, erasing functional traces---making it unsuitable for mechanistic analysis.

Wiring replacement is a form of eval-time intervention: the encoder is unchanged, only the Pointer's input combinations are modified. This directly measures ``how the model currently uses each iteration product.'' Since Pointer weights are trained on specific input distributions, replacement introduces distribution mismatch---this inherent limitation is mitigated in Experiment~2 through Pointer-weight-independent linear probes.

\ReportSubsection{11.3 Pre-validation: how large is the asymmetric design's advantage?}

We directly test the necessity of the asymmetric design through wiring replacement. At eval time, we change the Pointer's input combinations (encoder unchanged), evaluating on the full dev set of 10{,}465 samples. The large sample size ensures that even small F1 differences (e.g., 0.23) are measured with high precision.

\small
\begin{center}
{\bfseries Table 11.1. Pointer wiring replacement results}\\[0.25em]
\begin{tabularx}{0.96\linewidth}{>{\raggedright\arraybackslash}p{0.10\linewidth}>{\centering\arraybackslash}p{0.12\linewidth}>{\centering\arraybackslash}p{0.12\linewidth}>{\centering\arraybackslash}p{0.08\linewidth}>{\centering\arraybackslash}p{0.08\linewidth}>{\raggedright\arraybackslash}X}
\toprule
\textbf{Config} & \textbf{Start input} & \textbf{End input} & \textbf{F1} & \textbf{$\Delta$F1} & \textbf{Hypothesis tested}\\
\midrule
original & M1+M2 & M1+M3 & 70.24 & --- & Baseline (asymmetric)\\
sym\_M2 & M1+M2 & M1+M2 & 70.01 & $-0.23$ & Is M3 necessary?\\
sym\_M3 & M1+M3 & M1+M3 & 69.70 & $-0.54$ & Is M2 better?\\
no\_M1 & M2+M3 & M2+M3 & 69.67 & $-0.58$ & Is M1 necessary?\\
swap & M1+M3 & M1+M2 & 69.42 & $-0.82$ & Does direction matter?\\
only\_M1 & M1+M1 & M1+M1 & 63.97 & $-6.28$ & Are M2/M3 necessary?\\
\bottomrule
\end{tabularx}
\end{center}
\normalsize

Test results for the three design assumptions:

\textbf{Assumption~1 (Depth specialisation)---barely holds}: The asymmetric design (original) outperforms the symmetric design (sym\_M2) by only \textbf{0.23~F1}. The supposed ``specialisation'' from assigning $M_3$ to end prediction provides near-negligible marginal benefit. More notably, $M_2$ doing both tasks (sym\_M2: $-0.23$) \emph{outperforms} $M_3$ doing both tasks (sym\_M3: $-0.54$)---the more processed product is actually worse.

\textbf{Assumption~2 (Shared anchor)---near-redundant}: Removing $M_1$ (no\_M1: $-0.58$) barely affects performance, but $M_1$ alone (only\_M1: $-6.28$) is catastrophic. $M_1$ participates in both predictions yet contributes the least---the shared anchor provides negligible incremental value.

\textbf{Assumption~3 (Iterative gain)---partially holds with diminishing returns}: $M_2$/$M_3$ are essential (only\_M1: $-6.28$), but $M_2$ already captures the vast majority of incremental value. The second iteration's increment (only\_M1 vs sym\_M2: 6.05~F1) far exceeds the third iteration's marginal contribution (near zero or negative).

\textbf{Note on swap.} swap ($-0.82$) is the largest drop among non-only\_M1 configs, but this does not indicate $M_2$/$M_3$ carry different information. swap causes \emph{both} $w_1$ and $w_2$ to receive out-of-distribution inputs simultaneously ($w_1$ trained on $[M_1; M_2]$ but receives $[M_1; M_3]$; $w_2$ vice versa), whereas sym\_M2 only mismatches $w_2$ ($-0.23$) and sym\_M3 only mismatches $w_1$ ($-0.54$). If the mismatch effects are independent, expected swap $\approx 0.23 + 0.54 = 0.77$, actual $0.82$---highly consistent. The swap loss comes from \textbf{dual distribution mismatch}, not $M_2$/$M_3$ information difference.

\begin{center}
\SafeIncludeGraphics[width=0.78\linewidth]{prism-uploads/h3_fig1_wiring_replacement.png}\\[0.3em]
{\small\textbf{Figure 11.1.} Pointer wiring replacement results. The asymmetric design (original) outperforms the best symmetric alternative (sym\_M2) by only 0.23~F1---far smaller than the 6.28~F1 gap from removing $M_2$/$M_3$ entirely (only\_M1). The marginal value of asymmetric wiring is negligible.}
\end{center}

\ReportSubsection{11.4 Experiment design}

Three experiments progressively explain the three design failures observed in pre-validation:

\textbf{Experiment~1 (Why is the shared anchor redundant?)} Validate the cause of $M_1$ redundancy from parameter and information sides: (a)~Pointer weight $M_1$-quadrant norms---did the model learn to downweight $M_1$? (b)~$M_1$--$M_2$ CKA---has $M_2$ absorbed $M_1$'s information?

\textbf{Experiment~2 (Is $M_2$ and $M_3$'s information actually different?)} Use Pointer-weight-independent linear probes to directly measure $M_1$/$M_2$/$M_3$'s start/end boundary information, determining whether $M_2 > M_3$ reflects intrinsic information difference or adaptation bias.

\textbf{Experiment~3 (What did the third iteration do?)} Even if Experiment~2 shows $M_2$ and $M_3$'s boundary information is nearly identical, did the third iteration still change the representation structure? We characterise the $M_2 \to M_3$ transformation through representation structure analysis.

\ReportSubsection{11.5 Experiment 1: Why is the shared anchor redundant?}

\FieldLabel{11.5.1 Pointer weight quadrant analysis}

Pointer's $w_1$ has shape $[2d]$ ($d=96$), where the first $d$ dimensions correspond to $M_1$ features and the last $d$ to $M_2$ features (concatenation order: $[M_1;\, M_2]$). Similarly, $w_2$'s first $d$ correspond to $M_1$, last $d$ to $M_3$.

\small
\begin{center}
{\bfseries Table 11.2. Pointer weight quadrant L2 norms}\\[0.25em]
\begin{tabularx}{0.78\linewidth}{>{\centering\arraybackslash}p{0.12\linewidth}>{\centering\arraybackslash}p{0.28\linewidth}>{\centering\arraybackslash}p{0.28\linewidth}>{\centering\arraybackslash}X}
\toprule
\textbf{Weight} & \textbf{M1 quadrant L2} & \textbf{M2/M3 quadrant L2} & \textbf{Ratio}\\
\midrule
$w_1$ & 0.0902 & 0.1773 & 1.97$\times$\\
$w_2$ & 0.0728 & 0.1232 & 1.69$\times$\\
\bottomrule
\end{tabularx}
\end{center}
\normalsize

The $M_2$/$M_3$ quadrant norms in $w_1$ and $w_2$ are \textbf{1.97$\times$} and \textbf{1.69$\times$} their $M_1$ counterparts respectively. \textbf{Caveat}: in linear models, weight norms are influenced by input feature magnitudes---if $M_1$ tokens have systematically larger average magnitude than $M_2$/$M_3$, smaller weights may reflect scale compensation rather than reduced dependence. The quadrant ratio should be interpreted as \emph{directional support} (consistent with behavioural evidence), not independent strong evidence.

\FieldLabel{11.5.2 CKA information overlap}

Compute linear CKA for $M_1$--$M_2$, $M_2$--$M_3$, $M_1$--$M_3$ (full dev set).

\small
\begin{center}
{\bfseries Table 11.3. CKA between iteration products}\\[0.25em]
\begin{tabularx}{0.80\linewidth}{>{\centering\arraybackslash}p{0.12\linewidth}>{\centering\arraybackslash}p{0.10\linewidth}>{\raggedright\arraybackslash}X}
\toprule
\textbf{Pair} & \textbf{CKA} & \textbf{Interpretation}\\
\midrule
$M_1$--$M_2$ & 0.871 & $M_2$ contains most of $M_1$'s linearly readable information\\
$M_2$--$M_3$ & 0.954 & $M_2$ and $M_3$ share nearly identical representational structure\\
$M_1$--$M_3$ & 0.717 & $M_3$ is most distant from $M_1$---iterations progressively transform representations\\
\bottomrule
\end{tabularx}
\end{center}
\normalsize

CKA($M_1$, $M_2$) = 0.871 indicates some linear structure difference between $M_1$ and $M_2$ (not the $>0.9$ near-identity), meaning $M_2$ does not perfectly preserve $M_1$'s linear format. However, no\_M1 causes only $-0.58$ F1---indicating that the information in $M_1$'s ``unique'' linear structure (the $\sim$13\% CKA gap) is not important for the Pointer's boundary predictions.

CKA($M_2$, $M_3$) = 0.954 shows the third iteration barely changed the linear structure of representations---$M_2$ and $M_3$ essentially carry the same information.

\FieldLabel{11.5.3 Experiment 1 summary}

Weight analysis and CKA explain $M_1$ redundancy from two independent angles:

\begin{itemize}
\item \textbf{Parameter side}: Pointer weights' $M_2$/$M_3$ quadrant norms are 1.7--2.0$\times$ the $M_1$ quadrant---directionally consistent with low $M_1$ importance, but weight norms may be influenced by input feature scale and are not treated as independent strong evidence.
\item \textbf{Information side}: CKA($M_1$, $M_2$) = 0.871 shows $M_2$ and $M_1$ share substantial linear structure. Combined with no\_M1's $-0.58$ F1 behavioural evidence, $M_1$ provides negligible incremental value in the $[M_1;\, M_2]$ concatenation.
\end{itemize}

The shared anchor design exists but provides negligible incremental value---this is the \textbf{first unsupported assumption} of the asymmetric design.

\ReportSubsection{11.6 Experiment 2: Is $M_2$ and $M_3$'s information actually different?}

\textbf{Goal}: Directly measure $M_2$ vs $M_3$ boundary information difference, independent of Pointer weights. This tests the core assumption of the asymmetric design---whether different-depth iteration products carry different information.

\textbf{Method}: Train logistic regression on $M_1$/$M_2$/$M_3$'s per-token representations for \texttt{is\_start\_token} and \texttt{is\_end\_token}. 500 samples, 80/20 split, ROC-AUC. The Pointer itself is a linear projection ($w \cdot [M;\, M']$), so linear probes test the \emph{same type} of information readability.

\small
\begin{center}
{\bfseries Table 11.4. Linear probe AUC for boundary detection}\\[0.25em]
\begin{tabularx}{0.68\linewidth}{>{\centering\arraybackslash}p{0.22\linewidth}>{\centering\arraybackslash}p{0.20\linewidth}>{\centering\arraybackslash}X}
\toprule
\textbf{Representation} & \textbf{is\_start AUC} & \textbf{is\_end AUC}\\
\midrule
$M_1$ & 0.977 & 0.979\\
$M_2$ & \textbf{0.979} & \textbf{0.981}\\
$M_3$ & 0.978 & 0.980\\
\bottomrule
\end{tabularx}
\end{center}
\normalsize

\textbf{Core finding: no functional specialisation detected.} $M_2$ achieves the highest AUC for \emph{both} start and end, while $M_3$ is slightly lower for both. The expected ``$M_2 \to$ start, $M_3 \to$ end'' specialisation pattern \textbf{does not appear}---all differences are within $\Delta$AUC = 0.001, well below the estimated standard error ($\sim$0.005--0.01), in the range of statistical noise. The three iteration products' boundary information content is statistically indistinguishable.

\textbf{Cross-validation with wiring replacement}: Probes (information level) and wiring replacement (behavioural level) consistently point to $M_2 \geq M_3$. The convergence of two methods strengthens the conclusion's robustness: within the detection scope of our experiments, \textbf{no meaningful functional specialisation is observed} between $M_2$ and $M_3$, and the asymmetric wiring lacks information-level support---this is the \textbf{second unsupported assumption} of the asymmetric design.

\FieldLabel{11.6.1 Limitations}

\texttt{is\_start} and \texttt{is\_end} labels are extremely imbalanced ($\sim$0.70\% positive). Linear probes operate on $M_2$ or $M_3$ individually, but the Pointer processes $[M_1;\, M_2]$ concatenation---joint information may exceed the sum of parts. Probes detect information \emph{presence} in the representation, not how much the Pointer actually \emph{uses}.

\ReportSubsection{11.7 Experiment 3: What did the third iteration do?}

Experiment~2 shows $M_2$ and $M_3$'s boundary information is nearly identical ($\Delta$AUC = 0.001), yet CKA($M_2$, $M_3$) = 0.954, not 1.0---the third iteration did change something. What did it change?

Each encoder block's pipeline is Conv\_0 $\to$ Conv\_1 $\to$ Self-Attention $\to$ FFN. Self-Attention is the only global operation. If the third iteration's Self-Attention aggregates token representations toward the global mean (``over-smoothing''), adjacent tokens become more similar, making precise boundary positions harder to distinguish.

\textbf{Method}: Compute the following representation structure metrics on $M_1$/$M_2$/$M_3$ (200 samples):

\begin{enumerate}
\item \textbf{Local contrast}: $1 - \cos(\mathbf{h}_t, \mathbf{h}_{t+1})$, averaged. High = adjacent tokens more distinguishable.
\item \textbf{Answer boundary sharpness}: Cosine distance between answer boundary positions ($y_1$, $y_2$) and adjacent positions. High = sharper boundaries.
\item \textbf{Token norm coefficient of variation (CoV)}: Std/mean of per-token $\ell_2$ norms. High = more position-selective activation.
\end{enumerate}

\small
\begin{center}
{\bfseries Table 11.5. Representation structure analysis}\\[0.25em]
\begin{tabularx}{0.92\linewidth}{>{\raggedright\arraybackslash}p{0.26\linewidth}>{\centering\arraybackslash}p{0.12\linewidth}>{\centering\arraybackslash}p{0.12\linewidth}>{\centering\arraybackslash}p{0.12\linewidth}>{\centering\arraybackslash}X}
\toprule
\textbf{Metric} & \textbf{$M_1$} & \textbf{$M_2$} & \textbf{$M_3$} & \textbf{$M_2 \to M_3$ change}\\
\midrule
Local contrast & 0.1162 & 0.0669 & 0.0480 & $-28.2\%$\\
Answer boundary sharpness & 0.2213 & 0.1354 & 0.1017 & $-24.9\%$\\
Norm CoV & 0.1627 & 0.1361 & 0.1273 & $-6.5\%$\\
\bottomrule
\end{tabularx}
\end{center}
\normalsize

\textbf{All three metrics consistently indicate $M_3$ is smoother}: local contrast drops 28\%, answer boundary sharpness drops 25\%, norm diversity drops 6.5\%. The progressive pattern $M_1 \to M_2 \to M_3$ shows continuous smoothing---each iteration reduces local discriminability.

\begin{center}
\SafeIncludeGraphics[width=0.72\linewidth]{prism-uploads/h3_fig2_repr_structure.png}\\[0.3em]
{\small\textbf{Figure 11.2.} Representation structure across three iteration products. All metrics decrease monotonically $M_1 \to M_2 \to M_3$, with the $M_2 \to M_3$ reduction (local contrast $-28\%$, boundary sharpness $-25\%$) indicating continued smoothing after the second iteration has already captured all incremental boundary information.}
\end{center}

\FieldLabel{11.7.1 Why does $M_1 \to M_2$ smooth more yet improve information?}

\small
\begin{center}
{\bfseries Table 11.6. Smoothing vs information: progressive pattern}\\[0.25em]
\begin{tabularx}{0.80\linewidth}{>{\raggedright\arraybackslash}p{0.16\linewidth}>{\centering\arraybackslash}p{0.28\linewidth}>{\centering\arraybackslash}X}
\toprule
\textbf{Iteration} & \textbf{Local contrast change} & \textbf{Probe AUC change (start)}\\
\midrule
$M_1 \to M_2$ & $0.116 \to 0.067$ (\textbf{$-42\%$}) & $0.977 \to 0.979$ (\textbf{$+0.002$})\\
$M_2 \to M_3$ & $0.067 \to 0.048$ ($-28\%$) & $0.979 \to 0.978$ ($-0.001$)\\
\bottomrule
\end{tabularx}
\end{center}
\normalsize

The second iteration smooths \emph{more} than the third (42\% vs 28\%), yet simultaneously improves boundary information. This is not contradictory: $M_1$ is the CQ Attention output after only one encoding pass, retaining substantial local noise. The second iteration's smoothing is \textbf{beneficial information compression}---filtering noise while integrating cross-position semantic relationships. By $M_2$, the representation has approached an \textbf{approximate fixed point} of the shared transformation function (CKA($M_2$, $M_3$) = 0.954), and useful information compression is complete. The third iteration continues smoothing but has \textbf{no new information to integrate}---it is ``smoothing without return.''

\FieldLabel{11.7.2 Alternative hypothesis: stochastic depth causing insufficient training}

The ``smoothing without return'' explanation above assumes the third iteration's transformation is functionally redundant. A competing explanation exists: during training, stochastic depth skips sub-layers with probability $p_{\text{drop}} \times l/L$ (deeper layers have higher skip probability). Layers 15--21 (traversed only by $M_3$) have the highest skip probability, which may cause $M_3$'s representations to be \textbf{insufficiently trained} rather than ``over-smoothed''---the third iteration's parameters may not have received enough gradient updates to learn useful transformations. If this explanation holds, the $M_2 \to M_3$ structural changes reflect a training regularisation side-effect, not an inherent deficiency of weight-shared iteration. Current experiments cannot distinguish the two hypotheses, as both predict $M_3$ representation quality below $M_2$. We retain both explanations in parallel.

\FieldLabel{11.7.3 Key tension: structure changes substantially, but information barely changes}

The third iteration's representation structure changes are significant (20--28\%), yet linear probe AUC changes by only 0.001. The third iteration substantially smooths the geometric structure of representations---adjacent tokens become more similar, boundaries become blurrier, norms become more uniform---but boundary position information is \textbf{preserved} in a more distributed encoding, remaining extractable by linear classifiers. The information persists, but the ``format'' changes.

This explains the mere 0.31~F1 difference between sym\_M2 and sym\_M3: $M_3$'s representations are structurally smoother than $M_2$'s, but the Pointer's linear projection can still extract sufficient boundary signal---just slightly less efficiently.

\FieldLabel{11.7.4 Connection to H1 findings}

H1 found that Conv\_1 is the most important sub-layer in the Model Encoder---it produces high-discriminability local features for Self-Attention to use. The third iteration's Self-Attention global aggregation smooths the local contrast that Conv\_1 established ($-28\%$), echoing H1: Conv\_1 builds local features $\to$ Self-Attention utilises and progressively consumes them $\to$ too many iterations cause structural degradation, even though information is not lost.

\FieldLabel{11.7.5 Limitations}

Representation structure analysis is descriptive (correlational), not causal. ``$M_3$ is smoother'' does not equal ``$M_3$ is worse \emph{because} it is smoother''---confounding variables may exist. Conv and FFN in encoder blocks also affect representations; attributing changes to Self-Attention's ``over-aggregation'' requires finer-grained per-sub-layer analysis. $M_1$, $M_2$, $M_3$ use the same encoder parameters; each iteration applies the same function to different input distributions. ``Over-smoothing'' is not a problem of Self-Attention itself, but a cumulative effect of iteratively applying the same transformation.

\ReportSubsection{11.8 H3 synthesis}

\textbf{Direct answer to H3.} The functional specialisation hypothesis implied by the Pointer's asymmetric wiring \textbf{is not supported in the trained model}. All three assumptions of the asymmetric design fail to find support in our experiments:

\textbf{1.\ Depth specialisation---not detected.} Pre-validation observes a behavioural advantage of only 0.23~F1 for the asymmetric design; the following information-level analysis explains why. Linear probes show $M_2$ achieves the highest AUC for both start (0.979) and end (0.981), with $M_3$ slightly lower for both (start = 0.978, end = 0.980). The expected ``$M_2 \to$ start, $M_3 \to$ end'' specialisation pattern \textbf{does not appear}---all differences are within $\Delta$AUC = 0.001, well below the estimated standard error ($\sim$0.005--0.01), in the range of statistical noise. Wiring replacement (behavioural) and probes (informational) consistently point to $M_2 \geq M_3$.

\textbf{2.\ Shared anchor---near-redundant.} Replacing $M_1$ causes only $-0.58$~F1 (behavioural). Pointer weights' $M_2$/$M_3$ quadrant norms are 1.7--2.0$\times$ $M_1$ quadrants (parameter level; $w_1$: 1.97$\times$, $w_2$: 1.69$\times$). CKA($M_1$, $M_2$) = 0.871 shows $M_2$ and $M_1$ share substantial linear structure (information level). All three angles point in the same direction: \textbf{$M_1$ provides negligible incremental value in the $[M_1;\, M_2]$ concatenation}. Note that no\_M1 substitutes $M_2$/$M_3$ into $M_1$'s position (not pure zeroing); the $-0.58$ loss partly reflects representational similarity.

\textbf{3.\ Iterative gain---second iteration captures the vast majority of value; third iteration's marginal contribution is near zero.} The only\_M1 vs sym\_M2 gap of 6.05~F1 shows the second iteration is essential; but the third iteration's marginal F1 contribution is negligible. Experiment~3 observes a noteworthy phenomenon: the third iteration further smooths representation structure (local contrast $-28\%$, boundary sharpness $-25\%$), yet linearly readable boundary information barely changes ($\Delta$AUC = 0.001). Information persists in a more distributed encoding---\textbf{structural change does not equal information loss}. However, $M_3$'s structural changes may partly stem from stochastic depth causing insufficient training (see Experiment~3 limitations), not solely from inherent effects of iteration.

The asymmetric wiring ($M_2 \to$ start, $M_3 \to$ end) outperforms symmetric wiring (sym\_M2) by only 0.23~F1. This marginal advantage finds no support in the boundary information measured by linear probes ($\Delta$AUC = 0.001), and more likely reflects Pointer weight adaptation to training distributions.

\small
\begin{center}
{\bfseries Table 11.7. Chain of evidence for H3}\\[0.25em]
\begin{tabularx}{0.96\linewidth}{>{\raggedright\arraybackslash}p{0.22\linewidth}>{\raggedright\arraybackslash}p{0.30\linewidth}>{\raggedright\arraybackslash}X}
\toprule
\textbf{Evidence layer} & \textbf{Main result} & \textbf{Contribution to the final answer}\\
\midrule
Behavioural (pre-validation) & Asymmetric advantage only 0.23~F1; $M_1$ redundant ($-0.58$); $M_2 > M_3$ universality; swap's $-0.82$ explained by dual distribution mismatch. & Establishes that the asymmetric design's behavioural advantage is negligible.\\
Parameter (Exp~1a) & $w_1$/$w_2$ $M_2$/$M_3$ quadrant norms 1.7--2.0$\times$ $M_1$ quadrants (directional, not independent evidence). & Consistent with low $M_1$ importance at the parameter level.\\
Information (Exp~1b + Exp~2) & CKA: $M_1$--$M_2$ overlap 0.871, $M_2$--$M_3$ near-identity 0.954. Probes: no specialisation ($\Delta$AUC = 0.001). & Confirms $M_1$ redundancy and absence of $M_2$/$M_3$ functional specialisation.\\
Mechanism (Exp~3) & Third iteration smooths structure (20--28\%) without destroying information. & Explains \emph{why} specialisation fails: the shared transformation converges to a near-fixed point.\\
\bottomrule
\end{tabularx}
\end{center}
\normalsize

\FieldLabel{11.8.1 Connection to other hypotheses}

H1 demonstrated that Conv\_1 is an indispensable component in the Model Encoder (removal causes $-32.54$~F1); H3 demonstrates that $M_1$ is redundant in the Pointer (removal causes only $-0.58$~F1). Together they show that, in the trained model, QANet's architecture contains \textbf{coexisting critical bottlenecks and redundancies}. H2 demonstrated that C2Q is irreplaceable while Q2C is redundant in CQ Attention---the same theme of ``irreplaceability vs redundancy.'' All three hypotheses reveal the uneven distribution of functional importance within QANet from different levels, as observed through eval-time mechanistic analysis.

The direct echo between Experiment~3 and H1 is particularly notable: H1 found that Conv\_1 building local features is the core function of the Model Encoder; H3 found that iteratively applied Self-Attention progressively consumes these local features ($M_1 \to M_2 \to M_3$ local contrast drops from 0.116 to 0.048). This forms a complete picture: Conv\_1 produces, Self-Attention consumes, and too many iterations lead to over-consumption.

\FieldLabel{11.8.2 Limitations}

Three limitations should be kept explicit. First, \textbf{distribution mismatch in eval-time wiring replacement}: Pointer weights are trained on specific input distributions; wiring replacement changes the input distribution. Both sym\_M2 and sym\_M3 have 50\% input distribution mismatch. How much of their difference (0.31~F1) comes from intrinsic $M_2$/$M_3$ difference vs asymmetric distribution mismatch effects cannot be fully separated. Experiment~2's linear probes partially mitigate this.

Second, \textbf{causal attribution}: Experiment~3's representation analysis is descriptive. The co-occurrence of ``$M_3$ is smoother'' and ``boundary information format changes'' is observed, not causally established. Full causal verification would require selectively disabling Self-Attention in the third pass---high implementation complexity.

Third, \textbf{stochastic depth training-time effects}: As discussed in Experiment~3's main text, $M_3$'s structural changes may partly stem from stochastic depth causing insufficient training rather than iterative over-smoothing. Current experiments cannot distinguish the two. This is the primary alternative explanation for Conclusion~3 (iterative gain exhaustion).

\end{document}
